% Traduction : les locutions adverbiales ; redacción : ensayo argumentativo sobre el medioambiente — Espagnol Tle A — Cours signé OnBuch+
% Programme officiel MINESEC. Compiler avec : tectonic E10-locuciones-adverbiales-ensayo-ambiente.tex
\newcommand{\DOCMATIERE}{Espagnol}
\newcommand{\DOCNIVEAU}{Tle A}
\newcommand{\DOCMODULE}{Módulo 2 : Environnement, santé et bien-être}
\newcommand{\DOCLECON}{E10}
\newcommand{\DOCTITRE}{Traduction : les locutions adverbiales ; redacción : ensayo argumentativo sobre el medioambiente}
% OnBuch+ — préambule « cours signés » ENRICHI (compilé avec tectonic / XeLaTeX)
% Système visuel « Pop » d'OnBuch+ : fond crème (jamais blanc), bordures encre
% épaisses, ombres « dures » décalées, titres Archivo Black en capitales.
% Version MATHÉMATIQUES : graphes (pgfplots/tikz), boîtes pédagogiques
% (définition, propriété, méthode, attention, le savais-tu, exercice résolu,
% exercice, TP, bilan), page de garde et signature OnBuch+ ancrée en bas.
%
% Macros attendues dans main.tex AVANT \input{preamble} :
%   \DOCMATIERE  \DOCNIVEAU  \DOCMODULE  \DOCLECON (numéro)  \DOCTITRE
\documentclass[11pt]{article}

\usepackage{fontspec}
\usepackage[spanish,french]{babel} % français = langue principale ; l'espagnol via \esp{..} / espagnol
\usepackage[a4paper,margin=2cm,top=3.4cm,bottom=2.8cm,headheight=1.4cm,footskip=1.3cm]{geometry}
\usepackage{amsmath,amssymb,amsfonts}
\usepackage{mathtools} % \xmapsto, \coloneqq... souvent utilisés par les modèles
\usepackage{cancel} % simplifications barrées (bilans de forces)
% \abs{z} (module, valeur absolue) et \norm{u} (norme), utilisés par les modèles
\providecommand{\abs}{}\let\abs\relax\DeclarePairedDelimiter{\abs}{\lvert}{\rvert}
\providecommand{\norm}{}\let\norm\relax\DeclarePairedDelimiter{\norm}{\lVert}{\rVert}
\usepackage[table]{xcolor} % [table] : fournit aussi \rowcolors (alternance de lignes)
\usepackage{pagecolor}
\usepackage{tikz}
\usetikzlibrary{arrows.meta,positioning,calc,shapes.geometric,shapes.misc,shapes.arrows,decorations.pathmorphing,decorations.pathreplacing,decorations.markings,patterns,shadows,mindmap,trees,fit,backgrounds,matrix,intersections,angles,quotes}
\usepackage{pgfplots}
\usepackage[version=4]{mhchem} % équations nucléaires \ce{^{235}_{92}U}
\usepackage[european,siunitx]{circuitikz} % schémas électriques
\pgfplotsset{compat=1.18}
\usepgfplotslibrary{fillbetween,groupplots} % aires entre courbes (calcul intégral)
\usepackage{enumitem}
\usepackage{fancyhdr}
% [most] charge toutes les bibliothèques tcolorbox (listings, minted, raster,
% documentation...) et gonfle inutilement la mémoire TeX sur un document long
% (~300 boîtes) : on ne charge que ce qui est réellement utilisé.
\usepackage[skins,breakable]{tcolorbox}
\usepackage{graphicx}
\usepackage{colortbl}
\usepackage{booktabs}
\usepackage{tabularx}
\usepackage{diagbox} % case d'en-tête diagonale des tableaux à double entrée
% Colonnes extensibles centrée / gauche / droite (C, L, R), souvent utilisées par les modèles
\newcolumntype{C}{>{\centering\arraybackslash}X}
\newcolumntype{L}{>{\raggedright\arraybackslash}X}
\newcolumntype{R}{>{\raggedleft\arraybackslash}X}
\usepackage{array}
\usepackage{multicol}
\usepackage{multirow}
\usepackage{siunitx}
% parse-numbers=false : accepte aussi \SI{2,5 \times 10^{-3}}{...}
\sisetup{locale=FR,output-decimal-marker={,},group-minimum-digits=4,parse-numbers=false}
\DeclareSIUnit\ohm{\ensuremath{\symup{\Omega}}} % U+2126 absent de la police
\DeclareSIUnit\division{div} % réglages d'oscilloscope (ms/div, V/div)
\DeclareSIUnit\tr{tr} % tours (vitesse de rotation, ex. \SI{1500}{\tr\per\minute})
\DeclareSIUnit\molar{M} % concentration molaire (ex. \SI{3}{\molar}, \SI{1}{\micro\molar})
\DeclareSIUnit\an{an} % année (le modèle confond parfois avec \year, un compteur TeX) — ex. \SI{5}{\kilo\meter\per\an}
\DeclareSIUnit\FCFA{FCFA} % franc CFA (ex. \SI{500}{\FCFA\per\kilo\gram})
\DeclareSIUnit\degreeBrix{\textdegree Brix} % teneur en sucre d'un sirop/jus (réfractométrie)
\DeclareSIUnit\month{mois} % le modèle utilise parfois \month, un compteur TeX, comme unité de durée
\usepackage{qrcode}
% \degree : commande fréquemment utilisée par les modèles pour un angle ou une
% température en dehors de \SI{}{} (non fournie par les paquets chargés).
\providecommand{\degree}{^{\circ}}
% \qed (fin de démonstration), utilisé par les modèles sans amsthm
\providecommand{\qed}{\hfill$\square$}
% \barre : double barre des valeurs interdites dans un tableau de variations (tkz-tab)
\providecommand{\barre}{\ensuremath{\|}}
% environnement proof (amsthm non chargé) utilisé par les modèles
\ifdefined\proof\else\newenvironment{proof}[1][Démonstration]{\par\noindent\textit{#1.}\ }{\qed\par}\fi
\usepackage{lastpage}
\usepackage{hyperref}
\hypersetup{hidelinks}

% ============================================================
% Palette « Pop » OnBuch+ — mêmes hex que la classe P (plus_ui.dart)
% ============================================================
\definecolor{poporange}{HTML}{F59321}
\definecolor{popdark}{HTML}{E07A0C}
\definecolor{popcream}{HTML}{FFF6E8}
\definecolor{popink}{HTML}{1C1714}
\definecolor{poppurple}{HTML}{7A5AE0}
\definecolor{popgreen}{HTML}{2E9E6A}
\definecolor{poppink}{HTML}{E0457B}
\definecolor{popblue}{HTML}{2F7DE1}
\definecolor{popgold}{HTML}{C98A12}
\definecolor{popmuted}{HTML}{8A7F76}
\definecolor{popline}{HTML}{EADFD2}
\definecolor{popgray}{HTML}{8A7F76}
\colorlet{popgrey}{popgray}
% Alias tolérants (le modèle nomme parfois les couleurs différemment)
\colorlet{popwhite}{white}
\colorlet{popblack}{popink}
\colorlet{popred}{poppink}
\colorlet{popyellow}{popgold}
% Teintes claires (fonds de tableaux/encadrés) — le modèle nomme parfois
% les couleurs avec un suffixe L (ex. popblueL) sans les avoir définies.
\colorlet{poporangeL}{poporange!20}
\colorlet{popdarkL}{popdark!20}
\colorlet{poppurpleL}{poppurple!20}
\colorlet{popgreenL}{popgreen!20}
\colorlet{poppinkL}{poppink!20}
\colorlet{popblueL}{popblue!20}
\colorlet{popgoldL}{popgold!20}
\colorlet{popredL}{popred!20}
\colorlet{popgrayL}{popgray!20}
% Teintes claires pour fonds de tableaux / zones
\colorlet{popblueL}{popblue!10!white}
\colorlet{poporangeL}{poporange!14!white}
\colorlet{popgreenL}{popgreen!12!white}
\colorlet{poppurpleL}{poppurple!12!white}
\colorlet{poppinkL}{poppink!12!white}
\colorlet{popgoldL}{popgold!14!white}

\pagecolor{popcream}

% ============================================================
% Typographie
% ============================================================
\defaultfontfeatures{Ligatures=TeX}
\setmainfont{PlusJakartaSans-Medium.ttf}[Path=../../../fonts/,BoldFont=PlusJakartaSans-Bold.ttf]
% Symboles, API et emojis absents de la police principale : repli sur DejaVu Sans (caractères actifs)
\newfontface\symfont{DejaVuSans.ttf}[Path=../../../fonts/]
\makeatletter
\newcommand\symactive[1]{\catcode#1=13 \begingroup\lccode`\~=#1 \lowercase{\endgroup\def~}{{\symfont\char#1}}}
\newcommand\symblank[1]{\catcode#1=13 \begingroup\lccode`\~=#1 \lowercase{\endgroup\def~}{}}
\newcommand\symhyph[1]{\catcode#1=13 \begingroup\lccode`\~=#1 \lowercase{\endgroup\def~}{\mbox{-}}}
\newcount\symc
\newcommand\symrange[2]{\symc=#1 \@whilenum\symc<#2 \do{\expandafter\symactive\expandafter{\the\symc}\advance\symc by 1 }}
\newcommand\symblankrange[2]{\symc=#1 \@whilenum\symc<#2 \do{\expandafter\symblank\expandafter{\the\symc}\advance\symc by 1 }}
\makeatother
\symrange{"0250}{"02FF}\symactive{"2713}\symactive{"2714}\symactive{"2717}\symactive{"2718}\symactive{"2610}\symactive{"2611}\symactive{"2612}
\symhyph{"2011}\symhyph{"2E1A}\symblankrange{"1F300}{"1FAFF}\symblank{"FE0F}
% Maths en sans-serif (Fira Math) avec chiffres et lettres droites de Plus
% Jakarta, pour que les formules se fondent dans le texte.
\usepackage[math-style=ISO,bold-style=ISO]{unicode-math}
\setmathfont{FiraMath-Regular.otf}[Path=../../../fonts/]
\setmathfont{PlusJakartaSans-Medium.ttf}[Path=../../../fonts/,range={up/num,up/{Latin,latin},\mathup}]
\usepackage{icomma} % virgule décimale sans espace en mode math
% Caractères mathématiques Unicode absents des polices de texte (Plus Jakarta,
% Archivo Black) : rendus par leur équivalent mathématique, en texte comme en
% formule — sinon ils s'affichent en carré vide (ex. « Arithmétique dans ℤ »).
\usepackage{newunicodechar}
\newunicodechar{ℝ}{\ensuremath{\mathbb{R}}}
\newunicodechar{ℤ}{\ensuremath{\mathbb{Z}}}
\newunicodechar{ℂ}{\ensuremath{\mathbb{C}}}
\newunicodechar{ℕ}{\ensuremath{\mathbb{N}}}
\newunicodechar{ℚ}{\ensuremath{\mathbb{Q}}}
\newunicodechar{𝔻}{\ensuremath{\mathbb{D}}}
\newunicodechar{α}{\ensuremath{\alpha}}
\newunicodechar{β}{\ensuremath{\beta}}
\newunicodechar{γ}{\ensuremath{\gamma}}
\newunicodechar{δ}{\ensuremath{\delta}}
\newunicodechar{ε}{\ensuremath{\varepsilon}}
\newunicodechar{θ}{\ensuremath{\theta}}
\newunicodechar{λ}{\ensuremath{\lambda}}
\newunicodechar{μ}{\ensuremath{\mu}}
\newunicodechar{σ}{\ensuremath{\sigma}}
\newunicodechar{τ}{\ensuremath{\tau}}
\newunicodechar{φ}{\ensuremath{\varphi}}
\newunicodechar{ψ}{\ensuremath{\psi}}
\newunicodechar{ω}{\ensuremath{\omega}}
\newunicodechar{Σ}{\ensuremath{\Sigma}}
% majuscules produites par \MakeUppercase dans les titres (α-aminés -> Α)
\newunicodechar{Α}{\ensuremath{\alpha}}
\newunicodechar{Β}{\ensuremath{\beta}}
\newunicodechar{Γ}{\ensuremath{\Gamma}}
\newunicodechar{Δ}{\ensuremath{\Delta}}
\newunicodechar{Ε}{\ensuremath{\varepsilon}}
\newunicodechar{Θ}{\ensuremath{\Theta}}
\newunicodechar{Λ}{\ensuremath{\Lambda}}
\newunicodechar{Μ}{\ensuremath{\mu}}
\newunicodechar{Τ}{\ensuremath{\tau}}
\newunicodechar{Φ}{\ensuremath{\Phi}}
\newunicodechar{Ψ}{\ensuremath{\Psi}}
\newunicodechar{Ω}{\ensuremath{\Omega}}
\newunicodechar{↦}{\ensuremath{\mapsto}}
\newunicodechar{∈}{\ensuremath{\in}}
\newunicodechar{∉}{\ensuremath{\notin}}
\newunicodechar{∧}{\ensuremath{\wedge}}
\newunicodechar{⇔}{\ensuremath{\Leftrightarrow}}
\newunicodechar{⟺}{\ensuremath{\Longleftrightarrow}}
\newunicodechar{⇒}{\ensuremath{\Rightarrow}}
\newunicodechar{⟹}{\ensuremath{\Longrightarrow}}
\newunicodechar{∘}{\ensuremath{\circ}}
\newunicodechar{∪}{\ensuremath{\cup}}
\newunicodechar{∩}{\ensuremath{\cap}}
\newunicodechar{≡}{\ensuremath{\equiv}}
\newunicodechar{⊂}{\ensuremath{\subset}}
\newunicodechar{⊥}{\ensuremath{\perp}}
\newunicodechar{∃}{\ensuremath{\exists}}
\newunicodechar{∀}{\ensuremath{\forall}}
\newunicodechar{∛}{\ensuremath{\sqrt[3]{\,}}}
\newunicodechar{∜}{\ensuremath{\sqrt[4]{\,}}}
\newunicodechar{⃗}{\ensuremath{{}^{\to}}}
\newunicodechar{ⁿ}{\ifmmode^{n}\else\textsuperscript{\ensuremath{n}}\fi}
\newunicodechar{ˣ}{\ifmmode^{x}\else\textsuperscript{\ensuremath{x}}\fi}
\newunicodechar{ᵇ}{\ifmmode^{b}\else\textsuperscript{\ensuremath{b}}\fi}
\newunicodechar{ʳ}{\ifmmode^{r}\else\textsuperscript{\ensuremath{r}}\fi}
\newunicodechar{ᵘ}{\ifmmode^{u}\else\textsuperscript{\ensuremath{u}}\fi}
\newunicodechar{ʸ}{\ifmmode^{y}\else\textsuperscript{\ensuremath{y}}\fi}
\newunicodechar{ᵃ}{\ifmmode^{a}\else\textsuperscript{\ensuremath{a}}\fi}
\newunicodechar{ᵉ}{\ifmmode^{e}\else\textsuperscript{\ensuremath{e}}\fi}
\newunicodechar{ᵏ}{\ifmmode^{k}\else\textsuperscript{\ensuremath{k}}\fi}
\newunicodechar{ᵖ}{\ifmmode^{p}\else\textsuperscript{\ensuremath{p}}\fi}
\newunicodechar{ᴺ}{\ifmmode^{N}\else\textsuperscript{\ensuremath{N}}\fi}
\newunicodechar{ᶜ}{\ifmmode^{c}\else\textsuperscript{\ensuremath{c}}\fi}
\newunicodechar{ʰ}{\ifmmode^{h}\else\textsuperscript{\ensuremath{h}}\fi}
\newunicodechar{ᵠ}{\ifmmode^{q}\else\textsuperscript{\ensuremath{q}}\fi}
\newunicodechar{ₙ}{\ifmmode_{n}\else\textsubscript{\ensuremath{n}}\fi}
\newunicodechar{ₐ}{\ifmmode_{a}\else\textsubscript{\ensuremath{a}}\fi}
\newunicodechar{ᵢ}{\ifmmode_{i}\else\textsubscript{\ensuremath{i}}\fi}
\newunicodechar{ᵣ}{\ifmmode_{r}\else\textsubscript{\ensuremath{r}}\fi}
\newunicodechar{ₖ}{\ifmmode_{k}\else\textsubscript{\ensuremath{k}}\fi}
\newunicodechar{ᵦ}{\ifmmode_{\beta}\else\textsubscript{\ensuremath{\beta}}\fi}
\newunicodechar{ₓ}{\ifmmode_{x}\else\textsubscript{\ensuremath{x}}\fi}
\newfontfamily\popdisplay{ArchivoBlack-Regular.ttf}[Path=../../../fonts/]
\newfontfamily\poplogo{SpaceGrotesk-Bold.ttf}[Path=../../../fonts/]
\newfontfamily\popsemi{PlusJakartaSans-SemiBold.ttf}[Path=../../../fonts/]
\newfontfamily\popxbold{PlusJakartaSans-ExtraBold.ttf}[Path=../../../fonts/]
\newfontfamily\popxboldls{PlusJakartaSans-ExtraBold.ttf}[Path=../../../fonts/,LetterSpace=3.0]

\setlist[enumerate]{leftmargin=1.7em,itemsep=5pt,topsep=4pt}
\setlist[itemize]{leftmargin=1.7em,itemsep=4pt,topsep=4pt,label={\color{poporange}\textbullet}}
\setlength{\parindent}{0pt}
\setlength{\parskip}{5pt}
\renewcommand{\arraystretch}{1.25}

% pgfplots : style Pop par défaut
\pgfplotsset{
  every axis/.append style={axis line style={popink,line width=1pt},tick style={popink},
    grid style={popline,line width=0.5pt},label style={font=\small},tick label style={font=\footnotesize}},
  popcurve/.style={poporange,line width=1.8pt,no markers,smooth},
}
% Même style disponible dans un \draw TikZ simple (hors pgfplots)
\tikzset{popcurve/.style={poporange,line width=1.8pt}}
\tikzset{popgrid/.style={popline,very thin}}
\colorlet{popcurve}{poporange} % le nom du style est parfois utilisé comme couleur

% ============================================================
% Pill / squiggle
% ============================================================
\newcommand{\poppill}[3]{%
  \tikz[baseline=-0.55ex]\node[draw=popink,line width=1.2pt,rounded corners=2.6pt,%
    fill=#2,inner xsep=6.5pt,inner ysep=3pt]{%
    \popxboldls\fontsize{7}{7}\selectfont\color{#3}\MakeUppercase{#1}};%
}
\newcommand{\popsquiggle}[1]{%
  \tikz[baseline=-0.4ex]{
    \draw[#1,line width=2.6pt,line cap=round]
      plot [smooth] coordinates {(0,0.09) (0.34,0.01) (0.68,0.21) (1.02,0.01) (1.36,0.10)};
  }%
}
% Mot-clé surligné (marqueur orange sous le texte)
\newcommand{\cle}[1]{{\popxbold\color{popdark}#1}}

% ============================================================
% En-tête / pied
% ============================================================
\pagestyle{fancy}
\fancyhf{}
\renewcommand{\headrulewidth}{0pt}
\renewcommand{\footrulewidth}{0pt}
\lhead{\raisebox{-0.15ex}{\poplogo\fontsize{13}{13}\selectfont\color{popink}OnBuch\color{poporange}+}}
\rhead{\poppill{\DOCMATIERE\ \textbullet\ \DOCNIVEAU\ \textbullet\ Leçon \DOCLECON}{white}{popink}}
\lfoot{\popxbold\fontsize{7.6}{7.6}\selectfont\color{popmuted}\MakeUppercase{Cours signé OnBuch+}\ \textbullet\ {\popsemi\DOCTITRE}}
\rfoot{\tikz[baseline=-0.6ex]\node[rounded corners=8.5pt,fill=popink,minimum height=17pt,minimum width=17pt,inner xsep=5pt,inner sep=0]{\popxbold\fontsize{8}{8}\selectfont\color{popcream}\thepage\,/\,\pageref*{LastPage}};}

% ============================================================
% Titres
% ============================================================
\newcommand{\chaptitre}[2]{%
  \begin{tcolorbox}[enhanced,colback=poporange,colframe=popink,boxrule=2.6pt,arc=11pt,
      drop shadow=popink,left=15pt,right=15pt,top=12pt,bottom=11pt,before skip=3pt,after skip=18pt]
    \poppill{#2}{popink}{popcream}\par\vspace{9pt}
    {\raggedright\hyphenpenalty=10000\exhyphenpenalty=10000\popdisplay\fontsize{22}{24}\selectfont\color{popcream}\MakeUppercase{#1}\par}
  \end{tcolorbox}
}
% Section numérotée : pastille orange avec le numéro + titre Archivo
\newcounter{popsec}
\newcommand{\coursec}[1]{%
  \stepcounter{popsec}\setcounter{popsub}{0}%
  \par\vspace{16pt}\noindent
  \begin{minipage}{\linewidth}
  \tikz[baseline=-0.7ex]\node[draw=popink,line width=1.6pt,fill=poporange,rounded corners=4pt,minimum size=22pt,inner sep=0]{\popdisplay\fontsize{12}{12}\selectfont\color{popcream}\thepopsec};%
  \hspace{8pt}{\popdisplay\fontsize{15}{17}\selectfont\color{popink}\MakeUppercase{#1}}\par
  \vspace{4pt}\noindent{\color{poporange}\rule{36pt}{3.6pt}}
  \end{minipage}\par\nopagebreak\vspace{8pt}
}
\newcounter{popsub}
\newcommand{\courssub}[1]{%
  \stepcounter{popsub}%
  \par\vspace{9pt}\noindent{\popxbold\fontsize{11.5}{13}\selectfont\color{popink}{\color{poporange}\thepopsec.\thepopsub}\hspace{6pt}#1}\par\nopagebreak\vspace{4pt}
}

% ============================================================
% Boîtes pédagogiques — même recette : fond blanc, bordure épaisse colorée,
% ombre dure encre, badge plein. Titre optionnel [#1].
% ============================================================
\tcbset{popbase/.style={breakable,enhanced,colback=white,boxrule=2.3pt,arc=9pt,
  shadow={3pt}{-3pt}{0pt}{popink},left=14pt,right=14pt,top=11pt,bottom=11pt,before skip=11pt,after skip=15pt,
  fontupper=\popsemi\fontsize{10.6}{14.5}\selectfont\color{popink},
  fontlower=\popsemi\fontsize{10.6}{14.5}\selectfont\color{popink}}}
% Coche dessinée (le glyphe ✓ n'existe pas dans les polices du document)
\newcommand{\popcheck}[1]{\tikz[baseline=-0.5ex]\draw[#1,line width=1.6pt,line cap=round,line join=round](-0.13,0.0)--(-0.04,-0.09)--(0.13,0.1);}
\providecommand{\coche}{\popcheck{popgreen}}
\providecommand{\resultat}[1]{\par\smallskip\noindent\fcolorbox{popgreen}{popgreenL}{\parbox{\dimexpr\linewidth-2\fboxsep-2\fboxrule\relax}{\textbf{Résultat.} #1}}\par\smallskip}
\newcommand{\popboxhead}[3]{\poppill{#1}{#2}{white}\ifx\relax#3\relax\else\hspace{6pt}{\popxbold\fontsize{10.6}{12}\selectfont\color{popink}#3}\fi\par\vspace{7pt}}

\newtcolorbox{aretenir}[1][]{popbase,colframe=popgreen,before upper={\popboxhead{À retenir}{popgreen}{#1}}}
% Texte en espagnol (document de lecture, dialogue, poème, extrait) : cadre
% d'étude. Un titre optionnel (source, auteur) s'affiche dans l'en-tête.
\newtcolorbox{textebox}[1][]{popbase,colframe=popblue,before upper={\popboxhead{Texto}{popblue}{#1}\begin{otherlanguage}{spanish}},after upper={\end{otherlanguage}}}
% Marques ✓ / ✗ (forme correcte / fautive) souvent utilisées par les modèles
\providecommand{\cross}{\textcolor{poppink}{\ensuremath{\times}}}
\providecommand{\xmark}{\cross}\providecommand{\crossmark}{\cross}\providecommand{\cmark}{\ensuremath{\checkmark}}
% Ligne de réponse / justification à compléter (inventée par les modèles)
\providecommand{\justification}[1]{\par\noindent\textit{Justification~:} \dotfill\par}
% \textipa (package tipa non chargé) : la police ne couvre pas l'API -> simple passage
\providecommand{\textipa}[1]{#1}
% Soulignement (ulem non chargé) : \uline, \ul
\providecommand{\uline}[1]{\underline{#1}}\providecommand{\ul}[1]{\underline{#1}}
% \glossaire{\item ...} inventée par les modèles : liste de glossaire
\providecommand{\glossaire}[1]{\begin{itemize}#1\end{itemize}}
% \alert (beamer) inventée par les modèles : texte mis en évidence
\providecommand{\alert}[1]{\textbf{\textcolor{poppink}{#1}}}
% variantes ASCII d'ordinaux écrites en macro
\providecommand{\iere}{\textsuperscript{re}}\providecommand{\ieme}{\textsuperscript{e}}\providecommand{\ere}{\textsuperscript{re}}\providecommand{\eme}{\textsuperscript{e}}\providecommand{\er}{\textsuperscript{er}}\providecommand{\ie}{\textsuperscript{e}}
% Ordinaux français écrits en macro par les modèles : 1\ière, 3\ième
\newcommand{\ière}{\textsuperscript{re}}\newcommand{\ième}{\textsuperscript{e}}
% \exemple (inventée par les modèles) : étiquette « Exemple : »
\providecommand{\exemple}{\textbf{Exemple~:}\ }
% Mot ou phrase en espagnol à l'intérieur d'un paragraphe français
\newcommand{\esp}[1]{\ifmmode\text{\textit{#1}}\else\begingroup\selectlanguage{spanish}\itshape #1\endgroup\fi}
\newtcolorbox{exemplebox}[1][]{popbase,colframe=poporange,before upper={\popboxhead{Exemple}{poporange}{#1}}}
\newtcolorbox{definition}[1][]{popbase,colframe=popblue,before upper={\popboxhead{Définition}{popblue}{#1}}}
\newtcolorbox{propriete}[1][]{popbase,colframe=poppurple,before upper={\popboxhead{Propriété}{poppurple}{#1}}}
\newtcolorbox{methode}[1][]{popbase,colframe=popgold,before upper={\popboxhead{Méthode}{popgold}{#1}}}
\newtcolorbox{attention}[1][]{popbase,colframe=poppink,before upper={\popboxhead{Attention}{poppink}{#1}}}
\newtcolorbox{experience}[1][]{popbase,colframe=popblue,colback=popblueL,before upper={\popboxhead{Expérience}{popblue}{#1}}}
% « Le savais-tu ? » : carte violette pleine (contexte, culture, Cameroun)
\newtcolorbox{savaistu}[1][]{popbase,colback=poppurple,colframe=popink,
  fontupper=\popsemi\fontsize{10.4}{14.2}\selectfont\color{white},
  before upper={\poppill{Le savais-tu\,?}{white}{poppurple}\ifx\relax#1\relax\else\hspace{6pt}{\popxbold\color{white}#1}\fi\par\vspace{7pt}}}
% Exercice résolu : énoncé en haut, \tcblower puis la solution sur fond crème
\newcounter{popexo}\newcounter{popexr}
\newtcolorbox{exoresolu}[1][]{popbase,colframe=poporange,colbacklower=popcream,
  segmentation style={popink,line width=1.2pt,dashed},
  before upper={\stepcounter{popexr}\popboxhead{Exercice résolu \thepopexr}{poporange}{#1}},
  before lower={\poppill{Solution}{popink}{popcream}\par\vspace{6pt}}}
% Exercice (non résolu en place) — la correction est dans la partie Corrigés
\newtcolorbox{exercice}[1][]{popbase,colframe=popink,
  before upper={\stepcounter{popexo}\popboxhead{Exercice \thepopexo}{popink}{#1}}}
% Corrigé
\newtcolorbox{corrige}[1][]{popbase,colframe=popgreen,colback=popgreenL,
  before upper={\popboxhead{Corrigé}{popgreen}{#1}}}
% Objectifs / prérequis / bilan
\newtcolorbox{objectifs}{popbase,colframe=popink,colback=poporangeL,
  before upper={\popboxhead{Objectifs de la leçon}{popink}{}}}
\newtcolorbox{prerequis}{popbase,colframe=popink,colback=popblueL,
  before upper={\popboxhead{Prérequis}{popblue}{}}}
\newtcolorbox{bilan}{popbase,colframe=popink,colback=white,boxrule=2.8pt,
  before upper={\popboxhead{Fiche bilan}{poporange}{L'essentiel à connaître pour l'examen}}}
% Figure encadrée avec légende
\newtcolorbox{popfigure}[1][]{enhanced,breakable=false,colback=white,colframe=popline,boxrule=1.4pt,arc=8pt,
  left=10pt,right=10pt,top=10pt,bottom=8pt,before skip=10pt,after skip=12pt,halign=center,
  lower separated=false,halign lower=center,
  fontlower=\popsemi\fontsize{9}{11.5}\selectfont\color{popmuted},
  #1}
\newcommand{\legende}[1]{\par\vspace{6pt}{\popsemi\fontsize{9}{11.5}\selectfont\color{popmuted}#1}}

% ============================================================
% Signature OnBuch+
% ============================================================
\newcommand{\poplogoinline}[1]{{\poplogo\fontsize{#1}{#1}\selectfont\color{popink}OnBuch\color{poporange}+}}
\newcommand{\signatureonbuch}{%
  \begin{tcolorbox}[enhanced,colback=popink,colframe=popink,boxrule=2.2pt,arc=10pt,
      drop shadow=poporange,left=14pt,right=14pt,top=10pt,bottom=10pt,before skip=0pt,after skip=0pt]
    \tikz[baseline=-0.7ex]\node[circle,fill=popgreen,draw=popcream,line width=1.3pt,minimum size=22pt,inner sep=0]{\popcheck{white}};%
    \hspace{9pt}\begin{minipage}[c]{0.62\linewidth}
      {\popxbold\fontsize{12}{13}\selectfont\color{popcream}Signé\ \textbullet\ {\poplogo OnBuch\color{poporange}+}}\par\vspace{2pt}
      {\raggedright\popsemi\fontsize{8}{10.5}\selectfont\color{popline}\MakeUppercase{Équipe pédagogique OnBuch+}\\\MakeUppercase{Conforme au programme officiel MINESEC}\par}
    \end{minipage}\hfill
    {\poplogo\fontsize{22}{22}\selectfont\color{popcream}OnBuch\color{poporange}+}
  \end{tcolorbox}%
}

% ============================================================
% Page de garde
% #1 titre · #2 matière · #3 niveau · #4 module · #5 n° leçon · #6 durée ·
% #7 description / objectifs (texte ou itemize)
% ============================================================
\newcommand{\pagegarde}[7]{%
  \thispagestyle{empty}
  \newgeometry{a4paper,margin=1.7cm,top=1.6cm,bottom=1.4cm}%
  \begin{tikzpicture}[remember picture,overlay]
    \fill[poporange!18] ([xshift=-3.2cm,yshift=-3.6cm]current page.north east) circle (4.6cm);
    \fill[poppurple!14] ([xshift=2.4cm,yshift=7.6cm]current page.south west) circle (3.8cm);
  \end{tikzpicture}%
  {\poplogo\fontsize{30}{30}\selectfont\color{popink}OnBuch\color{poporange}+}\hfill
  \raisebox{0.6ex}{\poppill{Cours signé \textbullet\ Premium}{popink}{popcream}}\par
  \vspace{1pt}\popsquiggle{poppurple}\par
  \vspace{8pt}
  \begin{tcolorbox}[enhanced,colback=poporange,colframe=popink,boxrule=2.8pt,arc=13pt,
      drop shadow=popink,left=18pt,right=18pt,top=12pt,bottom=13pt,after skip=10pt]
    \poppill{#2 \textbullet\ #3}{popink}{popcream}\hspace{5pt}\poppill{#4}{white}{popink}\par\vspace{8pt}
    {\popdisplay\fontsize{14}{14}\selectfont\color{popink}LEÇON #5}\par\vspace{4pt}
    {\raggedright\hyphenpenalty=10000\exhyphenpenalty=10000\popdisplay\fontsize{25}{27}\selectfont\color{popcream}\MakeUppercase{#1}\par}
    \vspace{8pt}
    {\popxbold\fontsize{9.5}{9.5}\selectfont\color{popink}\MakeUppercase{Durée conseillée : #6}}
  \end{tcolorbox}
  \begin{tcolorbox}[enhanced,colback=white,colframe=popink,boxrule=2.2pt,arc=10pt,
      drop shadow=popink,left=16pt,right=16pt,top=10pt,bottom=10pt,after skip=8pt]
    \poppill{Dans cette leçon}{poppurple}{white}\par\vspace{5pt}
    \popsemi\fontsize{9.3}{12.6}\selectfont\color{popink}#7
  \end{tcolorbox}
  \noindent\begin{minipage}[t]{0.487\linewidth}
    \begin{tcolorbox}[enhanced,colback=popgreen,colframe=popink,boxrule=2.2pt,arc=10pt,
        drop shadow=popink,left=11pt,right=11pt,top=10pt,bottom=10pt,height=3.1cm,valign=center]
      \tcbox[colback=white,colframe=popink,boxrule=1.3pt,arc=5pt,boxsep=0pt,nobeforeafter,
        left=3pt,right=3pt,top=3pt,bottom=3pt,valign=center]{\qrcode[height=1.9cm]{https://play.google.com/store/apps/details?id=com.luvvix.onbuch&hl=fr}}%
      \hspace{8pt}\begin{minipage}[c]{0.5\linewidth}
        {\popdisplay\fontsize{11}{12.5}\selectfont\color{white}\MakeUppercase{Télécharge OnBuch}}\par\vspace{4pt}
        {\raggedright\popsemi\fontsize{8.4}{11}\selectfont\color{white}Gratuit \textbullet\ Google Play\\Tuteur IA, annales, concours, résultats\par}
      \end{minipage}
    \end{tcolorbox}
  \end{minipage}\hfill
  \begin{minipage}[t]{0.487\linewidth}
    \begin{tcolorbox}[enhanced,colback=poppurple,colframe=popink,boxrule=2.2pt,arc=10pt,
        drop shadow=popink,left=11pt,right=11pt,top=10pt,bottom=10pt,height=3.1cm,valign=center]
      \tikz[baseline=-0.7ex]\node[circle,fill=white,draw=popink,line width=1.3pt,minimum size=1.6cm,inner sep=0]{\popdisplay\fontsize{24}{24}\selectfont\color{poppurple}+};%
      \hspace{8pt}\begin{minipage}[c]{0.6\linewidth}
        {\popdisplay\fontsize{11}{12.5}\selectfont\color{white}\MakeUppercase{Passe à OnBuch+}}\par\vspace{4pt}
        {\raggedright\popsemi\fontsize{8.4}{11}\selectfont\color{white}Tous les cours signés, Tuteur IA illimité, fiches PDF — onglet OnBuch+ de l'app\par}
      \end{minipage}
    \end{tcolorbox}
  \end{minipage}
  \vspace{8pt}
  \popcontenu
  \vfill
  \signatureonbuch
  \restoregeometry
  \newpage
}

% Rangée « Ce cours contient » (page de garde)
\newcommand{\poptile}[3]{% #1 couleur #2 gros texte #3 légende
  \begin{tcolorbox}[enhanced,colback=#1,colframe=popink,boxrule=2pt,arc=9pt,shadow={2.5pt}{-2.5pt}{0pt}{popink},
      left=6pt,right=6pt,top=9pt,bottom=9pt,halign=center,valign=center,height=1.75cm,width=\linewidth,nobeforeafter]
    {\popdisplay\fontsize{12.5}{13}\selectfont\color{white}\MakeUppercase{#2}}\par\vspace{3pt}
    {\centering\popsemi\fontsize{7.8}{9.5}\selectfont\color{white}#3\par}
  \end{tcolorbox}}
\newcommand{\popcontenu}{%
  \par\noindent{\popxboldls\fontsize{7.5}{7.5}\selectfont\color{popmuted}CE COURS SIGNÉ CONTIENT}\par\vspace{7pt}
  \noindent\begin{minipage}[t]{0.235\linewidth}\poptile{popblue}{Cours}{complet, illustré, pas à pas}\end{minipage}\hfill
  \begin{minipage}[t]{0.235\linewidth}\poptile{poporange}{Méthodes}{et exercices résolus}\end{minipage}\hfill
  \begin{minipage}[t]{0.235\linewidth}\poptile{poppink}{Exercices}{type examen + corrigés}\end{minipage}\hfill
  \begin{minipage}[t]{0.235\linewidth}\poptile{popgreen}{Fiche bilan}{l'essentiel à retenir}\end{minipage}\par}

% Sommaire
\newtcolorbox{sommaire}{enhanced,colback=white,colframe=popink,boxrule=2.2pt,arc=10pt,
  drop shadow=popink,left=18pt,right=18pt,top=13pt,bottom=13pt,before skip=6pt,after skip=20pt,
  before upper={\poppill{Sommaire}{poppurple}{white}\par\vspace{9pt}\popsemi\fontsize{10.6}{14.5}\selectfont\color{popink}}}

% Signature de fin de cours — poussée tout en bas de la dernière page
\newcommand{\findecours}{%
  \par\vfill\nopagebreak
  \begin{center}\popsquiggle{poporange}\end{center}\vspace{4pt}
  \signatureonbuch
}
\AtBeginDocument{\def\llbracket{\mathopen{[\mkern-3.5mu[}}\def\rrbracket{\mathclose{]\mkern-3.5mu]}}} % intervalles d'entiers (glyphe ⟦ absent de la police)
% Glyphes absents de Fira Math : redéfinis à partir de glyphes disponibles
\AtBeginDocument{%
  \def\setminus{\mathbin{\backslash}}\let\smallsetminus\setminus
  \def\vdots{\mathord{\vbox{\baselineskip3.2pt\lineskiplimit0pt\kern4pt\hbox{.}\hbox{.}\hbox{.}}}}%
  \def\ddots{\mathinner{\mkern1mu\raise6.5pt\hbox{.}\mkern2mu\raise3.6pt\hbox{.}\mkern2mu\raise0.7pt\hbox{.}\mkern1mu}}%
  \def\triangle{\mathord{\scalebox{0.9}{$\symup{\Delta}$}}}%
  \def\nabla{\mathord{\raisebox{\depth}{\scalebox{1}[-1]{$\symup{\Delta}$}}}}%
  \def\checkmark{\popcheck{popdark}}%
  \def\star{\mathbin{\ast}}\def\bigstar{\mathord{\ast}}%
  \def\diamond{\mathbin{\scalebox{0.6}{$\Diamond$}}}%
  \def\bigcup{\mathop{\mathchoice{\raisebox{-0.3em}{\scalebox{1.7}{$\cup$}}}{\scalebox{1.25}{$\cup$}}{\cup}{\cup}}\displaylimits}%
  \def\bigcap{\mathop{\mathchoice{\raisebox{-0.3em}{\scalebox{1.7}{$\cap$}}}{\scalebox{1.25}{$\cap$}}{\cap}{\cap}}\displaylimits}%
  \def\lesseqqgtr{\mathrel{\lessgtr}}\def\bigoplus{\mathop{\oplus}\displaylimits}%
}
\providecommand{\ssi}{si et seulement si} % abréviation fréquente
\AtBeginDocument{\providecommand{\wideparen}{\overparen}} % arc de cercle

%% FIGFIX-BEGIN v4 — correctifs globaux des figures (docs/onbuch-generation/tools/figfix)
\usepackage{adjustbox}\usepackage{etoolbox}
% toute figure TikZ/pgfplots plus large que la ligne est réduite à la largeur disponible
\BeforeBeginEnvironment{tikzpicture}{\begin{adjustbox}{max width=\linewidth}}
\AfterEndEnvironment{tikzpicture}{\end{adjustbox}}
% légendes pgfplots lisibles par-dessus les courbes
\pgfplotsset{every axis/.append style={legend style={fill=white,fill opacity=0.92,text opacity=1,draw=black!15,inner sep=3pt}}}
% étiquettes d'arêtes (syntaxe edge["..."]) avec fond blanc
\tikzset{every edge quotes/.append style={fill=white,inner sep=1.2pt}}
%% FIGFIX-END
\begin{document}

\pagegarde{Traduction : les locutions adverbiales ; redacción : ensayo argumentativo sobre el medioambiente}{Espagnol}{Tle A}{Módulo 2 : Environnement, santé et bien-être}{E10}{2 h}{Leçon mixte du Módulo 2 (Environnement, santé et bien-être) : maîtriser les locutions adverbiales espagnoles de temps, lieu, manière et quantité et leurs correspondances françaises, puis construire un essai argumentatif structuré sur une question environnementale, avec connecteurs logiques et arguments exemplifiés.
\vspace{4pt}
\begin{multicols}{2}\small
\begin{itemize}
\item À la fin de cette leçon, je sais identifier et classer les locutions adverbiales espagnoles (temps, lieu, manière, quantité)
\item je sais traduire correctement les locutions adverbiales courantes de l'espagnol vers le français et inversement
\item je sais éviter les faux amis et les calques entre locutions françaises et espagnoles
\item je sais employer les locutions adverbiales dans une phrase et choisir leur place correcte
\item je sais reconnaître la structure d'un essai argumentatif (introduction, thèse, arguments, exemples, conclusion)
\item je sais utiliser les connecteurs logiques d'addition, d'opposition, de cause et de conséquence
\end{itemize}\end{multicols}}

\begin{sommaire}
\begin{multicols}{2}
\begin{enumerate}
\item Texte support et découverte
\item Les locutions adverbiales : classement et sens
\item Correspondances français-espagnol et pièges de traduction
\item L'essai argumentatif : structure et connecteurs
\item Lexique du medioambiente et modèle d'essai
\item Entraînement guidé : thème, version et rédaction
\item Activité d'intégration
\item Méthodes et exercices résolus
\item Exercices
\item Corrigés des exercices
\item Fiche bilan
\end{enumerate}
\end{multicols}
\end{sommaire}

\begin{objectifs}
\begin{itemize}
\item À la fin de cette leçon, je sais identifier et classer les locutions adverbiales espagnoles (temps, lieu, manière, quantité)
\item je sais traduire correctement les locutions adverbiales courantes de l'espagnol vers le français et inversement
\item je sais éviter les faux amis et les calques entre locutions françaises et espagnoles
\item je sais employer les locutions adverbiales dans une phrase et choisir leur place correcte
\item je sais reconnaître la structure d'un essai argumentatif (introduction, thèse, arguments, exemples, conclusion)
\item je sais utiliser les connecteurs logiques d'addition, d'opposition, de cause et de conséquence
\item je sais rédiger un essai argumentatif complet sur un sujet environnemental
\item je sais mobiliser le lexique du medioambiente pour illustrer mes arguments
\end{itemize}
\end{objectifs}

\begin{prerequis}
\begin{itemize}
\item Les adverbes simples en -mente et les adverbes de fréquence (Première)
\item Le lexique de base de l'environnement (naturaleza, contaminación, reciclaje)
\item Les connecteurs simples : y, pero, porque, entonces
\item La structure d'un paragraphe argumentatif en français
\item Le présent de l'indicatif et le futur pour exprimer des opinions et des prévisions
\end{itemize}
\end{prerequis}

\newpage

\coursec{Texte support et découverte}

\courssub{Situation d'accroche et problématique}

Imaginez la scène : il est seize heures à Yaoundé, le ciel s'assombrit \emph{de repente} au-dessus du marché de Mokolo. Les vendeuses de plantains et de piments lèvent la tête ; \emph{en seguida}, elles replient leurs bâches et courent mettre leurs marchandises à l'abri. \emph{Poco a poco}, les rues en pente se transforment en torrents de boue rouge ; \emph{de vez en cuando}, un bendskin slalome entre les flaques. Ces petites expressions figées — \esp{de repente}, \esp{poco a poco}, \esp{a menudo}, \esp{en seguida} — donnent du rythme, de la couleur et de l'authenticité à l'espagnol. Le jury du Baccalauréat les attend dans vos traductions, et leur absence ou leur mauvaise équivalence coûte cher.

Mais savoir traduire ne suffit pas. Face à la fonte du lac Tchad, à la déforestation du bassin du Congo ou à la pollution plastique des plages de Kribi, il faut aussi savoir \emph{défendre une opinion} par écrit, en espagnol, de manière structurée : c'est l'essai argumentatif, exercice central de l'expression écrite en Terminale A.

\textbf{Problématique de la leçon :} comment reconnaître, comprendre et traduire les locutions adverbiales espagnoles, et comment s'en servir pour construire un essai argumentatif convaincant sur l'environnement ?

\courssub{Lecture du texte : \esp{El lago que desaparece}}

Lisez d'abord le texte en silence, puis à voix haute. Ne cherchez pas encore à tout traduire : repérez seulement les petites expressions qui rythment le récit.

\begin{textebox}[El lago que desaparece — texto original]
Poco a poco, el lago Chad se seca. A menudo, los pescadores esperan en vano junto a una orilla que retrocede cada año, y sus redes vuelven vacías con demasiada frecuencia. Hace medio siglo, la cuenca alimentaba a millones de personas; hoy, de vez en cuando, algún niño pregunta a su abuelo dónde estaba el agua, y el anciano señala un horizonte de arena.

En Douala, lejos de allí, el problema tiene otro rostro pero la misma raíz. Los vertederos crecen sin control junto a los barrios húmedos, y cuando llueve, de repente las calles se convierten en ríos de plástico. A propósito o por descuido, miles de bolsas terminan en el Wouri, y en seguida el mar las devuelve a las playas de Kribi.

De repente, un día, comprendimos que el problema del norte era también el nuestro: la escasez de agua, la basura, el clima que cambia sin pedir permiso. Por fin, algunas asociaciones han empezado a concienciar a los jóvenes; sin embargo, las promesas oficiales se cumplen rara vez. De vez en cuando, las autoridades anuncian un plan; sin embargo, poco a poco, todo vuelve a la rutina.

En seguida debemos reaccionar, porque la naturaleza no espera. A menudo se dice que África es la gran reserva verde del planeta; por suerte, todavía es verdad, pero ¿hasta cuándo?
\end{textebox}

\begin{definition}[Glossaire du texte]
\begin{itemize}
\item \esp{la cuenca} : le bassin (hydrographique)
\item \esp{la orilla} : la rive, le bord (du lac, de la mer)
\item \esp{retroceder} : reculer, se retirer
\item \esp{el vertedero} : la décharge, le dépôt d'ordures
\item \esp{la escasez} : la pénurie, le manque
\item \esp{concienciar} : sensibiliser (quelqu'un à un problème)
\item \esp{la red} : le filet (de pêche)
\item \esp{el descuido} : la négligence ; \esp{por descuido} : par négligence
\item \esp{cumplir (una promesa)} : tenir (une promesse)
\item \esp{la reserva verde} : le poumon vert, la réserve écologique
\item \esp{la raíz} : la racine ; \esp{la misma raíz} : la même origine
\item \esp{rara vez} : rarement (locution adverbiale de fréquence)
\end{itemize}
\end{definition}

\begin{savaistu}[Le lac Tchad, une crise majeure]
Le lac Tchad, partagé entre le Tchad, le Niger, le Nigeria et le Cameroun (région de l'Extrême-Nord), a perdu l'essentiel de sa superficie en un demi-siècle, sous l'effet combiné de la sécheresse, des prélèvements d'eau et de la pression démographique. C'est l'une des plus graves crises écologiques d'Afrique centrale, et un sujet de plus en plus fréquent dans les textes du Bac : retenez le vocabulaire \esp{la sequía} (la sécheresse), \esp{el cambio climático} (le changement climatique), \esp{los refugiados climáticos} (les réfugiés climatiques).
\end{savaistu}

\courssub{Observation : repérage des locutions adverbiales}

Relisez le texte et soulignez toutes les expressions figées qui fonctionnent comme un adverbe. Vous devriez en trouver neuf, souvent répétées :

\begin{center}
\esp{\textbf{poco a poco} (×2) \quad \textbf{a menudo} (×2) \quad \textbf{de vez en cuando} (×2) \quad \textbf{de repente} (×2) \quad \textbf{a propósito} \quad \textbf{en seguida} (×2) \quad \textbf{por fin} \quad \textbf{sin embargo} (×2) \quad \textbf{por suerte} \quad \textbf{rara vez}}
\end{center}

\begin{definition}[Qu'est-ce qu'une locution adverbiale ?]
Une \cle{locution adverbiale} (\esp{locución adverbial}) est un \textbf{groupe de mots figé} qui joue, à lui tout seul, le rôle d'un adverbe : il modifie le verbe, l'adjectif ou la phrase entière. Ainsi \esp{de repente} équivaut à l'adverbe \esp{repentinamente} (= soudain), et \esp{a menudo} équivaut à \esp{frecuentemente} (= souvent). La locution est \emph{invariable} : on ne peut ni changer l'ordre des mots, ni mettre un mot au pluriel (\esp{de repentes} n'existe pas). Remarque : \esp{sin embargo} est, à proprement parler, une \emph{locution conjonctive} (connecteur d'opposition) ; on l'étudie ici avec les locutions adverbiales car elle se traduit elle aussi en bloc et joue un rôle essentiel dans l'argumentation.
\end{definition}

\begin{methode}[Comprendre une locution inconnue : la méthode du contexte]
\begin{enumerate}
\item \textbf{Encadrez} la locution : lisez la phrase entière, puis les phrases qui précèdent et qui suivent.
\item \textbf{Interrogez le verbe} : la locution répond-elle à \emph{quand ?} (temps), \emph{où ?} (lieu), \emph{comment ?} (manière), \emph{combien de fois ?} (fréquence) ?
\item \textbf{Remplacez} mentalement par un adverbe simple que vous connaissez (\esp{rápidamente}, \esp{finalmente}...) et vérifiez que le sens de la phrase tient toujours.
\item \textbf{Formulez une hypothèse} de traduction française, puis seulement après, vérifiez au dictionnaire.
\item \textbf{Testez} votre équivalent français dans la phrase traduite : s'il produit une phrase bizarre, recommencez.
\end{enumerate}
\end{methode}

\courssub{Questions de compréhension globale}

\begin{enumerate}
\item \textbf{Thème :} de quoi parle le texte ? Quels sont les deux espaces géographiques évoqués et quel problème commun les relie ?
\item \textbf{Problème posé :} quelle contradiction l'auteur dénonce-t-il entre les paroles et les actes ? Relevez la phrase qui la résume.
\item \textbf{Point de vue de l'auteur :} le narrateur est-il optimiste, pessimiste, ou les deux ? Justifiez avec deux expressions du texte (\esp{por suerte}, \esp{¿hasta cuándo?}).
\item \textbf{Genre du texte :} s'agit-il d'un récit, d'une description ou d'un texte argumentatif ? Quels indices vous le montrent (le \esp{comprendimos} à la première personne du pluriel, les questions rhétoriques, la conclusion incitant à agir) ?
\end{enumerate}

\courssub{Questions de détail : le sens de chaque locution dans son contexte}

Pour chaque locution, demandez-vous : \emph{que signifie-t-elle ici, dans CETTE phrase ?}

\begin{itemize}
\item \esp{Poco a poco, el lago Chad se seca} : le dessèchement est-il brutal ou progressif ?
\item \esp{A menudo, los pescadores esperan en vano} : s'agit-il d'un fait rare ou habituel ?
\item \esp{De repente, un día, comprendimos...} : la prise de conscience est-elle lente ou soudaine ?
\item \esp{De vez en cuando, las autoridades anuncian un plan} : les annonces sont-elles régulières ou occasionnelles ?
\item \esp{A propósito o por descuido} : les sacs arrivent-ils au fleuve volontairement ou accidentellement ?
\item \esp{En seguida debemos reaccionar} : faut-il agir plus tard ou tout de suite ?
\item \esp{Por fin, algunas asociaciones han empezado...} : ce début d'action est-il attendu depuis longtemps ?
\item \esp{Sin embargo, todo vuelve a la rutina} : la locution introduit-elle une conséquence ou une opposition ?
\end{itemize}

\courssub{Hypothèses de traduction : tableau des équivalents français}

Appliquez la méthode du contexte et complétez mentalement ce tableau avant de regarder la colonne « Équivalent français ».

\begin{center}
\begin{tabularx}{\linewidth}{|l|X|X|}
\hline
\rowcolor{popblueL} \textbf{Locución} & \textbf{Équivalent français} & \textbf{Valeur / piège} \\
\hline
\esp{poco a poco} & peu à peu, petit à petit & manière : exprime la progression lente, pas la quantité \\
\hline
\esp{a menudo} & souvent, fréquemment & fréquence ; jamais « au menu » (faux ami graphique !) \\
\hline
\esp{de repente} & soudain, tout à coup & temps ; invariable, pas de pluriel \\
\hline
\esp{de vez en cuando} & de temps en temps, parfois & fréquence faible ; ne pas confondre avec \esp{a veces} (parfois, plus fréquent) \\
\hline
\esp{a propósito} & exprès, volontairement & manière ; attention : \esp{a propósito de} + nom = « à propos de » \\
\hline
\esp{en seguida} & aussitôt, tout de suite & temps ; s'écrit aussi \esp{enseguida} en un seul mot \\
\hline
\esp{por fin} & enfin, finalement & marque le soulagement après une attente \\
\hline
\esp{sin embargo} & cependant, pourtant & connecteur d'opposition ; équivalent de \esp{no obstante} \\
\hline
\esp{por suerte} & heureusement, par chance & évaluation ; contraire : \esp{por desgracia} \\
\hline
\esp{rara vez} & rarement & fréquence ; variante : \esp{raras veces} \\
\hline
\end{tabularx}
\end{center}

\begin{attention}[Pièges du francophone]
\begin{itemize}
\item \esp{A menudo} ne veut pas dire « au menu » : c'est un faux ami purement graphique. Il signifie « souvent ».
\item \esp{En seguida} ne signifie pas « dans la suite (du texte) » mais « immédiatement ». Pour « dans la suite », dites \esp{a continuación}.
\item \esp{A propósito} seul = « exprès » ; \esp{a propósito de} + nom = « à propos de ». Le contexte tranche.
\item Ne traduisez jamais une locution mot à mot : \esp{de vez en cuando} (« de fois en quand ») n'a de sens qu'en bloc = « de temps en temps ».
\item Attention à l'accent : \esp{a propósito} prend un accent sur le \emph{o} ; \esp{proposito} sans accent est une faute.
\end{itemize}
\end{attention}

\courssub{Carte mentale : classer les locutions du texte}

\begin{popfigure}
\begin{tikzpicture}[
  mindmap,
  grow cyclic,
  every node/.style={concept, align=center, font=\small},
  root concept/.append style={concept color=popblue, font=\small\bfseries},
  level 1/.append style={level distance=3.2cm, sibling angle=90},
  level 2/.append style={level distance=2.6cm, sibling angle=45},
  text=white
]
\node[root concept, align=center]{Locuciones\\adverbiales}
  child[concept color=popgreen]{ node[align=center]{tiempo\\(¿cuándo?)}
    child { node[concept color=popgreen!70, align=center]{de repente\\en seguida} }
    child { node[concept color=popgreen!70]{por fin} }
  }
  child[concept color=poporange]{ node[align=center]{frecuencia\\(¿con qué\\frecuencia?)}
    child { node[concept color=poporange!70, align=center]{a menudo\\rara vez} }
    child { node[concept color=poporange!70, align=center]{de vez\\en cuando} }
  }
  child[concept color=poppurple]{ node[align=center]{modo\\(¿cómo?)}
    child { node[concept color=poppurple!70]{poco a poco} }
    child { node[concept color=poppurple!70, align=center]{a propósito\\por suerte} }
  }
  child[concept color=poppink]{ node[align=center]{conexión\\(oposición)}
    child { node[concept color=poppink!70]{sin embargo} }
  };
\end{tikzpicture}
\legende{Carte mentale des locutions adverbiales du texte, classées selon la question à laquelle elles répondent.}
\end{popfigure}

\begin{exoresolu}[Traduire une phrase à locutions]
Énoncé : traduisez en français la phrase \esp{De vez en cuando, las autoridades prometen actuar; sin embargo, a propósito o no, nada cambia.}
\tcblower
Solution détaillée, étape par étape :
\begin{enumerate}
\item \textbf{Découpage :} je repère trois locutions : \esp{de vez en cuando} (fréquence), \esp{sin embargo} (opposition), \esp{a propósito} (manière).
\item \textbf{Contexte :} les autorités promettent, mais rien ne change ; la phrase oppose une promesse à une réalité.
\item \textbf{Équivalents :} \esp{de vez en cuando} = « de temps en temps » ; \esp{sin embargo} = « cependant » ; \esp{a propósito o no} = « exprès ou non ».
\item \textbf{Traduction finale :} « De temps en temps, les autorités promettent d'agir ; cependant, exprès ou non, rien ne change. »
\item \textbf{Vérification :} la phrase française est naturelle, l'opposition est bien rendue par « cependant », et aucune locution n'a été traduite mot à mot.
\end{enumerate}
\end{exoresolu}

\begin{exercice}[À vous de jouer]
Traduisez en français, en appliquant la méthode du contexte :
\begin{enumerate}
\item \esp{Poco a poco, las redes vuelven vacías y los pescadores pierden la esperanza.}
\item \esp{Cuando llueve en Douala, de repente las calles se convierten en ríos de plástico.}
\item \esp{En seguida debemos reaccionar, porque la naturaleza no espera.}
\end{enumerate}
Puis, à l'oral, reformulez chaque phrase espagnole en remplaçant la locution par un adverbe simple (\esp{gradualmente}, \esp{repentinamente}, \esp{inmediatamente}).
\end{exercice}

\begin{corrige}[Exercice « À vous de jouer »]
\begin{enumerate}
\item « Peu à peu, les filets reviennent vides et les pêcheurs perdent espoir. » Reformulation : \esp{Gradualmente, las redes vuelven vacías y los pescadores pierden la esperanza.}
\item « Quand il pleut à Douala, soudain les rues se transforment en rivières de plastique. » Reformulation : \esp{Cuando llueve en Douala, las calles se convierten repentinamente en ríos de plástico.}
\item « Nous devons réagir tout de suite, car la nature n'attend pas. » Reformulation : \esp{Debemos reaccionar inmediatamente, porque la naturaleza no espera.}
\end{enumerate}
\end{corrige}

\begin{aretenir}[L'essentiel de cette première étape]
Une locution adverbiale est un bloc figé, invariable, qui se traduit \emph{en bloc} et jamais mot à mot. Pour la comprendre, on s'appuie d'abord sur le contexte et sur la question à laquelle elle répond (quand ? comment ? à quelle fréquence ?). Les locutions du texte — \esp{poco a poco}, \esp{a menudo}, \esp{de repente}, \esp{de vez en cuando}, \esp{a propósito}, \esp{en seguida}, \esp{por fin}, \esp{sin embargo}, \esp{por suerte}, \esp{rara vez} — sont des classiques du Bac : apprenez-les avec leur équivalent français et un exemple. Dans la prochaine section, nous les classerons systématiquement par catégorie avant de les réinvestir dans l'essai argumentatif.
\end{aretenir}

\coursec{Les locutions adverbiales : classement et sens}

Dans la section précédente, nous avons lu un texte sur la déforestation et repéré des expressions comme \esp{poco a poco}, \esp{de vez en cuando} ou \esp{hoy en día}. Ces petits groupes de mots, que l'on ne peut pas traduire mot à mot, s'appellent des \cle{locutions adverbiales} (\esp{locuciones adverbiales}). Elles sont partout : dans la presse, dans les conversations, dans les essais. Les maîtriser, c'est à la fois mieux comprendre les textes et enrichir immédiatement ses propres productions. Cette section les classe par sens, donne leurs équivalents français et signale les pièges à éviter.

\courssub{Qu'est-ce qu'une locution adverbiale ?}

\begin{definition}[Locución adverbial]
Une \cle{locution adverbiale} est un groupe figé de plusieurs mots qui joue, ensemble, le rôle d'un seul adverbe : elle précise le temps, le lieu, la manière ou la quantité. Elle se construit le plus souvent avec une préposition (\esp{a}, \esp{de}, \esp{en}, \esp{por}, \esp{sin}) suivie d'un nom, d'un adverbe ou d'une autre structure.
\end{definition}

Observez les paires suivantes :

\begin{itemize}
\item \esp{Poco a poco, el bosque desaparece.} « Petit à petit, la forêt disparaît. »
\item \esp{A menudo llueve en la región del Litoral.} « Il pleut souvent dans la région du Littoral. »
\item \esp{Lo hizo a propósito.} « Il l'a fait exprès. »
\end{itemize}

Dans chaque phrase, on pourrait remplacer la locution par un adverbe simple : \esp{poco a poco} $\rightarrow$ \esp{gradualmente} ; \esp{a menudo} $\rightarrow$ \esp{frecuentemente} ; \esp{a propósito} $\rightarrow$ \esp{intencionadamente}. La locution vaut donc un adverbe, mais elle est plus expressive et plus naturelle à l'oral comme à l'écrit.

\begin{propriete}[Invariabilité de la locution]
Une locution adverbiale est \cle{invariable} et \cle{figée} : on ne change ni l'article, ni la préposition, ni le nombre des mots qui la composent.
\begin{itemize}
\item On dit \esp{de vez en cuando} et jamais « de vez en cuando\emph{s} » ni « de una vez en cuando ».
\item On dit \esp{a menudo} et jamais « al menudo » ni « a menudos ».
\item On dit \esp{en voz alta} et jamais « en una voz alta ».
\end{itemize}
Toute modification détruit la locution ou change son sens. Il faut donc les apprendre comme des blocs, exactement comme du vocabulaire.
\end{propriete}

\begin{attention}[Piège du mot à mot]
On ne traduit jamais une locution adverbiale mot à mot. \esp{A pie} ne renvoie pas au sens anatomique du mot « pied » mais au déplacement : \esp{Voy al liceo a pie} = « Je vais au lycée à pied » ; \esp{de sobra} ne veut pas dire « de l'ombre » mais « en trop, largement ». Cherchez toujours l'équivalent global, pas la somme des mots.
\end{attention}

\courssub{Les locutions de temps}

Elles répondent à la question \esp{¿cuándo?} (quand ?) ou \esp{¿con qué frecuencia?} (à quelle fréquence ?). Ce sont les plus fréquentes dans les textes d'actualité sur l'environnement.

\begin{center}
\begin{tabularx}{\linewidth}{|l|X|X|}
\hline
\rowcolor{popblueL} \textbf{Locución} & \textbf{Équivalent français} & \textbf{Ejemplo y traducción} \\
\hline
\esp{de repente} & soudain, tout à coup & \esp{De repente, empezó a llover.} « Soudain, il s'est mis à pleuvoir. » \\
\hline
\esp{de pronto} & soudain, tout à coup & \esp{De pronto, el río se desbordó.} « Soudain, la rivière a débordé. » \\
\hline
\esp{a menudo} & souvent & \esp{A menudo hablamos del clima.} « Nous parlons souvent du climat. » \\
\hline
\esp{a veces} & parfois & \esp{A veces reciclamos el plástico.} « Parfois nous recyclons le plastique. » \\
\hline
\esp{de vez en cuando} & de temps en temps & \esp{De vez en cuando planto árboles.} « De temps en temps, je plante des arbres. » \\
\hline
\esp{en seguida} & tout de suite, aussitôt & \esp{En seguida vinieron los bomberos.} « Les pompiers sont arrivés aussitôt. » \\
\hline
\esp{por fin} & enfin & \esp{Por fin llueve en el Norte.} « Enfin il pleut dans le Nord. » \\
\hline
\esp{al mismo tiempo} & en même temps & \esp{Al mismo tiempo, suben las temperaturas.} « En même temps, les températures montent. » \\
\hline
\esp{hoy en día} & aujourd'hui, de nos jours & \esp{Hoy en día, el planeta está en peligro.} « De nos jours, la planète est en danger. » \\
\hline
\esp{cada vez más / menos} & de plus en plus / de moins en moins & \esp{Hace cada vez más calor.} « Il fait de plus en plus chaud. » \\
\hline
\end{tabularx}
\end{center}

\begin{attention}[De pronto, de repente et les faux amis]
\begin{itemize}
\item \esp{De pronto} et \esp{de repente} sont \cle{synonymes} : tous deux signifient « soudain ». Ne cherchez pas de différence de sens.
\item \esp{En seguida} signifie « tout de suite » (délai très court), et non « ensuite ». Pour « ensuite », on dit \esp{luego} ou \esp{después}.
\item \esp{A veces} (parfois) s'écrit en deux mots ; ne pas confondre avec \esp{a ver} (voyons) ni avec \esp{avezés}, qui n'existe pas.
\item \esp{Hoy en día} s'écrit en trois mots, sans trait d'union ; on peut dire aussi \esp{hoy día} ou \esp{actualmente}.
\end{itemize}
\end{attention}

\courssub{Les locutions de lieu}

Elles répondent à la question \esp{¿dónde?} (où ?).

\begin{center}
\begin{tabularx}{\linewidth}{|l|X|X|}
\hline
\rowcolor{popgreenL} \textbf{Locución} & \textbf{Équivalent français} & \textbf{Ejemplo y traducción} \\
\hline
\esp{a lo lejos} & au loin & \esp{A lo lejos se ve el monte Camerún.} « Au loin, on voit le mont Cameroun. » \\
\hline
\esp{por todas partes} & partout & \esp{Hay basura por todas partes.} « Il y a des déchets partout. » \\
\hline
\esp{en ninguna parte} & nulle part & \esp{No encuentro mi cuaderno en ninguna parte.} « Je ne trouve mon cahier nulle part. » \\
\hline
\esp{de arriba abajo} & de haut en bas & \esp{Recorrieron el río de arriba abajo.} « Ils ont parcouru la rivière de haut en bas. » \\
\hline
\esp{aquí mismo} & ici même & \esp{Aquí mismo nació este árbol.} « Cet arbre est né ici même. » \\
\hline
\esp{alrededor de} & autour de & \esp{Hay flores alrededor de la escuela.} « Il y a des fleurs autour de l'école. » \\
\hline
\end{tabularx}
\end{center}

Notez la double négation espagnole : \esp{No encuentro nada en ninguna parte} (je ne trouve rien nulle part). Contrairement au français, l'espagnol cumule les négations sans faute : \esp{no} + \esp{nada} + \esp{ninguna parte} est la tournure correcte et obligatoire.

\courssub{Les locutions de manière}

Elles répondent à la question \esp{¿cómo?} (comment ?). Ce sont les championnes de la conversation quotidienne.

\begin{center}
\begin{tabularx}{\linewidth}{|l|X|X|}
\hline
\rowcolor{poporangeL} \textbf{Locución} & \textbf{Équivalent français} & \textbf{Ejemplo y traducción} \\
\hline
\esp{poco a poco} & petit à petit & \esp{Poco a poco aprendo español.} « Petit à petit, j'apprends l'espagnol. » \\
\hline
\esp{a propósito} & exprès ; à propos & \esp{Lo rompió a propósito.} « Il l'a cassé exprès. » \\
\hline
\esp{sin querer} & sans le vouloir, par inadvertance & \esp{Sin querer, tiré el papel al suelo.} « Sans le vouloir, j'ai jeté le papier par terre. » \\
\hline
\esp{a pie} & à pied & \esp{Voy al mercado a pie.} « Je vais au marché à pied. » \\
\hline
\esp{de memoria} & par cœur & \esp{Se sabe el poema de memoria.} « Il connaît le poème par cœur. » \\
\hline
\esp{en voz alta} & à voix haute & \esp{Lee el texto en voz alta.} « Lis le texte à voix haute. » \\
\hline
\esp{en voz baja} & à voix basse & \esp{Hablaban en voz baja.} « Ils parlaient à voix basse. » \\
\hline
\esp{a ciegas} & à l'aveuglette & \esp{No compres a ciegas.} « N'achète pas à l'aveuglette. » \\
\hline
\esp{de verdad} & vraiment & \esp{¿De verdad cuidas el agua?} « Tu économises vraiment l'eau ? » \\
\hline
\end{tabularx}
\end{center}

\begin{attention}[A propósito : le faux ami classique]
\esp{A propósito} ne signifie jamais « à proposition ». Selon le contexte, il a deux sens :
\begin{itemize}
\item « exprès, intentionnellement » : \esp{Lo hizo a propósito.} « Il l'a fait exprès. » (contraire : \esp{sin querer}).
\item « à propos, au fait » (pour introduire un sujet) : \esp{A propósito, ¿vienes mañana?} « À propos, tu viens demain ? »
\end{itemize}
De même, \esp{de memoria} = « par cœur » (et non « de mémoire » au sens de « de souvenir » : « de mémoire d'homme » se dit \esp{que se recuerde}).
\end{attention}

\courssub{Les locutions de quantité et de degré}

Elles répondent à la question \esp{¿cuánto?} (combien ?) ou nuancent l'intensité d'une affirmation.

\begin{center}
\begin{tabularx}{\linewidth}{|l|X|X|}
\hline
\rowcolor{popgoldL} \textbf{Locución} & \textbf{Équivalent français} & \textbf{Ejemplo y traducción} \\
\hline
\esp{por lo menos} & au moins & \esp{Debemos plantar por lo menos diez árboles.} « Nous devons planter au moins dix arbres. » \\
\hline
\esp{más o menos} & plus ou moins, à peu près & \esp{La clase dura más o menos dos horas.} « Le cours dure à peu près deux heures. » \\
\hline
\esp{en absoluto} & absolument pas & \esp{No estoy de acuerdo en absoluto.} « Je ne suis absolument pas d'accord. » \\
\hline
\esp{por supuesto} & bien sûr & \esp{—¿Reciclas? —Por supuesto.} « — Tu recycles ? — Bien sûr. » \\
\hline
\esp{de sobra} & largement, en trop & \esp{Sabemos de sobra que el agua es vida.} « Nous savons pertinemment que l'eau, c'est la vie. » \\
\hline
\esp{apenas} & à peine & \esp{Apenas llueve en el desierto.} « Il pleut à peine dans le désert. » \\
\hline
\end{tabularx}
\end{center}

\begin{attention}[En absoluto et apenas]
\begin{itemize}
\item \esp{En absoluto} s'emploie presque toujours dans une phrase \cle{négative} : \esp{No contamina en absoluto} (il ne pollue absolument pas). Pour répondre « absolument ! » de façon affirmative, on dit \esp{¡Desde luego!} ou \esp{¡Por supuesto!}, jamais « ¡en absoluto! ».
\item \esp{Apenas} signifie « à peine » : \esp{Apenas tiene agua} = « il a à peine de l'eau ». Ne pas confondre avec le verbe français « apprendre », qui se dit \esp{aprender} !
\end{itemize}
\end{attention}

\courssub{Les locutions d'attitude et d'opinion : les alliées de l'essai}

Ces locutions expriment le point de vue du locuteur. Elles seront précieuses dans l'essai argumentatif de la section 5, car elles permettent de marquer son adhésion, son regret ou sa certitude.

\begin{center}
\begin{tabularx}{\linewidth}{|l|X|X|}
\hline
\rowcolor{poppurpleL} \textbf{Locución} & \textbf{Équivalent français} & \textbf{Emploi dans l'essai} \\
\hline
\esp{sin duda} & sans doute & \esp{Sin duda, el reciclaje es necesario.} « Sans doute le recyclage est-il nécessaire. » \\
\hline
\esp{por supuesto} & bien sûr & \esp{Por supuesto, debemos actuar ya.} « Bien sûr, nous devons agir maintenant. » \\
\hline
\esp{en efecto} & en effet & \esp{En efecto, el clima ha cambiado.} « En effet, le climat a changé. » \\
\hline
\esp{desde luego} & évidemment, certes & \esp{Desde luego, no es suficiente.} « Certes, ce n'est pas suffisant. » \\
\hline
\esp{por desgracia} & malheureusement & \esp{Por desgracia, los bosques desaparecen.} « Malheureusement, les forêts disparaissent. » \\
\hline
\esp{afortunadamente} & heureusement & \esp{Afortunadamente, hay soluciones.} « Heureusement, il existe des solutions. » \\
\hline
\end{tabularx}
\end{center}

\begin{exemplebox}[Un paragraphe d'essai riche en locutions]
\esp{Hoy en día, la contaminación avanza poco a poco por todas partes. Sin duda, el problema es grave; en efecto, cada vez más ciudades sufren el smog. Por desgracia, muchos ciudadanos tiran basura sin querer o a propósito. Afortunadamente, de vez en cuando surgen iniciativas ecológicas que, por lo menos, nos dan esperanza.}\par
« De nos jours, la pollution avance petit à petit partout. Sans doute le problème est-il grave ; en effet, de plus en plus de villes souffrent du smog. Malheureusement, beaucoup de citoyens jettent des déchets sans le vouloir ou exprès. Heureusement, de temps en temps naissent des initiatives écologiques qui, au moins, nous donnent de l'espoir. »
\end{exemplebox}

\courssub{Carte mentale récapitulative}

\begin{popfigure}
\begin{center}
\begin{tikzpicture}[mindmap, grow cyclic, every node/.style={concept, align=center, font=\small}, concept color=popblue!40, level 1/.append style={level distance=3.4cm, sibling angle=72}, level 2/.append style={level distance=2.6cm, sibling angle=45, font=\scriptsize}]
\node[concept color=popblue!60, font=\small\bfseries] {Locuciones adverbiales}
  child[concept color=popgreen!40] { node {Tiempo}
    child { node[align=center]{de repente\\en seguida} }
    child { node[align=center]{a menudo\\de vez en cuando} }
    child { node[align=center]{hoy en día\\por fin} } }
  child[concept color=poporange!40] { node {Lugar}
    child { node[align=center]{a lo lejos\\aquí mismo} }
    child { node[align=center]{por todas partes\\alrededor de} } }
  child[concept color=poppurple!40] { node {Manera}
    child { node[align=center]{poco a poco\\a pie} }
    child { node[align=center]{a propósito\\sin querer} }
    child { node[align=center]{en voz alta\\de memoria} } }
  child[concept color=popgold!50] { node {Cantidad}
    child { node[align=center]{por lo menos\\más o menos} }
    child { node[align=center]{en absoluto\\apenas} } }
  child[concept color=poppink!40] { node {Opinión}
    child { node[align=center]{sin duda\\en efecto} }
    child { node[align=center]{por desgracia\\afortunadamente} } };
\end{tikzpicture}
\end{center}
\legende{Carte mentale des locutions adverbiales : cinq familles de sens à mémoriser comme des blocs invariables.}
\end{popfigure}

\courssub{La place de la locution dans la phrase}

\begin{propriete}[Position et ponctuation]
La locution adverbiale occupe une place \cle{libre} : début, milieu ou fin de phrase.
\begin{itemize}
\item Début : \esp{De vez en cuando, visito a mis abuelos.} « De temps en temps, je rends visite à mes grands-parents. »
\item Milieu : \esp{Yo de vez en cuando visito a mis abuelos.}
\item Fin : \esp{Visito a mis abuelos de vez en cuando.}
\end{itemize}
Règle d'orthographe : quand la locution est \cle{en tête de phrase}, on met une \cle{virgule} après elle, comme en français : \esp{Poco a poco, el bosque desaparece.} En milieu de phrase, elle peut être encadrée de virgules si elle constitue une incise.
\end{propriete}

En tête de phrase, la locution sert souvent de \cle{connecteur} : elle accroche l'attention du lecteur et annonce le ton. C'est une excellente habitude à prendre pour varier les débuts de phrases dans l'essai : \esp{Por desgracia, ...} ; \esp{Sin duda, ...} ; \esp{Hoy en día, ...}

\begin{textebox}[Mini-dialogue : au marché de Yaoundé]
— ¿Vas a menudo al mercado?\\
— No, solo de vez en cuando, pero en seguida encuentro lo que busco.\\
— ¿Y vas a pie o en moto?\\
— Normalmente a pie; poco a poco hago el camino y, a lo lejos, ya veo los puestos de fruta.\\
— ¿Compras sin pensar?\\
— ¡En absoluto! Miro los precios por todas partes y, por lo menos, comparo tres puestos. A propósito, mañana hay un mercado ecológico aquí mismo, alrededor de la escuela.\\
— ¡Por fin! Sin duda, iré temprano.\par
\emph{Glossaire :} \esp{el puesto} = l'étal, le stand ; \esp{comparar} = comparer ; \esp{temprano} = tôt ; \esp{sin pensar} = sans réfléchir.
\end{textebox}

Ce dialogue concentre douze locutions des cinq familles. Relevez-les, classez-les, puis remplacez chacune par un adverbe simple quand c'est possible : \esp{a menudo} $\rightarrow$ \esp{frecuentemente} ; \esp{en seguida} $\rightarrow$ \esp{inmediatamente}.

\begin{exoresolu}[Thème : traduisez en espagnol]
Énoncé : Traduisez en espagnol : « De nos jours, la forêt disparaît petit à petit. Heureusement, les élèves plantent des arbres de temps en temps, et ils le font parfois à pied, sans le vouloir... non, exprès ! »
\tcblower
Solution détaillée :
\begin{itemize}
\item « De nos jours » $\rightarrow$ \esp{Hoy en día} (locution de temps, en tête, suivie d'une virgule).
\item « petit à petit » $\rightarrow$ \esp{poco a poco} (locution de manière, invariable).
\item « Heureusement » $\rightarrow$ \esp{Afortunadamente} (locution d'opinion).
\item « de temps en temps » $\rightarrow$ \esp{de vez en cuando} (jamais « cuando\emph{s} »).
\item « parfois » $\rightarrow$ \esp{a veces} ; « à pied » $\rightarrow$ \esp{a pie} ; « sans le vouloir » $\rightarrow$ \esp{sin querer} ; « exprès » $\rightarrow$ \esp{a propósito}.
\end{itemize}
Traduction complète : \esp{Hoy en día, el bosque desaparece poco a poco. Afortunadamente, los alumnos plantan árboles de vez en cuando, y a veces lo hacen a pie, ¡sin querer... no, a propósito!}
\end{exoresolu}

\begin{exercice}[À vous de jouer]
\begin{enumerate}
\item Classez les locutions suivantes selon leur famille (temps, lieu, manière, quantité, opinion) : \esp{en seguida} ; \esp{a ciegas} ; \esp{por todas partes} ; \esp{sin duda} ; \esp{cada vez más} ; \esp{de memoria} ; \esp{apenas} ; \esp{a lo lejos}.
\item Traduisez en français : \esp{Por desgracia, apenas llueve; sin embargo, poco a poco la gente aprende a ahorrar agua.}
\item Traduisez en espagnol : « Il pleut souvent ici même, mais soudain tout s'arrête ; heureusement, nous savons pertinemment quoi faire. »
\end{enumerate}
\end{exercice}

\begin{corrige}[Exercice « À vous de jouer »]
\begin{enumerate}
\item Temps : \esp{en seguida}, \esp{cada vez más}. Lieu : \esp{por todas partes}, \esp{a lo lejos}. Manière : \esp{a ciegas}, \esp{de memoria}. Quantité : \esp{apenas}. Opinion : \esp{sin duda}.
\item « Malheureusement, il pleut à peine ; cependant, petit à petit, les gens apprennent à économiser l'eau. »
\item \esp{A menudo llueve aquí mismo, pero de repente todo se detiene; afortunadamente, sabemos de sobra qué hacer.}
\end{enumerate}
\end{corrige}

\begin{aretenir}[L'essentiel de la section]
\begin{itemize}
\item Une locution adverbiale est un bloc \cle{invariable} qui vaut un adverbe : on l'apprend entière, sans jamais changer ses mots.
\item Cinq familles : temps (\esp{de repente}, \esp{a menudo}, \esp{hoy en día}), lieu (\esp{a lo lejos}, \esp{por todas partes}), manière (\esp{poco a poco}, \esp{a propósito}, \esp{sin querer}), quantité (\esp{por lo menos}, \esp{en absoluto}, \esp{apenas}), opinion (\esp{sin duda}, \esp{por desgracia}, \esp{afortunadamente}).
\item Place libre dans la phrase ; en tête, elle est suivie d'une \cle{virgule}.
\item Faux amis à retenir : \esp{a propósito} = exprès / à propos ; \esp{en seguida} = tout de suite ; \esp{de sobra} = largement ; \esp{apenas} = à peine.
\end{itemize}
\end{aretenir}

Maintenant que les locutions sont classées et comprises, la section suivante confrontera systématiquement les correspondances entre le français et l'espagnol, afin de traduire sans calquer.

\coursec{Correspondances français-espagnol et pièges de traduction}

Après avoir classé les locutions adverbiales selon leur sens (temps, lieu, manière, quantité) dans la section précédente, nous abordons maintenant l'étape décisive pour le bachelier : \textbf{traduire fidèlement ces blocs figés sans tomber dans les pièges du calque littéral}. En version comme en thème, une locution adverbiale se comporte comme une \cle{unité de sens indivisible} : on ne la traduit jamais mot à mot. Cette section vous donne les correspondances officielles, signale les asymétries entre les deux langues, dresse la liste des erreurs fatales repérées chaque année au Bac, et vous propose une \cle{méthode} pas à pas pour les neutraliser.

\courssub{Correspondances régulières : le socle incontournable}

Les locutions ci-dessous constituent le \cle{noyau dur} que tout élève de Terminale A doit maîtriser pour l'épreuve de traduction. Elles apparaissent dans les sujets de version (texte journalistique, littéraire ou scientifique) et sont exigées en thème (rédaction d'essai, lettre formelle, article d'opinion).

\begin{center}
\begin{tabularx}{\linewidth}{|X|X|X|}
\hline
\rowcolor{popblueL}
\textbf{Français} & \textbf{Espagnol} & \textbf{Classe / Nuance} \\
\hline
souvent & \esp{a menudo} & fréquence élevée, registre courant \\
\hline
fréquemment & \esp{con frecuencia} & plus formel, écrit \\
\hline
soudain / tout à coup / brusquement & \esp{de repente} / \esp{de pronto} & \esp{de repente} : plus courant ; \esp{de pronto} : légèrement plus littéraire \\
\hline
peu à peu / progressivement & \esp{poco a poco} & processus graduel, très usité \\
\hline
de temps en temps / par moments & \esp{de vez en cuando} & fréquence faible, irrégulière \\
\hline
aussitôt / immédiatement / sur-le-champ & \esp{en seguida} / \esp{al instante} & \esp{en seguida} : oral et écrit ; \esp{al instante} : plus formel \\
\hline
exprès / intentionnellement & \esp{a propósito} / \esp{aposta} & \esp{a propósito} : « à dessein » ; \esp{aposta} : familier (Péninsule) \\
\hline
en général / globalement & \esp{por lo general} / \esp{en general} & \esp{por lo general} : locution figée ; \esp{en general} : plus libre \\
\hline
à long terme & \esp{a largo plazo} & \textbf{jamais} \esp{a largo término} \\
\hline
à court terme & \esp{a corto plazo} & symétrique du précédent \\
\hline
en vain & \esp{en vano} & locution unique, pas *\esp{en vanidad} \\
\hline
par hasard / par chance & \esp{por casualidad} / \esp{de casualidad} & \esp{por casualidad} : standard \\
\hline
\end{tabularx}
\end{center}

\begin{exemplebox}[Exemples authentiques en contexte]
\esp{Los jóvenes cameruneses usan a menudo el bendskin para desplazarse en Douala.} \\
« Les jeunes Camerounais utilisent souvent le bendskin pour se déplacer à Douala. »

\esp{De repente, el cielo se oscureció y empezó a llover torrencialmente.} \\
« Soudain, le ciel s'est assombri et il a commencé à pleuvoir torrentiellement. »

\esp{La deforestación avanza poco a poco en la cuenca del Congo.} \\
« La déforestation progresse peu à peu dans le bassin du Congo. »

\esp{De vez en cuando, los elefantes cruzan la carretera cerca de Bertoua.} \\
« De temps en temps, les éléphants traversent la route près de Bertoua. »

\esp{Ante la alarma, los bomberos llegaron en seguida.} \\
« Face à l'alarme, les pompiers sont arrivés aussitôt. »

\esp{Lo hice a propósito para que lo vieras.} \\
« Je l'ai fait exprès pour que tu le voies. »

\esp{Por lo general, la temporada de cacao comienza en octubre.} \\
« En général, la saison du cacao commence en octobre. »

\esp{Los efectos del cambio climático se verán a largo plazo.} \\
« Les effets du changement climatique se verront à long terme. »

\esp{Buscamos agua en vano durante tres horas.} \\
« Nous avons cherché de l'eau en vain pendant trois heures. »
\end{exemplebox}

\courssub{Asymétries : quand le français utilise un adverbe simple et l'espagnol une locution (et inversement)}

C'est une source majeure d'erreurs : l'apprenant francophone cherche un adverbe simple en espagnol là où la norme exige une locution, ou l'inverse. Le tableau suivant récapitule les cas les plus fréquents au programme de Terminale.

\begin{center}
\begin{tabularx}{\linewidth}{|X|X|X|}
\hline
\rowcolor{poporangeL}
\textbf{Français (adverbe simple)} & \textbf{Espagnol (locution)} & \textbf{Remarque} \\
\hline
bientôt & \esp{dentro de poco} / \esp{pronto} & \esp{pronto} est adverbe, mais \esp{dentro de poco} très courant \\
\hline
déjà & \esp{ya} (adverbe) mais \esp{desde ya} = « dès maintenant » & attention : \esp{ya} seul = « déjà / maintenant » \\
\hline
encore & \esp{todavía} / \esp{aún} (adverbes) mais \esp{de nuevo} = « à nouveau » & \esp{otra vez} = « une fois de plus » \\
\hline
ensemble & \esp{juntos} (adverbe) mais \esp{al mismo tiempo} = « en même temps » & \\
\hline
vraiment & \esp{realmente} / \esp{de veras} (locution) & \esp{de veras} plus expressif, oral \\
\hline
malheureusement & \esp{desgraciadamente} / \esp{por desgracia} (locution) & \esp{por desgracia} très usité en presse \\
\hline
heureusement & \esp{afortunadamente} / \esp{por suerte} (locution) & \esp{por suerte} = « par chance » \\
\hline
\rowcolor{popgreenL}
\textbf{Espagnol (adverbe simple)} & \textbf{Français (locution)} & \textbf{Remarque} \\
\hline
\esp{aprisa} & à la hâte / rapidement & littéraire \\
\hline
\esp{despacio} & lentement / tout doucement & courant \\
\hline
\esp{anteayer} & avant-hier & mot unique ! \\
\hline
\esp{antier} & avant-hier (Amérique) & variant régional \\
\hline
\end{tabularx}
\end{center}

\begin{propriete}[Règle d'or : traduire le bloc, jamais le mot]
Une locution adverbiale (espagnole ou française) est une \cle{unité lexicale figée} dont le sens n'est pas la somme de ses composants. En traduction, on cherche \emph{l'\cle{équivalent fonctionnel}} dans la langue cible, c'est-à-dire la locution (ou l'adverbe) qui remplit exactement le même rôle sémantique, syntaxique et stylistique.
\end{propriete}

\courssub{Faux amis, calques et pièges classiques : la liste noire du Bac}

Voici les erreurs qui reviennent chaque session et qui font perdre \emph{au minimum un point} sur l'exercice de version (sur 4 ou 5 points). Mémorisez-les comme des \emph{interdits absolus}.

\begin{attention}[Erreur fatale n°1 : « de vez en cuando » ≠ « de fois en quand »]
\textbf{Jamais} n'écrivez : *\esp{de veces en cuando}, *\esp{de vez en cuando} (orthographe correcte mais calque syntaxique du français « de fois en quand » qui n'existe pas), *\esp{de tiempo en tiempo}. \\
\emph{Correct} : \esp{de vez en cuando} (invariable, figé). \\
\emph{Variante soutenue} : \esp{de cuando en cuando}.
\end{attention}

\begin{attention}[Erreur fatale n°2 : « à long terme » ≠ *\esp{a largo término}]
\textbf{Término} signifie « terme » au sens de \emph{mot, expression} (ex. : \esp{término técnico}) ou « fin, limite » (ex. : \esp{el término del plazo}). Pour « délai, échéance, horizon temporel », on utilise exclusivement \esp{plazo}. \\
\emph{Correct} : \esp{a largo plazo} / \esp{a corto plazo} / \esp{a mediano plazo}.
\end{attention}

\begin{propriete}[Piège « en seguida » vs « ensuite » : distinction cruciale]
\esp{En seguida} signifie \textbf{immédiatement, tout de suite, sans délai} (valeur temporelle de contiguïté). \\
Pour exprimer « ensuite, par la suite » (séquence, non-contiguïté), on utilise \esp{después}, \esp{luego}, \esp{posteriormente}, \esp{más tarde}.
\end{propriete}

\begin{exemplebox}[Contraste \esp{en seguida} / \esp{después} / \esp{luego}]
\esp{En seguida vengo.} \quad « Je reviens \textbf{tout de suite} / \textbf{aussitôt}. » (action immédiate) \\
\esp{Después fuimos al río.} \quad « \textbf{Ensuite} / \textbf{Par la suite}, nous sommes allés à la rivière. » (étape suivante) \\
\esp{Terminé el examen, luego salimos.} \quad « Une fois l'examen fini, \textbf{ensuite} nous sommes sortis. »
\end{exemplebox}

\begin{attention}[Erreur fatale n°3 : *\esp{sín embargo} (accent mal placé) ou confusion avec « sans embargo »]
\begin{itemize}
\item L'adverbe de concession s'écrit \textbf{sans accent} : \esp{sin embargo} (deux mots, accent tonique sur \emph{em} → \emph{sín embargo} phonétiquement, mais \textbf{pas d'accent graphique}).
\item Ne le confondez jamais avec l'expression \esp{sin embargo} = « cependant, néanmoins » et la locution \esp{sin embargo de que} = « malgré le fait que » (+ \cle{subjonctif}).
\item \emph{Calque interdit} : *\esp{sans embargo} (anglicisme / francisation).
\end{itemize}
\end{attention}

\begin{exemplebox}[\esp{sin embargo} en contexte argumentatif (environnement)]
\esp{Las empresas prometen reforestar; sin embargo, la tala ilegal continúa.} \\
« Les entreprises promettent de reboiser ; \textbf{cependant}, le braconnage forestier continue. »

\esp{Sin embargo de que la ley \textbf{prohíba} la caza, los furtivos actúan con impunidad.} \\
« \textbf{Bien que / Malgré le fait que} la loi interdise la chasse, les braconniers agissent en toute impunité. » (Après \esp{sin embargo de que} → \cle{subjonctif} : \esp{prohíba} ; \esp{actúan} reste à l'\cle{indicatif} car fait réel).
\end{exemplebox}

\courssub{Tableau récapitulatif des pièges : traduction mot à mot (barrée) vs traduction validée}

\begin{popfigure}
\begin{tikzpicture}[
  node distance=0.8cm and 1.2cm,
  false/.style={rectangle, draw=popred, thick, fill=popred!10, rounded corners, text width=5.5cm, align=center, font=\small},
  true/.style={rectangle, draw=popgreen, thick, fill=popgreen!10, rounded corners, text width=5.5cm, align=center, font=\small},
  arrow/.style={-Stealth, thick, poporange},
  title/.style={font=\bfseries\large, text=popdark}
]
\node[title] (t1) at (0,5.5) {Traduction mot à mot (FAUSSE)};
\node[title] (t2) at (12,5.5) {Traduction correcte (VALIDÉE)};

% Row 1
\node[false] (f1) at (0,4.2) {*de vez en cuando → *de fois en quand};
\node[true] (v1) at (12,4.2) {de vez en cuando → de temps en temps};
\draw[arrow] (f1.east) -- (v1.west) node[midway, above, font=\tiny, poporange] {✗ calque};

% Row 2
\node[false] (f2) at (0,3.0) {*a largo término → à long terme};
\node[true] (v2) at (12,3.0) {a largo plazo → à long terme};
\draw[arrow] (f2.east) -- (v2.west) node[midway, above, font=\tiny, poporange] {✗ faux ami};

% Row 3
\node[false] (f3) at (0,1.8) {*en seguida → ensuite};
\node[true] (v3) at (12,1.8) {en seguida → tout de suite / aussitôt};
\draw[arrow] (f3.east) -- (v3.west) node[midway, above, font=\tiny, poporange] {✗ contresens};

% Row 4
\node[false] (f4) at (0,0.6) {*sin embargo → sans embargo};
\node[true] (v4) at (12,0.6) {sin embargo → cependant / néanmoins};
\draw[arrow] (f4.east) -- (v4.west) node[midway, above, font=\tiny, poporange] {✗ anglicisme};

% Row 5
\node[false] (f5) at (0,-0.6) {*de tiempo en tiempo → de temps en temps};
\node[true] (v5) at (12,-0.6) {de vez en cuando → de temps en temps};
\draw[arrow] (f5.east) -- (v5.west) node[midway, above, font=\tiny, poporange] {✗ invention};

% Row 6
\node[false] (f6) at (0,-1.8) {*todo a coup → tout à coup};
\node[true] (v6) at (12,-1.8) {de repente / de pronto → tout à coup};
\draw[arrow] (f6.east) -- (v6.west) node[midway, above, font=\tiny, poporange] {✗ calque phonétique};
\end{tikzpicture}
\legende{Les six pièges mortels en traduction de locutions adverbiales : à gauche, les productions erronées observées en copie d'élèves ; à droite, les formes attendues par le jury.}
\end{popfigure}

\courssub{Connecteurs logiques et locutions argumentatives pour l'essai sur l'environnement}

Au-delà des locutions adverbiales pures, l'essai argumentatif (voir Section 4) exige une maîtrise des \emph{connecteurs logiques} qui sont souvent des locutions figées. Voici les indispensables pour structurer un essai sur le \esp{medioambiente} :

\begin{center}
\begin{tabularx}{\linewidth}{|X|X|X|}
\hline
\rowcolor{poppurpleL}
\textbf{Fonction} & \textbf{Locution espagnole} & \textbf{Français} \\
\hline
Introduire la thèse & \esp{en primer lugar}, \esp{para empezar}, \esp{ante todo} & en premier lieu, pour commencer, avant tout \\
\hline
Ajouter un argument & \esp{además}, \esp{por añadidura}, \esp{asimismo}, \esp{del mismo modo} & de plus, en outre, de même \\
\hline
Opposer / Concéder & \esp{sin embargo}, \esp{no obstante}, \esp{por el contrario}, \esp{aunque} & cependant, néanmoins, au contraire, bien que \\
\hline
Cause / Conséquence & \esp{por consiguiente}, \esp{por tanto}, \esp{como resultado}, \esp{de ahí que} + subj. & par conséquent, donc, de sorte que \\
\hline
Exemplifier & \esp{por ejemplo}, \esp{así}, \esp{en particular}, \esp{un caso ilustrativo es} & par exemple, ainsi, en particulier, un cas illustratif est \\
\hline
Conclure & \esp{en conclusión}, \esp{por último}, \esp{en definitiva}, \esp{en resumen} & en conclusion, enfin, en définitive, en résumé \\
\hline
\end{tabularx}
\end{center}

\begin{methode}[Stratégie de traduction d'une phrase contenant des locutions (Version)]
\begin{enumerate}
\item \textbf{Repérage} : Surlignez \emph{toutes} les locutions adverbiales et connecteurs de la phrase source (espagnol).
\item \textbf{Blocage} : Pour chacune, écrivez au-dessus son équivalent français \emph{figé} (tableau de correspondances), sans analyser les mots isolés.
\item \textbf{Squelette} : Traduisez le reste de la phrase (sujet, verbe, compléments) en respectant la syntaxe française.
\item \textbf{Assemblage} : Insérez les locutions françaises aux bons endroits (place normale de l'adverbe en français : souvent en début de phrase ou après le verbe).
\item \textbf{Relecture} : Vérifiez la cohérence temporelle, logique et le registre de langue.
\end{enumerate}
\end{methode}

\courssub{Entraînement guidé : Version (3 phrases espagnoles → français)}

Appliquez la méthode ci-dessus. Les locutions cibles sont en \textbf{gras} dans l'énoncé.

\begin{exoresolu}[Version guidée : locutions adverbiales]
\textbf{Énoncé :} Traduisez en français les trois phrases suivantes.
\begin{enumerate}
\item \esp{Los agricultores de la región \textbf{de vez en cuando} queman los rastrojos, \textbf{sin embargo} \textbf{por lo general} saben que eso degrada el suelo.}
\item \esp{\textbf{De repente}, el nivel del río subió \textbf{en seguida} tras las lluvias torrenciales, y \textbf{poco a poco} inundó los campos de maíz.}
\item \esp{El gobierno prometió actuar \textbf{a largo plazo}, \textbf{por consiguiente} lanzó un plan de reforestación \textbf{a propósito} para la COP28.}
\end{enumerate}

\tcblower

\textbf{Solution détaillée :}

\textbf{Phrase 1 :}
\begin{itemize}
\item Locutions repérées : \esp{de vez en cuando} = « de temps en temps » ; \esp{sin embargo} = « cependant / néanmoins » ; \esp{por lo general} = « en général / généralement ».
\item Squelette : « Les agriculteurs de la région ... brûlent les chaumes, ... savent que cela dégrade le sol. »
\item Assemblage : « Les agriculteurs de la région \textbf{de temps en temps} brûlent les chaumes, \textbf{cependant} \textbf{généralement} savent que cela dégrade le sol. »
\item \textbf{Version finale} : « Les agriculteurs de la région brûlent \textbf{de temps en temps} les chaumes, \textbf{cependant}, \textbf{en général}, ils savent que cela dégrade le sol. » (Ajustement ordre : \esp{por lo general} devant le verbe → « en général, ils savent »).
\end{itemize}

\textbf{Phrase 2 :}
\begin{itemize}
\item Locutions : \esp{de repente} = « soudain / tout à coup » ; \esp{en seguida} = « aussitôt / immédiatement » ; \esp{poco a poco} = « peu à peu / progressivement ».
\item Squelette : « ... le niveau du fleuve monta ... après les pluies torrentielles, et ... inonda les champs de maïs. »
\item Assemblage : « \textbf{Soudain}, le niveau du fleuve monta \textbf{aussitôt} après les pluies torrentielles, et \textbf{peu à peu} inonda les champs de maïs. »
\item \textbf{Version finale} : « \textbf{Soudain}, le niveau du fleuve a monté \textbf{aussitôt} après les pluies torrentielles, et \textbf{peu à peu} a envahi les champs de maïs. » (\cle{Passé composé} naturel en français pour actions achevées).
\end{itemize}

\textbf{Phrase 3 :}
\begin{itemize}
\item Locutions : \esp{a largo plazo} = « à long terme » ; \esp{por consiguiente} = « par conséquent / c'est pourquoi » ; \esp{a propósito} = « exprès / à dessein / pour l'occasion ».
\item Squelette : « Le gouvernement a promis d'agir ..., ... a lancé un plan de reboisement ... pour la COP28. »
\item Assemblage : « Le gouvernement a promis d'agir \textbf{à long terme}, \textbf{par conséquent} a lancé un plan de reboisement \textbf{exprès} pour la COP28. »
\item \textbf{Version finale} : « Le gouvernement a promis d'agir \textbf{à long terme} ; \textbf{c'est pourquoi} il a lancé un plan de reboisement \textbf{spécialement} pour la COP28. » (\esp{a propósito} ici = « à dessein, spécialement pour » ; \esp{por consiguiente} → « c'est pourquoi » plus fluide que « par conséquent » en milieu de phrase).
\end{itemize}
\end{exoresolu}

\courssub{Entraînement guidé : Thème (3 phrases françaises → espagnol avec locutions imposées)}

\begin{exercice}[Thème : locutions adverbiales imposées]
Traduisez en espagnol en \textbf{utilisant obligatoirement la locution indiquée entre parenthèses}.
\begin{enumerate}
\item Les autorités camerounaises contrôlent \textbf{de temps en temps} les exploitations forestières illégales. (\esp{de vez en cuando})
\item \textbf{Soudain}, une tempête de sable a recouvert la ville de Maroua \textbf{tout de suite} après le coucher du soleil. (\esp{de repente} / \esp{en seguida})
\item Nous devons agir \textbf{à long terme} pour protéger le bassin du Congo ; \textbf{par conséquent}, le ministère a lancé \textbf{exprès} un programme d'éducation environnementale. (\esp{a largo plazo} / \esp{por consiguiente} / \esp{a propósito})
\end{enumerate}
\end{exercice}

\begin{corrige}[Exercice Thème n°3]
\begin{enumerate}
\item \esp{Las autoridades camerunesas controlan \textbf{de vez en cuando} las explotaciones forestales ilegales.}
\item \esp{\textbf{De repente}, una tormenta de arena cubrió la ciudad de Maroua \textbf{en seguida} tras la puesta de sol.}
\item \esp{Debemos actuar \textbf{a largo plazo} para proteger la cuenca del Congo; \textbf{por consiguiente}, el ministerio ha lanzado \textbf{a propósito} un programa de educación ambiental.}
\end{enumerate}
\textbf{Notes :}
\begin{itemize}
\item Phrase 1 : \esp{controlan} (\cle{présent indicatif}) ou \esp{controlaron} (\cle{prétérit}) selon contexte ; \esp{de vez en cuando} placé avant le verbe (standard).
\item Phrase 2 : \esp{cubrió} (\cle{prétérit}, action ponctuelle achevée) ; \esp{en seguida} après le verbe ; \esp{tras} + infinitif = « après le coucher » (registre écrit).
\item Phrase 3 : \esp{Debemos} (devoir moral, présent) ; \esp{a largo plazo} (jamais \esp{término}) ; \esp{por consiguiente} suivi de point-virgule ou point ; \esp{a propósito} = « exprès, à dessein » (ici « spécialement pour »).
\end{itemize}
\end{corrige}

\courssub{Le savais-tu ? : La statistique du jury}

\begin{savaistu}[Au Bac, une locution mal traduite = un point en moins]
Chaque année, la commission de correction du Baccalauréat (session juin et session septembre) signale dans son \emph{rapport jury} que les items « locutions adverbiales » et « connecteurs logiques » sont \textbf{les plus discriminants} en version. Sur 5 points attribués à l'exercice de version, \textbf{1 point entier} est souvent réservé à la traduction correcte de 2 à 3 locutions ciblées dans le texte. Les copies qui donnent *\esp{de tiempo en tiempo}, *\esp{a largo término} ou *\esp{sans embargo} perdent ce point \emph{systématiquement}, même si le reste de la phrase est parfait. \textbf{Moralité :} apprenez le tableau des correspondances par cœur ; c'est un investissement à rendement garanti.
\end{savaistu}

\courssub{Texte support d'application : « Alerta en el bosque de Dja »}

Le texte suivant, original, intègre \emph{toutes} les locutions étudiées dans cette section (et quelques-unes de la section 2). Il sert de base à l'exercice de version complet de la Section 6. Lisez-le, repérez les locutions (surlignez-les), puis consultez le glossaire.

\begin{textebox}[Alerta en el bosque de Dja — Reportage simulé, Yaoundé, 2024]
En el sureste de Camerún, la reserva de la biosfera del Dja, patrimonio de la humanidad, enfrenta una presión creciente. \textbf{De vez en cuando}, los guardabosques descubren campamentos furtivos donde se ahuma carne de gorila y pangolín. \textbf{De repente}, el sonido de una motosierra rompe el silencio del dosel: los taladores ilegales entran \textbf{en seguida} al amanecer y huyen \textbf{en seguida} antes del mediodía. \textbf{Poco a poco}, los corredores ecológicos se fragmentan y los elefantes pierden sus rutas migratorias ancestrales.

Las ONG locales denuncian que \textbf{por lo general} las multas son irrisorias y \textbf{a largo plazo} la impunidad destruye el capital natural. \textbf{Sin embargo}, un nuevo proyecto comunitario ofrece esperanza: los pueblos bakas y badjoués participan \textbf{a propósito} en la vigilancia con drones y sensores acústicos. \textbf{Por consiguiente}, las alertas tempranas permiten interceptar a los cazadores \textbf{antes de que} dañen la fauna. \textbf{En conclusión}, la conservación eficaz exige justicia rápida, educación ambiental y alternativas económicas para las poblaciones ribereñas.
\end{textebox}

\begin{definition}[Glossaire du texte (mots difficiles)]
\begin{itemize}
\item \esp{la reserva de la biosfera} : la réserve de biosphère
\item \esp{patrimonio de la humanidad} : patrimoine de l'humanité (UNESCO)
\item \esp{enfrenta una presión creciente} : fait face à une pression croissante
\item \esp{guardabosques} : garde-forestier, éco-garde
\item \esp{furtivos} : braconniers (adj. ou n.m.)
\item \esp{ahumar} : fumer (viande), faire du boucan
\item \esp{el dosel} : la canopée (toit de la forêt)
\item \esp{taladores} : abatteurs, bûcherons (illégaux ici)
\item \esp{corredores ecológicos} : couloirs écologiques
\item \esp{irrisorias} : dérisoires (peu élevées)
\item \esp{la impunidad} : l'impunité
\item \esp{capital natural} : capital naturel
\item \esp{pueblos bakas y badjoués} : peuples Baka et Badjoué (autochtones)
\item \esp{sensores acústicos} : capteurs acoustiques
\item \esp{alertas tempranas} : alertes précoces
\item \esp{interceptar} : intercepter
\item \esp{poblaciones ribereñas} : populations riveraines (vivant au bord de la rivière/forêt)
\end{itemize}
\end{definition}

\begin{exemplebox}[Questions de compréhension / repérage (préparation Section 6)]
\begin{enumerate}
\item Combien de locutions adverbiales de \textbf{fréquence} trouvez-vous ? Lesquelles ?
\item Identifiez les deux locutions de \textbf{temps ponctuel} (soudain / aussitôt).
\item Quelle locution exprime la \textbf{concession} ? Quelle en est la traduction exacte ?
\item Repérez la locution de \textbf{conséquence} et celle de \textbf{finalité} (\emph{exprès / à dessein}).
\item Traduisez la dernière phrase (\esp{En conclusión...}) en appliquant la méthode du bloc.
\end{enumerate}
\end{exemplebox}

\courssub{Synthèse visuelle : carte mentale des correspondances et pièges}

\begin{popfigure}
\begin{tikzpicture}[
  mindmap,
  concept color=popblue,
  level 1 concept/.append style={sibling angle=90, level distance=4.5cm},
  level 2 concept/.append style={sibling angle=45, level distance=3.2cm},
  level 3 concept/.append style={sibling angle=30, level distance=2.5cm},
  every node/.style={concept, font=\small},
  grow cyclic
]
\node[concept, text width=2.8cm] {Correspondances et Pièges}
  child[concept color=popgreen] {node[concept, text width=2.5cm] {Régulières}
    child {node[concept, text width=2.2cm] {souvent = a menudo}}
    child {node[concept, text width=2.2cm] {soudain = de repente}}
    child {node[concept, text width=2.2cm] {peu à peu = poco a poco}}
    child {node[concept, text width=2.2cm] {de temps en temps = de vez en cuando}}
    child {node[concept, text width=2.2cm] {aussitôt = en seguida}}
    child {node[concept, text width=2.2cm] {à long terme = a largo plazo}}
  }
  child[concept color=poporange] {node[concept, text width=2.5cm] {Asymétries}
    child {node[concept, text width=2.2cm] {bientôt → dentro de poco}}
    child {node[concept, text width=2.2cm] {malheureusement → por desgracia}}
    child {node[concept, text width=2.2cm] {anteayer → avant-hier}}
  }
  child[concept color=popred] {node[concept, text width=2.5cm] {Pièges mortels}
    child {node[concept, text width=2.2cm] {*de tiempo en tiempo}}
    child {node[concept, text width=2.2cm] {*a largo término}}
    child {node[concept, text width=2.2cm] {en seguida ≠ ensuite}}
    child {node[concept, text width=2.2cm] {*sin embargo (accent)}}
    child {node[concept, text width=2.2cm] {*sans embargo}}
    child {node[concept, text width=2.2cm] {*todo a coup}}
  }
  child[concept color=poppurple] {node[concept, text width=2.5cm] {Connecteurs essai}
    child {node[concept, text width=2.2cm] {sin embargo / no obstante}}
    child {node[concept, text width=2.2cm] {por consiguiente}}
    child {node[concept, text width=2.2cm] {por ejemplo}}
    child {node[concept, text width=2.2cm] {en conclusión}}
  };
\end{tikzpicture}
\legende{Carte mentale récapitulative : correspondances régulières (vert), asymétries adverbe/locution (orange), pièges interdits (rouge), connecteurs pour l'essai argumentatif (pourpre).}
\end{popfigure}

Cette section vous a doté de la \emph{boîte à outils} indispensable pour aborder sereinement les exercices de version et de thème du Bac, ainsi que la rédaction de l'essai argumentatif (Section 4). Mémorisez le tableau des pièges, entraînez-vous sur le texte « Alerta en el bosque de Dja » et passez à la structure de l'essai.

\coursec{L'essai argumentatif : structure et connecteurs}

Dans la section précédente, nous avons appris à repérer les pièges de traduction des locutions adverbiales (\esp{de repente}, \esp{a menudo}, \esp{poco a poco}...) et à choisir la bonne correspondance française. Mais une bonne traduction ne suffit pas : au baccalauréat, on vous demandera aussi de \cle{produire} un texte cohérent, structuré et argumenté. C'est précisément l'objet de cette section : apprendre à construire un \esp{ensayo argumentativo}, notamment sur le thème du \esp{medio ambiente}, en utilisant les connecteurs logiques — et en y intégrant, avec élégance, les locutions adverbiales étudiées plus haut.

\courssub{Qu'est-ce qu'un texte argumentatif ?}

\begin{definition}[El texto argumentativo]
Un \esp{texto argumentativo} (ou \esp{ensayo argumentativo}) est un texte dans lequel on \cle{défend une thèse} — c'est-à-dire une opinion personnelle — sur un \cle{sujet polémique} (qui prête à débat), en s'appuyant sur des \cle{arguments} justifiés par des \cle{exemples}. Son but n'est pas seulement d'informer, mais de \emph{convaincre} le lecteur.
\end{definition}

Observons d'abord un court extrait d'essai pour voir la mécanique en action.

\begin{textebox}[Extracto de un ensayo sobre el medio ambiente]
Hoy en día, la protección del medio ambiente es uno de los grandes debates de nuestra sociedad. ¿Es posible conciliar el desarrollo económico y el respeto de la naturaleza? En mi opinión, sí, pero a condición de cambiar nuestros hábitos.

En primer lugar, la deforestación destruye el hábitat de los animales. Por ejemplo, en la cuenca del Congo, cada año desaparecen miles de hectáreas de bosque. Además, la contaminación del aire provoca enfermedades graves, sobre todo en las grandes ciudades como Douala o Ciudad de México.

Sin embargo, algunos afirman que proteger el medio ambiente cuesta demasiado dinero. Aunque cueste mucho, estoy convencido de que no actuar costará todavía más: las sequías y las inundaciones destruyen cosechas y viviendas.

En conclusión, debemos actuar en seguida. Poco a poco, con pequeños gestos diarios y con leyes más estrictas, podemos construir un futuro más limpio para nuestros hijos.
\end{textebox}

\begin{exemplebox}[Glosario del texto]
\begin{itemize}
\item \esp{conciliar} = concilier
\item \esp{el hábitat} = l'habitat
\item \esp{la cuenca} = le bassin (hydrographique)
\item \esp{la sequía} = la sécheresse
\item \esp{la inundación} = l'inondation
\item \esp{la cosecha} = la récolte
\item \esp{la vivienda} = le logement
\item \esp{el gesto} = le geste
\item \esp{estricto/a} = strict(e)
\item \esp{a condición de} = à condition de (suivi de l'infinitif ou de \esp{que} + subjonctif)
\end{itemize}
\end{exemplebox}

Observez : le texte ne se contente pas d'affirmer. Il \emph{prouve} : chaque idée (argument) est illustrée par un fait concret (exemple), et les paragraphes sont reliés par des mots-charnières (\esp{en primer lugar}, \esp{además}, \esp{sin embargo}, \esp{en conclusión}). C'est toute l'architecture de l'essai.

\begin{definition}[Tesis, argumento, ejemplo]
\begin{itemize}
\item La \cle{tesis} (thèse) est l'\textbf{opinion que l'on défend} : \esp{En mi opinión, debemos proteger los bosques.} « À mon avis, nous devons protéger les forêts. »
\item L'\cle{argumento} (argument) est la \textbf{raison qui justifie} la thèse : \esp{...porque los bosques producen el oxígeno que respiramos.} « ...parce que les forêts produisent l'oxygène que nous respirons. »
\item L'\cle{ejemplo} (exemple) est le \textbf{fait concret qui rend l'argument crédible} : \esp{Por ejemplo, el Amazonas absorbe enormes cantidades de CO2.} « Par exemple, l'Amazonie absorbe d'énormes quantités de CO2. »
\end{itemize}
Sans argument, la thèse n'est qu'une affirmation gratuite ; sans exemple, l'argument reste abstrait. Les trois sont indissociables.
\end{definition}

\courssub{La structure en quatre parties}

Un essai argumentatif au baccalauréat suit un plan rigide qu'il faut respecter scrupuleusement :

\begin{enumerate}
\item \textbf{Introducción} : présentation du sujet (accroche avec \esp{hoy en día}, \esp{actualmente}...) + \cle{problématique} sous forme de question ouverte.
\item \textbf{Desarrollo 1} : annonce de la thèse + premier argument + exemple.
\item \textbf{Desarrollo 2} : deuxième argument + exemple + éventuelle \cle{contre-argumentation} (on présente l'avis adverse avec \esp{sin embargo}, puis on le réfute).
\item \textbf{Conclusión} : bilan de la réflexion + ouverture (une question, un souhait, une perspective).
\end{enumerate}

\begin{popfigure}
\begin{tikzpicture}[
node distance=0.55cm,
etape/.style={rectangle, rounded corners=4pt, draw=popblue, fill=popblueL, align=center, text width=6.2cm, inner sep=5pt, font=\small},
fleche/.style={-{Stealth[length=2.5mm]}, thick, poporange}
]
\node[etape, align=center] (intro){\textbf{Introducción}\\sujet + problemática};
\node[etape, below=of intro, align=center] (tesis){\textbf{Tesis}\\mi opinión};
\node[etape, below=of tesis] (arg1) {\textbf{Argumento 1 + ejemplo}};
\node[etape, below=of arg1] (arg2) {\textbf{Argumento 2 + ejemplo}};
\node[etape, below=of arg2, align=center] (contra){\textbf{Contraargumento}\\sin embargo...};
\node[etape, below=of contra, fill=popgreenL, draw=popgreen, align=center] (concl){\textbf{Conclusión}\\balance + apertura};
\draw[fleche] (intro) -- (tesis);
\draw[fleche] (tesis) -- (arg1);
\draw[fleche] (arg1) -- (arg2);
\draw[fleche] (arg2) -- (contra);
\draw[fleche] (contra) -- (concl);
\end{tikzpicture}
\legende{Organigramme de la structure de l'essai argumentatif}
\end{popfigure}

\courssub{La problématique : poser la bonne question}

La problématique (\esp{la problemática}) est une \cle{question ouverte} qui annonce le débat. Elle ne doit pas appeler une simple réponse par « oui » ou « non » sans discussion. En espagnol, n'oubliez jamais le point d'interrogation ouvrant \esp{¿}.

\begin{exemplebox}[Ejemplos de problemáticas]
\esp{¿Es posible conciliar el desarrollo económico y la protección del medio ambiente?} « Est-il possible de concilier développement économique et protection de l'environnement ? »

\esp{¿Deben los gobiernos prohibir los plásticos de un solo uso?} « Les gouvernements doivent-ils interdire les plastiques à usage unique ? »

\esp{¿Pueden los pequeños gestos individuales salvar el planeta?} « Les petits gestes individuels peuvent-ils sauver la planète ? »
\end{exemplebox}

\courssub{Les connecteurs logiques par fonction}

Les connecteurs (\esp{los conectores lógicos}) sont les articulations du texte. Voici le tableau récapitulatif à apprendre par cœur, avec les correspondances françaises.

\begin{center}
\begin{tabularx}{\linewidth}{|l|X|X|}
\hline
\rowcolor{popblueL} \textbf{Fonction} & \textbf{Español} & \textbf{Français} \\
\hline
Introduire & \esp{en primer lugar} ; \esp{para empezar} ; \esp{hoy en día} ; \esp{actualmente} & premièrement ; pour commencer ; de nos jours ; actuellement \\
\hline
Ajouter & \esp{además} ; \esp{asimismo} ; \esp{también} ; \esp{por otra parte} & de plus ; de même ; aussi ; d'autre part \\
\hline
Opposer & \esp{sin embargo} ; \esp{no obstante} ; \esp{por el contrario} ; \esp{aunque} & cependant ; néanmoins ; au contraire ; bien que \\
\hline
Causer / conclure logiquement & \esp{porque} ; \esp{ya que} ; \esp{puesto que} ; \esp{por lo tanto} ; \esp{por eso} ; \esp{en consecuencia} & parce que ; puisque ; puisque ; par conséquent ; c'est pourquoi ; en conséquence \\
\hline
Conclure & \esp{en conclusión} ; \esp{para terminar} ; \esp{en resumen} ; \esp{en definitiva} & en conclusion ; pour terminer ; en résumé ; en définitive \\
\hline
Opinion & \esp{en mi opinión} ; \esp{a mi juicio} ; \esp{creo que} ; \esp{estoy convencido de que} ; \esp{es evidente que} & à mon avis ; à mon sens ; je crois que ; je suis convaincu que ; il est évident que \\
\hline
\end{tabularx}
\end{center}

\begin{popfigure}
\begin{tikzpicture}[
func/.style={rectangle, rounded corners=4pt, draw=popdark, fill=popgoldL, align=center, font=\small\bfseries, inner sep=4pt},
conn/.style={rectangle, rounded corners=3pt, draw=popblue, fill=white, align=center, font=\footnotesize, inner sep=3pt, text width=3.1cm},
arrow/.style={-{Stealth[length=2.5mm]}, thick, poporange}
]
\node[func] (f1) {Introducir};
\node[conn, below=0.3cm of f1, align=center] (c1){en primer lugar\\para empezar\\hoy en día\\actualmente};
\node[func, right=2.2cm of f1] (f2) {Añadir};
\node[conn, below=0.3cm of f2, align=center] (c2){además\\asimismo\\por otra parte};
\node[func, right=2.2cm of f2] (f3) {Oponer};
\node[conn, below=0.3cm of f3, align=center] (c3){sin embargo\\no obstante\\aunque};
\node[func, right=2.2cm of f3] (f4) {Concluir};
\node[conn, below=0.3cm of f4, align=center] (c4){por lo tanto\\en conclusión\\en definitiva};
\draw[arrow] (c1.east) -- (c2.west);
\draw[arrow] (c2.east) -- (c3.west);
\draw[arrow] (c3.east) -- (c4.west);
\node[font=\footnotesize\itshape, popmuted, below=0.8cm of c2] {le fil du raisonnement : introduction $\rightarrow$ développement $\rightarrow$ conclusion};
\end{tikzpicture}
\legende{Frise des connecteurs classés par fonction}
\end{popfigure}

\begin{exemplebox}[Ejemplos traducidos]
\esp{En primer lugar, la deforestación destruye el hábitat de los animales.} « Premièrement, la déforestation détruit l'habitat des animaux. »

\esp{Además, el transporte contamina el aire de las ciudades.} « De plus, les transports polluent l'air des villes. »

\esp{Sin embargo, muchas empresas no respetan las leyes.} « Cependant, beaucoup d'entreprises ne respectent pas les lois. »

\esp{Por lo tanto, debemos actuar en seguida.} « Par conséquent, nous devons agir tout de suite. »

\esp{En definitiva, proteger la naturaleza es proteger nuestra vida.} « En définitive, protéger la nature, c'est protéger notre vie. »
\end{exemplebox}

\begin{propriete}[Indicatif ou subjonctif après les expressions d'opinion]
Règle essentielle pour le Bac :
\begin{itemize}
\item Après \esp{creo que}, \esp{es evidente que}, \esp{estoy convencido de que} (opinion \textbf{affirmée}) : \textbf{indicatif}.\\
\esp{Creo que el gobierno \textbf{hace} esfuerzos.} « Je crois que le gouvernement fait des efforts. »
\item Après \esp{no creo que}, \esp{dudo que}, \esp{no es seguro que} (opinion \textbf{niée} ou doute) : \textbf{subjonctif}.\\
\esp{No creo que el gobierno \textbf{haga} suficiente.} « Je ne crois pas que le gouvernement en fasse assez. »
\item Après \esp{aunque} : \textbf{indicatif} si le fait est réel/constaté, \textbf{subjonctif} s'il est hypothétique/concessif.\\
\esp{Aunque \textbf{cuesta} mucho, hay que actuar.} (fait réel : ça coûte cher)\\
\esp{Aunque \textbf{cueste} mucho, hay que actuar.} (hypothétique : même si cela devait coûter cher)
\end{itemize}
\end{propriete}

\begin{attention}[Erreurs fréquentes des francophones]
\begin{itemize}
\item Ne commencez \textbf{jamais} tous vos paragraphes par \esp{y} (« et ») : c'est le calque du français scolaire « et... et... ». Variez : \esp{además}, \esp{por otra parte}, \esp{asimismo}.
\item Attention à \esp{en primer lugar} : on ne dit pas \esp{en primer sitio} dans ce sens (calque de « en premier lieu » ; \esp{sitio} = lieu physique).
\item \esp{Por eso} signifie « c'est pourquoi » (conséquence), pas « pour cela » au sens de « pour cette chose » (finalité = \esp{para eso}).
\item \esp{A propósito} signifie « à propos / exprès », jamais « à propos de » au sens de « concernant » : pour « concernant l'environnement », dites \esp{sobre el medio ambiente} ou \esp{acerca del medio ambiente}.
\item N'oubliez pas la virgule après les connecteurs en tête de phrase : \esp{Sin embargo, ...}
\item \esp{Medio ambiente} s'écrit en deux mots (norme RAE), bien que \esp{medioambiente} soit toléré.
\end{itemize}
\end{attention}

\courssub{Intégrer les locutions adverbiales pour enrichir le style}

Un bon essai réutilise les locutions adverbiales de la leçon : elles donnent du rythme et de la précision.

\begin{exemplebox}[Locuciones adverbiales en el ensayo]
\esp{Poco a poco, la conciencia ecológica avanza en nuestras sociedades.} « Petit à petit, la conscience écologique progresse dans nos sociétés. »

\esp{Cada vez más jóvenes participan en la limpieza de las playas.} « De plus en plus de jeunes participent au nettoyage des plages. »

\esp{Sin duda, el reciclaje debe convertirse en un hábito diario.} « Sans doute le recyclage doit-il devenir une habitude quotidienne. »

\esp{De vez en cuando, el gobierno organiza campañas de sensibilización.} « De temps en temps, le gouvernement organise des campagnes de sensibilisation. »

\esp{Debemos actuar en seguida, antes de que sea demasiado tarde.} « Nous devons agir tout de suite, avant qu'il ne soit trop tard. »
\end{exemplebox}

\courssub{Lexique clé : el medio ambiente}

\begin{definition}[Vocabulario esencial para el ensayo]
\begin{center}
\begin{tabularx}{\linewidth}{|X|X|X|}
\hline
\rowcolor{popgreenL} \textbf{Noms / Sustantivos} & \textbf{Verbes / Verbos} & \textbf{Adjectifs / Adjetivos} \\
\hline
\esp{el medio ambiente} (environnement) & \esp{contaminar} (polluer) & \esp{contaminado/a} (pollué) \\
\esp{la contaminación} (pollution) & \esp{destruir} (détruire) & \esp{sostenible} (durable) \\
\esp{la deforestación} (déforestation) & \esp{proteger} (protéger) & \esp{ecológico/a} (écologique) \\
\esp{el cambio climático} (changement climatique) & \esp{reciclar} (recycler) & \esp{renovable} (renouvelable) \\
\esp{la capa de ozono} (couche d'ozone) & \esp{concienciar} (sensibiliser) & \esp{grave} (grave) \\
\esp{el efecto invernadero} (effet de serre) & \esp{agotar} (épuiser) & \esp{urgente} (urgent) \\
\esp{los residuos / la basura} (déchets) & \esp{reducir} (réduire) & \esp{responsable} (responsable) \\
\esp{la energía renovable} (énergie renouvelable) & \esp{cuidar} (soigner/préserver) & \esp{consciente} (conscient) \\
\hline
\end{tabularx}
\end{center}
\end{definition}

\courssub{Méthode : le plan en cinq étapes}

\begin{methode}[Cómo redactar un ensayo argumentativo]
\begin{enumerate}
\item \textbf{Analyser le sujet} : souligner les mots-clés, identifier le thème (ici : \esp{medio ambiente}) et ce qu'on vous demande (opinion, débat, solution).
\item \textbf{Formuler la problématique} : transformer le sujet en question ouverte commençant par \esp{¿Es posible...?}, \esp{¿Deben...?}, \esp{¿Pueden...?}.
\item \textbf{Choisir deux ou trois arguments} qui défendent votre thèse (et éventuellement un contre-argument à réfuter).
\item \textbf{Trouver un exemple concret par argument} : un pays, une ville (Douala, Yaoundé, Malabo), un fait d'actualité, une expérience personnelle.
\item \textbf{Rédiger avec les connecteurs} : un connecteur différent en tête de chaque paragraphe, les expressions d'opinion pour la thèse, et deux ou trois locutions adverbiales pour le style. Relire : accents, accords, subjonctif après \esp{no creo que}.
\end{enumerate}
\end{methode}

\begin{savaistu}[Cameroun et Guinée équatoriale : poumons verts d'Afrique]
Le Cameroun et la Guinée équatoriale partagent le bassin du Congo (\esp{la cuenca del Congo}), deuxième poumon vert de la planète après l'Amazonie. La déforestation pour le cacao, le bois d'œuvre ou l'huile de palme menace la biodiversité (gorilles, éléphants de forêt). À Malabo comme à Yaoundé, des ONG (\esp{ONG}) mènent des campagnes \esp{de reforestación} et d'éducation environnementale. Connaître ces réalités locales donne de la force à vos exemples dans l'essai.
\end{savaistu}

\begin{experience}[Débat oral : « Plásticos de un solo uso »]
\begin{enumerate}
\item En binôme : l'un défend l'interdiction totale (\esp{prohibición total}), l'autre propose le recyclage et l'éducation (\esp{reciclaje y educación}).
\item Utilisez obligatoirement : \esp{en primer lugar}, \esp{sin embargo}, \esp{por lo tanto}, \esp{poco a poco}, \esp{en conclusión}.
\item Changez de rôle après 2 minutes. Le professeur note l'usage des connecteurs et du subjonctif (\esp{aunque}, \esp{no creo que}).
\end{enumerate}
\end{experience}

\begin{exoresolu}[Relier les phrases avec le connecteur approprié]
Énoncé. Complétez avec le connecteur qui convient : \esp{sin embargo}, \esp{por lo tanto}, \esp{en primer lugar}, \esp{además}, \esp{en conclusión}.

1. \esp{............, hay que reducir el uso del coche en la ciudad.}
2. \esp{El transporte público es barato; ............, contamina menos que el coche.}
3. \esp{Muchos ciudadanos conocen el problema; ............, pocos cambian sus hábitos.}
4. \esp{El aire está muy contaminado; ............, debemos actuar en seguida.}
5. \esp{............, proteger el medio ambiente es tarea de todos.}
\tcblower
Solution détaillée.

1. \esp{\textbf{En primer lugar}, hay que reducir el uso del coche en la ciudad.} « Premièrement, il faut réduire l'usage de la voiture en ville. » (ouverture du développement)

2. \esp{El transporte público es barato; \textbf{además}, contamina menos que el coche.} « Les transports publics sont bon marché ; de plus, ils polluent moins que la voiture. » (addition d'un argument)

3. \esp{Muchos ciudadanos conocen el problema; \textbf{sin embargo}, pocos cambian sus hábitos.} « Beaucoup de citoyens connaissent le problème ; cependant, peu changent leurs habitudes. » (opposition entre connaissance et action)

4. \esp{El aire está muy contaminado; \textbf{por lo tanto}, debemos actuar en seguida.} « L'air est très pollué ; par conséquent, nous devons agir tout de suite. » (conséquence)

5. \esp{\textbf{En conclusión}, proteger el medio ambiente es tarea de todos.} « En conclusion, protéger l'environnement est l'affaire de tous. » (clôture de l'essai)
\end{exoresolu}

\courssub{Entraînement guidé : thème, version et rédaction}

\begin{exercice}[Thème : locutions adverbiales et connecteurs]
Traduisez en espagnol en utilisant les locutions et connecteurs vus en cours.
\begin{enumerate}
\item De nos jours, la pollution est un problème grave. (hoy en día / grave)
\item Premièrement, nous devons recycler. (en primer lugar / debemos reciclar)
\item De plus, les usines contaminent les rivières. (además / las fábricas / los ríos)
\item Cependant, le gouvernement n'agit pas assez. (sin embargo / el gobierno / no actúa bastante)
\item Par conséquent, les citoyens doivent se mobiliser. (por lo tanto / los ciudadanos / movilizarse)
\item Petit à petit, la situation s'améliore. (poco a poco / la situación / mejorar)
\item En conclusion, protéger la nature est urgent. (en conclusión / proteger la naturaleza / urgente)
\end{enumerate}
\end{exercice}

\begin{corrige}[Exercice Thème]
\begin{enumerate}
\item \esp{Hoy en día, la contaminación es un problema grave.}
\item \esp{En primer lugar, debemos reciclar.}
\item \esp{Además, las fábricas contaminan los ríos.}
\item \esp{Sin embargo, el gobierno no actúa bastante.}
\item \esp{Por lo tanto, los ciudadanos deben movilizarse.}
\item \esp{Poco a poco, la situación mejora.} (ou \esp{va mejorando})
\item \esp{En conclusión, proteger la naturaleza es urgente.}
\end{enumerate}
\end{corrige}

\begin{exercice}[Version : compréhension et traduction]
Traduisez en français le paragraphe suivant :
\esp{«Actualmente, el cambio climático es una realidad innegable. En primer lugar, las temperaturas suben cada año. Además, los fenómenos meteorológicos extremos son más frecuentes. Sin embargo, algunos gobiernos niegan la evidencia. Por lo tanto, la sociedad civil debe presionar para que se tomen medidas urgentes. En definitiva, no hay planeta B.»}
\end{exercice}

\begin{corrige}[Exercice Version]
« Actuellement, le changement climatique est une réalité indéniable. Premièrement, les températures augmentent chaque année. De plus, les phénomènes météorologiques extrêmes sont plus fréquents. Cependant, certains gouvernements nient l'évidence. Par conséquent, la société civile doit faire pression pour que des mesures urgentes soient prises. En définitive, il n'y a pas de planète B. »
\end{corrige}

\begin{exercice}[Rédaction : Essai argumentatif complet]
\textbf{Sujet :} \esp{« ¿Deben las ciudades africanas prohibir los coches antiguos para reducir la contaminación? »}

Rédigez un essai de 150--180 mots environ. Respectez la structure en 4 parties. Utilisez au minimum 5 connecteurs différents et 3 locutions adverbiales. Soignez le subjonctif après \esp{no creo que} / \esp{es necesario que}.
\end{exercice}

\begin{corrige}[Exercice Rédaction -- Proposition de corrigé type]
\begin{textebox}[Ensayo modelo]
Actualmente, la calidad del aire en las grandes ciudades africanas como Douala o Yaoundé es preocupante. ¿Deben las ciudades africanas prohibir los coches antiguos para reducir la contaminación? En mi opinión, sí, pero con medidas de acompañamiento.

En primer lugar, los vehículos viejos emiten gases tóxicos que enferman a la población. Por ejemplo, en Douala, los casos de asma han aumentado mucho en los últimos años. Además, la contaminación acelera el cambio climático, que ya afecta a nuestras cosechas de cacao.

Sin embargo, algunos argumentan que prohibir estos coches perjudica a los más pobres, que no pueden comprarse uno nuevo. Aunque sea cierto, no creo que la solución sea no hacer nada; por el contrario, el Estado debe ofrecer transporte público barato y eficiente.

En conclusión, la prohibición progresiva es necesaria. Poco a poco, con autobuses modernos y carriles bici, podemos respirar un aire más limpio. En definitiva, la salud no tiene precio.
\end{textebox}
\end{corrige}

\begin{aretenir}[L'essentiel de la section]
Un essai argumentatif = une \cle{thèse} défendue par des \cle{arguments} illustrés d'\cle{exemples}, en quatre parties (introduction avec problématique, deux développements, conclusion). Les \cle{connecteurs} structurent le raisonnement : \esp{en primer lugar} pour introduire, \esp{además} pour ajouter, \esp{sin embargo} pour opposer, \esp{por lo tanto} pour conclure logiquement, \esp{en conclusión} pour terminer. Après \esp{creo que} : indicatif ; après \esp{no creo que} : subjonctif. Enrichissez le style avec des locutions adverbiales : \esp{poco a poco}, \esp{cada vez más}, \esp{sin duda}, \esp{en seguida}, \esp{de vez en cuando}. Vocabulaire clé : \esp{medio ambiente}, \esp{contaminación}, \esp{deforestación}, \esp{cambio climático}, \esp{sostenible}, \esp{reciclar}.
\end{aretenir}

\coursec{Lexique du medioambiente et modèle d'essai}

Dans la section précédente, nous avons appris à bâtir le squelette d'un essai argumentatif : introduction, thèse, arguments, exemples, conclusion, reliés par des connecteurs logiques. Mais un squelette sans chair n'est rien : pour convaincre sur le thème de l'environnement, il faut disposer d'un \cle{vocabulaire} précis et riche. C'est l'objet de cette section : d'abord le lexique, organisé en trois champs (problèmes, solutions, acteurs), puis la lecture commentée d'un essai modèle complet, où vous verrez la théorie prendre vie.

\courssub{Le champ lexical des problèmes environnementaux}

Observons d'abord ces phrases, inspirées de la presse hispanophone et adaptées au contexte camerounais. Remarquez au passage les locutions adverbiales, héritées de la section précédente :

\esp{La deforestación avanza poco a poco en la región del Sur.} « La déforestation avance petit à petit dans la région du Sud. »

\esp{De vez en cuando aparecen nuevos vertederos en Douala.} « De temps en temps, de nouvelles décharges apparaissent à Douala. »

\esp{El calentamiento global amenaza las cosechas de cacao.} « Le réchauffement climatique menace les récoltes de cacao. »

\begin{definition}[Les problèmes : vocabulaire essentiel]
\begin{itemize}
\item \esp{la contaminación} : la pollution
\item \esp{la deforestación} : la déforestation
\item \esp{el calentamiento global} : le réchauffement climatique
\item \esp{los residuos plásticos} : les déchets plastiques
\item \esp{la sequía} : la sécheresse
\item \esp{la desertificación} : la désertification
\item \esp{los vertederos} : les décharges (sauvages ou publiques)
\item \esp{el humo de los tubos de escape} : la fumée des pots d'échappement
\end{itemize}
\end{definition}

\begin{attention}[Faux amis et pièges de traduction]
\esp{Contaminar} signifie \emph{polluer} (l'air, l'eau, le sol) ; il ne s'emploie pas au sens médical de « contaminer » (transmettre une maladie), qui se dit \esp{contagiar}. De même, \esp{el medioambiente} s'écrit souvent en un seul mot dans la presse moderne, mais la forme en deux mots \esp{medio ambiente} reste correcte. Attention aussi au genre : \esp{el problema} est masculin malgré sa terminaison en \emph{-a} (comme \esp{el clima}, \esp{el planeta}) ; on dira donc \esp{un problema grave}, jamais \esp{una problema}. Enfin, \esp{la sequía} prend un accent sur le \emph{í} et \esp{las energías} sur le \emph{í} également : ne les oubliez pas à l'écrit.
\end{attention}

\courssub{Le champ lexical des solutions}

\begin{definition}[Les solutions : vocabulaire essentiel]
\begin{itemize}
\item \esp{reciclar} : recycler ; \esp{reutilizar} : réutiliser
\item \esp{proteger el medioambiente} : protéger l'environnement
\item \esp{concienciar a la población} : sensibiliser la population
\item \esp{las energías renovables} : les énergies renouvelables
\item \esp{el desarrollo sostenible} : le développement durable
\item \esp{plantar árboles} : planter des arbres
\item \esp{reducir el consumo} : réduire la consommation
\item \esp{apostar por el transporte público} : miser sur le transport en commun
\end{itemize}
\end{definition}

\begin{propriete}[Verbes d'action : constructions à maîtriser]
Ces verbes se construisent ainsi :
\begin{itemize}
\item \esp{dañar / destruir / amenazar + complément direct} : \esp{Los residuos dañan el suelo.} « Les déchets abîment le sol. »
\item \esp{preservar / proteger + complément} : \esp{Debemos preservar los bosques.} « Nous devons préserver les forêts. »
\item \esp{fomentar / prohibir + nom ou infinitif} : \esp{El gobierno debe prohibir las bolsas de plástico.} « Le gouvernement doit interdire les sacs plastiques. »
\item \esp{concienciar / sensibilizar + a + personne} : \esp{Las ONG sensibilizan a los ciudadanos.} « Les ONG sensibilisent les citoyens. »
\item \esp{apostar por + nom} : \esp{Hay que apostar por las energías limpias.} « Il faut miser sur les énergies propres. »
\end{itemize}
\end{propriete}

\begin{attention}[Attention à la préposition \esp{a}]
En français on dit « sensibiliser la population » sans préposition ; en espagnol, la personne animée exige la préposition \esp{a} : \esp{concienciar \textbf{a} la población}. C'est la fameuse « préposition accusative » (ou \emph{a} personnel) espagnole, absolument obligatoire devant un complément direct de personne. Autre piège : « miser sur » ne se traduit pas mot à mot ; on dit \esp{apostar \textbf{por}}, avec \esp{por} et non \esp{sobre}.
\end{attention}

\courssub{Le champ lexical des acteurs}

Un bon essai ne dit jamais seulement « il faut » : il précise \emph{qui} doit agir. Voici les acteurs à mobiliser :

\begin{center}
\begin{tabularx}{\linewidth}{|X|X|X|}\hline
\rowcolor{popblueL} Español & Français & Exemple de rôle \\ \hline
\esp{el gobierno} & le gouvernement & \esp{aprobar leyes} (voter des lois) \\ \hline
\esp{las ONG} & les ONG & \esp{organizar campañas} (organiser des campagnes) \\ \hline
\esp{los ciudadanos} & les citoyens & \esp{reciclar a diario} (recycler chaque jour) \\ \hline
\esp{las empresas} & les entreprises & \esp{reducir los residuos} (réduire les déchets) \\ \hline
\esp{los jóvenes} & les jeunes & \esp{plantar árboles} (planter des arbres) \\ \hline
\esp{la comunidad internacional} & la communauté internationale & \esp{firmar acuerdos} (signer des accords) \\ \hline
\end{tabularx}
\end{center}

\begin{popfigure}
\begin{tikzpicture}[mindmap, grow cyclic, every node/.style={concept, align=center, font=\small},
root concept/.append style={concept color=popblue, font=\bfseries},
level 1/.append style={level distance=3.4cm, sibling angle=120},
level 2/.append style={level distance=2.4cm, sibling angle=50}]
\node[root concept]{medio\-ambiente}
  child[concept color=poporange]{ node {problemas}
    child { node[concept color=poporangeL, text=popdark] {contaminación} }
    child { node[concept color=poporangeL, text=popdark] {deforestación} }
    child { node[concept color=poporangeL, text=popdark] {sequía y vertederos} }
  }
  child[concept color=popgreen]{ node {soluciones}
    child { node[concept color=popgreenL, text=popdark] {reciclar y reutilizar} }
    child { node[concept color=popgreenL, text=popdark] {energías renovables} }
    child { node[concept color=popgreenL, text=popdark] {plantar árboles} }
  }
  child[concept color=poppurple]{ node {actores}
    child { node[concept color=poppurpleL, text=popdark] {gobierno y ONG} }
    child { node[concept color=poppurpleL, text=popdark] {ciudadanos} }
    child { node[concept color=poppurpleL, text=popdark] {jóvenes y empresas} }
  };
\end{tikzpicture}
\legende{Carte mentale du lexique de l'environnement : trois branches à mémoriser pour tout essai.}
\end{popfigure}

\courssub{Lecture commentée d'un essai modèle}

Voici maintenant un essai complet d'environ 150 mots, du niveau attendu au baccalauréat. Lisez-le d'abord en entier, puis nous l'analyserons phrase par phrase.

\begin{textebox}[¿Debemos prohibir las bolsas de plástico en Camerún?]
Hoy en día, las bolsas de plástico contaminan poco a poco nuestras ciudades. ¿Debemos prohibirlas definitivamente? En mi opinión, la respuesta es sí.

En primer lugar, estas bolsas tardan siglos en desaparecer; por eso, vemos montañas de residuos por todas partes en Douala. Además, dañan gravemente el suelo y los ríos, y amenazan la salud de los ciudadanos.

Sin embargo, algunos comerciantes dependen de ellas para vender sus productos, y muchas familias las reutilizan a menudo en casa.

Por lo tanto, el gobierno debe prohibir las bolsas no biodegradables y, al mismo tiempo, fomentar las bolsas de papel y de tela. Las ONG y los jóvenes pueden concienciar a la población poco a poco.

En conclusión, debemos actuar en seguida: proteger el medioambiente es proteger nuestra vida.
\end{textebox}

\textbf{Glossaire :} \esp{hoy en día} = de nos jours ; \esp{tardar siglos en desaparecer} = mettre des siècles à disparaître ; \esp{por todas partes} = partout ; \esp{la tela} = le tissu ; \esp{al mismo tiempo} = en même temps ; \esp{depender de} = dépendre de.

\textbf{Questions de compréhension :} 1) Quelle est la thèse de l'auteur ? 2) Quels deux arguments la soutiennent ? 3) Quelle objection est évoquée ? 4) Quelle solution concrète est proposée ?

\courssub{Analyse du modèle : la structure rendue visible}

Reprenons l'essai et soulignons ses rouages :

\begin{itemize}
\item \textbf{Introduction} : \esp{Hoy en día...} situe le sujet ; la question \esp{¿Debemos prohibirlas?} pose le problème ; \esp{En mi opinión, la respuesta es sí} annonce la \cle{thèse}.
\item \textbf{Arguments} : \esp{En primer lugar} (durée de décomposition, exemple de \cle{Douala}) ; \esp{Además} (dégâts sur le sol, les rivières, la santé — on peut citer aussi la plage de \cle{Kribi}, souillée par les plastiques).
\item \textbf{Objection} : \esp{Sin embargo} introduit le point de vue adverse (les commerçants), ce qui montre l'honnêteté intellectuelle de l'auteur.
\item \textbf{Contre-proposition} : \esp{Por lo tanto} + solutions nommant les \cle{acteurs} (gobierno, ONG, jóvenes).
\item \textbf{Conclusion} : \esp{En conclusión} + formule-chute chiasmatique facile à imiter : \esp{proteger el medioambiente es proteger nuestra vida}.
\item \textbf{Locutions adverbiales} : remarquez \esp{poco a poco}, \esp{a menudo}, \esp{en seguida}, \esp{por todas partes}, \esp{hoy en día} : elles donnent du naturel et de la précision au texte, et elles rapportent des points au Bac.
\end{itemize}

\begin{exemplebox}[Exemples traduits à réutiliser]
\esp{el calentamiento global} = le réchauffement climatique ; \esp{concienciar a la población} = sensibiliser la population ; \esp{los vertederos} = les décharges ; \esp{las bolsas tardan siglos en desaparecer} = les sacs mettent des siècles à disparaître ; \esp{fomentar las energías renovables} = encourager les énergies renouvelables ; \esp{actuar en seguida} = agir immédiatement.
\end{exemplebox}

\begin{savaistu}[Un exemple local en or pour le Bac]
Le Cameroun a interdit les sacs plastiques non biodégradables en 2014. Citer cette mesure dans votre essai (\esp{En 2014, el gobierno camerunés prohibió las bolsas de plástico no biodegradables}) montre au correcteur que vous reliez la langue à la réalité de votre pays : c'est exactement le type d'\cle{exemple} précis qui fait la différence entre une bonne et une excellente copie. Notez l'adjectif \esp{camerunés} (camerounais), avec son accent.
\end{savaistu}

\begin{exoresolu}[Compléter un paragraphe argumentatif]
Énoncé : complétez ce paragraphe avec les mots suivants : \esp{concienciar}, \esp{residuos}, \esp{por lo tanto}, \esp{renovables}, \esp{poco a poco}.

\esp{Los \dots\ plásticos llenan los vertederos de Yaoundé. \dots\ , las empresas deben reducir el embalaje y apostar por las energías \dots\ . Las escuelas pueden \dots\ a los alumnos, y así cambiaremos \dots\ nuestros hábitos.}
\tcblower
Solution détaillée : \esp{Los \textbf{residuos} plásticos llenan los vertederos de Yaoundé. \textbf{Por lo tanto}, las empresas deben reducir el embalaje y apostar por las energías \textbf{renovables}. Las escuelas pueden \textbf{concienciar} a los alumnos, y así cambiaremos \textbf{poco a poco} nuestros hábitos.}

Traduction : « Les déchets plastiques remplissent les décharges de Yaoundé. Par conséquent, les entreprises doivent réduire l'emballage et miser sur les énergies renouvelables. Les écoles peuvent sensibiliser les élèves, et ainsi nous changerons petit à petit nos habitudes. » Notez la logique : problème (\esp{residuos}) → connecteur de conséquence (\esp{por lo tanto}) → solutions (\esp{renovables}, \esp{concienciar}) → nuance temporelle (\esp{poco a poco}). Vérifiez aussi l'accord : \esp{energías renovables} (féminin pluriel).
\end{exoresolu}

\begin{exercice}[À vous de jouer]
1) Traduisez en espagnol : « De nos jours, la déforestation menace nos forêts ; par conséquent, les citoyens doivent planter des arbres immédiatement. » 2) Rédigez un paragraphe de cinq lignes sur la pollution des plages de Kribi, avec au moins deux locutions adverbiales et un connecteur d'opposition.
\end{exercice}

\begin{corrige}[Exercice]
1) \esp{Hoy en día, la deforestación amenaza nuestros bosques; por lo tanto, los ciudadanos deben plantar árboles en seguida.} 2) Modèle possible : \esp{Las playas de Kribi sufren a menudo la contaminación de los residuos plásticos. Sin embargo, los jóvenes organizan de vez en cuando campañas de limpieza. Poco a poco, la población toma conciencia del problema.} Vérifiez la présence d'au moins deux locutions et d'un connecteur d'opposition (\esp{sin embargo}, \esp{no obstante}).
\end{corrige}

\begin{aretenir}[L'essentiel de la section]
Pour réussir un essai sur l'environnement : mémorisez les trois champs lexicaux (problèmes, solutions, acteurs), construisez les verbes d'action avec les bonnes prépositions (\esp{concienciar a}, \esp{apostar por}, \esp{depender de}), citez un exemple local précis (Douala, Kribi, la loi de 2014) et parsemez votre texte de locutions adverbiales (\esp{poco a poco}, \esp{en seguida}, \esp{hoy en día}) qui prouvent votre maîtrise de la langue.
\end{aretenir}

\coursec{Entraînement guidé : thème, version et rédaction}

Dans la section précédente, nous avons constitué notre lexique du \esp{medioambiente} et analysé un modèle d'essai argumentatif complet, avec ses connecteurs et sa charpente en trois parties. Il est temps maintenant de passer de l'observation à la \cle{pratique} : c'est en traduisant et en rédigeant soi-même que les locutions adverbiales cessent d'être une liste à mémoriser pour devenir des outils vivants. Cette dernière section vous propose un parcours complet, exactement calqué sur les épreuves du baccalauréat : une \cle{version} guidée, un \cle{thème} guidé, une rédaction à trous, puis une correction collective justifiée et un barème pour vous auto-évaluer.

\courssub{Version guidée : un fleuve en danger}

\begin{methode}[Méthode de la version en quatre pas]
\begin{enumerate}
\item \textbf{Lire} tout le paragraphe une première fois sans traduire, pour saisir le sens global (ici : la pollution d'un fleuve).
\item \textbf{Repérer et souligner} les locutions adverbiales : ce sont des blocs indissociables qui se traduisent d'un seul tenant.
\item \textbf{Traduire} phrase par phrase, en commençant par le squelette (sujet-verbe-complément), puis en plaçant la locution à l'endroit naturel du français.
\item \textbf{Relire} la traduction française à voix haute : si une phrase sonne bizarre en français, c'est qu'elle est mal traduite, même si chaque mot est correct.
\end{enumerate}
\end{methode}

\begin{textebox}[El río Wouri, texto para la versión]
El río Wouri sufre \textbf{a causa de} la contaminación industrial. \textbf{De vez en cuando}, los vecinos de Douala ven peces muertos en la superficie. \textbf{Poco a poco}, las autoridades toman medidas, pero \textbf{de repente} una nueva fábrica vierte sus residuos en el agua. \textbf{Sin embargo}, los jóvenes no pierden la esperanza: \textbf{a menudo} organizan jornadas de limpieza en las orillas.
\end{textebox}

\textbf{Glossaire :} \esp{los vecinos} : les habitants, les voisins ; \esp{la superficie} : la surface ; \esp{verter} : déverser (conjugué ici : \esp{vierte}, 3\up{e} personne du singulier, verbe à diphtongaison e $\rightarrow$ ie) ; \esp{los residuos} : les déchets ; \esp{las orillas} : les rives, les berges ; \esp{perder la esperanza} : perdre espoir ; \esp{tomar medidas} : prendre des mesures.

\textbf{Questions de compréhension (avant de traduire) :}
\begin{enumerate}
\item ¿Qué problema sufre el río Wouri ?
\item ¿Qué hacen los jóvenes para ayudar ?
\item ¿Cuántas locuciones adverbiales aparecen destacadas en negrita ? Clasifícalas (tiempo, modo, causa, oposición).
\end{enumerate}

\begin{attention}[Ne jamais laisser une locution non traduite]
En version, l'examinateur vérifie systématiquement que chaque locution a été rendue en français. Si le sens exact vous échappe, \textbf{ne l'omettez jamais} : donnez un équivalent approché cohérent avec le contexte. Par exemple, si vous avez oublié que \esp{de vez en cuando} signifie « de temps en temps », écrivez au moins « parfois » ou « quelquefois » : vous perdez une fraction de point sur la nuance, mais vous sauvez le sens. Une locution sautée, en revanche, est sanctionnée comme une erreur de compréhension globale.
\end{attention}

\begin{exoresolu}[Version corrigée : El río Wouri]
Traduisez en français le paragraphe \esp{El río Wouri} ci-dessus.
\tcblower
\textbf{Traduction de référence :}

« Le fleuve Wouri souffre \textbf{à cause de} la pollution industrielle. \textbf{De temps en temps}, les habitants de Douala voient des poissons morts à la surface. \textbf{Petit à petit}, les autorités prennent des mesures, mais \textbf{soudain} une nouvelle usine déverse ses déchets dans l'eau. \textbf{Cependant}, les jeunes ne perdent pas espoir : ils organisent \textbf{souvent} des journées de nettoyage sur les berges. »

\textbf{Justification de chaque choix :}
\begin{itemize}
\item \esp{a causa de} : locution de cause, équivalent direct de « à cause de » (et non « grâce à », réservé à une cause positive : \esp{gracias a}).
\item \esp{de vez en cuando} : « de temps en temps » ; « parfois » est accepté.
\item \esp{poco a poco} : « petit à petit » ou « peu à peu » ; « progressivement » est également accepté si le contexte s'y prête.
\item \esp{de repente} : « soudain » ou « tout à coup » ; attention, « de repente » n'existe pas en français : la traduction mot à mot est une faute.
\item \esp{sin embargo} : « cependant » ou « pourtant » ; en tête de phrase, il est suivi d'une virgule en français comme en espagnol.
\item \esp{a menudo} : « souvent » ; placé après le verbe en français (« ils organisent souvent »), ce qui est la place la plus naturelle.
\end{itemize}
\end{exoresolu}

\courssub{Thème guidé : quatre phrases à traduire}

\begin{methode}[En thème : le squelette d'abord, la locution ensuite]
Face à une phrase française à traduire en espagnol, ne commencez jamais par la locution. Procédez ainsi :
\begin{enumerate}
\item Construisez le \textbf{squelette} de la phrase : sujet + verbe (conjugué correctement !) + compléments. Exemple : « les jeunes prennent conscience du problème » $\rightarrow$ \esp{los jóvenes toman conciencia del problema}.
\item \textbf{Insérez} la locution imposée à la bonne place : en tête (suivie d'une virgule) ou après le verbe. $\rightarrow$ \esp{Poco a poco, los jóvenes toman conciencia del problema}.
\item \textbf{Vérifiez} : accents, accords, ponctuation espagnole (¿ ¡), virgule après la locution initiale.
\end{enumerate}
\end{methode}

\begin{propriete}[Locution en tête de phrase : la virgule obligatoire]
Une locution adverbiale placée \textbf{en tête de phrase} est suivie d'une \cle{virgule} en espagnol, comme en français :

\esp{De repente, empezó a llover.} « Soudain, il s'est mis à pleuvoir. »

\esp{Sin embargo, el problema continúa.} « Cependant, le problème continue. »

En revanche, les locutions courtes placées \textbf{après le verbe} ne prennent pas de virgule : \esp{Aquí llueve a menudo.} « Ici, il pleut souvent. »
\end{propriete}

\begin{exemplebox}[Modèles traduits à imiter]
\begin{itemize}
\item « Il pleut souvent ici, mais soudain le soleil revient. » = \esp{Aquí llueve a menudo, pero de repente vuelve el sol.}
\item « Petit à petit, les jeunes prennent conscience du problème. » = \esp{Poco a poco, los jóvenes toman conciencia del problema.}
\end{itemize}
Observez : dans le premier exemple, \esp{a menudo} suit le verbe (place naturelle pour « souvent »), tandis que \esp{de repente} précède le verbe pour marquer la rupture. Dans le second, \esp{poco a poco} ouvre la phrase et appelle une virgule.
\end{exemplebox}

\begin{exercice}[Thème : quatre phrases]
Traduisez en espagnol en utilisant obligatoirement les locutions indiquées :
\begin{enumerate}
\item Les élèves de Terminale travaillent souvent ensemble pour protéger l'environnement. (\esp{a menudo})
\item Soudain, la pluie a commencé à tomber sur le marché de Yaoundé. (\esp{de repente})
\item Petit à petit, la forêt du Cameroun retrouve sa beauté grâce aux jeunes volontaires. (\esp{poco a poco})
\item L'air de la ville est pollué ; cependant, les habitants ne perdent pas espoir. (\esp{sin embargo})
\end{enumerate}
\end{exercice}

\begin{corrige}[Exercice de thème]
\begin{enumerate}
\item \esp{Los alumnos de Terminale trabajan a menudo juntos para proteger el medioambiente.} --- Squelette : \esp{los alumnos trabajan juntos} ; \esp{a menudo} s'insère après le verbe. « Ensemble » = \esp{juntos}, accordé au masculin pluriel. « Pour protéger » = \esp{para proteger} : \esp{para} + infinitif exprime le but.
\item \esp{De repente, la lluvia empezó a caer sobre el mercado de Yaoundé.} --- Locution en tête + virgule. « Commencer à » = \esp{empezar a} + infinitif ; prétérit indéfini \esp{empezó} pour une action ponctuelle passée (attention à l'accent).
\item \esp{Poco a poco, el bosque de Camerún recupera su belleza gracias a los jóvenes voluntarios.} --- Attention : \esp{gracias a} (cause positive), pas \esp{a causa de}. « Cameroun » se dit \esp{Camerún} avec accent sur la dernière syllabe.
\item \esp{El aire de la ciudad está contaminado; sin embargo, los habitantes no pierden la esperanza.} --- \esp{estar} + participe pour l'état résultant (\esp{está contaminado}) ; \esp{sin embargo} en tête de proposition, précédé du point-virgule et suivi de la virgule. En espagnol, on ne met pas d'espace avant le point-virgule.
\end{enumerate}
\end{corrige}

\begin{attention}[Pièges fréquents des francophones en thème]
\begin{itemize}
\item Ne déplacez pas la locution sans raison : si le français dit « les élèves travaillent souvent », écrivez \esp{los alumnos trabajan a menudo} et non \esp{a menudo los alumnos trabajan} ; respectez l'ordre de la phrase source.
\item \esp{sin embargo} s'écrit en \textbf{deux mots}, invariable : ne jamais écrire \esp{sinembargo} en un seul mot.
\item N'oubliez pas la \textbf{virgule} après la locution initiale : \esp{De repente empezó} (sans virgule) est une faute de ponctuation sanctionnée.
\item Faux ami : « actuellement » = \esp{actualmente} et « actuel » = \esp{actual} ; en revanche, \esp{de repente} ne signifie jamais « récemment » mais toujours « soudain ».
\end{itemize}
\end{attention}

\courssub{Rédaction guidée : l'essai à trous}

Voici le début d'un essai argumentatif sur l'environnement. Les connecteurs et les locutions ont disparu : à vous de les réintégrer. Banque de mots : \esp{sin embargo} ; \esp{en primer lugar} ; \esp{poco a poco} ; \esp{además} ; \esp{por eso} ; \esp{en conclusión} ; \esp{a menudo}.

\begin{textebox}[Ensayo a completar : Un planeta en peligro]
Hoy en día, el medioambiente es un tema que preocupa a todo el mundo. \textbf{(1) \ldots\ldots\ldots\ldots}, es necesario reconocer que la situación es grave: los ríos están contaminados y el aire de las grandes ciudades es irrespirable.

\textbf{(2) \ldots\ldots\ldots\ldots}, debemos reducir los residuos plásticos que \textbf{(3) \ldots\ldots\ldots\ldots} terminan en el océano. \textbf{(4) \ldots\ldots\ldots\ldots}, es fundamental plantar árboles: \textbf{(5) \ldots\ldots\ldots\ldots}, los bosques pueden recuperarse si actuamos juntos. \textbf{(6) \ldots\ldots\ldots\ldots}, los gobiernos deben crear leyes más estrictas.

\textbf{(7) \ldots\ldots\ldots\ldots}, proteger la naturaleza es proteger nuestra propia vida: cada gesto cuenta, y el futuro del planeta está en nuestras manos.
\end{textebox}

\textbf{Glossaire :} \esp{irrespirable} : irrespirable ; \esp{los residuos plásticos} : les déchets plastiques ; \esp{plantar árboles} : planter des arbres ; \esp{las leyes estrictas} : des lois strictes ; \esp{cada gesto cuenta} : chaque geste compte ; \esp{recuperarse} : se régénérer, se reconstituer.

\begin{corrige}[Essai à trous]
\begin{enumerate}
\item \esp{Sin embargo} : opposition entre « tout le monde s'inquiète » et « la situation est grave » ; connecteur d'opposition en tête + virgule.
\item \esp{En primer lugar} : ouverture de l'énumération des arguments.
\item \esp{a menudo} : locution de fréquence placée avant le verbe de la subordonnée (\esp{que a menudo terminan}) ; « souvent » se place ici naturellement.
\item \esp{Además} : ajout d'un deuxième argument (« de plus »).
\item \esp{poco a poco} : locution de manière exprimant la progression graduelle de la régénération des forêts.
\item \esp{Por eso} : conséquence logique (« c'est pourquoi ») reliant la nécessité d'agir à l'action des gouvernements.
\item \esp{En conclusión} : connecteur canonique d'ouverture de la conclusion d'un essai argumentatif.
\end{enumerate}
\textbf{Justification méthodologique :} chaque connecteur remplit une fonction précise (opposition, ordre, addition, conséquence, conclusion). Au baccalauréat, un essai sans connecteurs visibles perd des points sur le critère « structure », même si les idées sont bonnes.
\end{corrige}

\courssub{Tableau récapitulatif et barème type Bac}

\begin{center}
\begin{tabularx}{\linewidth}{|X|X|X|}
\hline
\rowcolor{popblueL} Español & Français & Fonction / emploi \\
\hline
a menudo & souvent & fréquence ; après le verbe \\
\hline
de repente & soudain, tout à coup & temps/rupture ; tête + virgule \\
\hline
poco a poco & petit à petit, peu à peu & manière progressive ; tête + virgule \\
\hline
de vez en cuando & de temps en temps & fréquence occasionnelle \\
\hline
sin embargo & cependant, pourtant & opposition ; tête + virgule \\
\hline
en seguida & tout de suite, aussitôt & temps ; après le verbe \\
\hline
a propósito & exprès / à propos & intention ou transition \\
\hline
a causa de / gracias a & à cause de / grâce à & cause négative / positive \\
\hline
\end{tabularx}
\end{center}

\begin{aretenir}[Barème type Bac : les quatre critères]
\begin{itemize}
\item \textbf{Exactitude lexicale} (environ 25 \%) : chaque locution traduite par un équivalent correct ; aucune locution omise.
\item \textbf{Grammaire} (environ 25 \%) : conjugaisons, accords, ser/estar, ponctuation espagnole, virgule après locution initiale.
\item \textbf{Structure} (environ 25 \%) : pour l'essai, introduction + développement argumenté (thèse, arguments, exemples) + conclusion.
\item \textbf{Connecteurs logiques} (environ 25 \%) : variété et pertinence (\esp{en primer lugar, además, sin embargo, por eso, en conclusión}).
\end{itemize}
\end{aretenir}

\begin{popfigure}
\begin{tikzpicture}
\node[draw=popblue, fill=popblueL, rounded corners, font=\bfseries, minimum width=4.2cm, minimum height=0.8cm] (t) at (0,0) {Grille d'auto-évaluation};
\node[draw=popdark, fill=popcream, rounded corners, align=left, anchor=north west, text width=11.5cm] at (-5.75,-0.7) {%
\textbf{1. Locutions adverbiales} : toutes traduites, aucune omise \hfill $\square$ oui \ $\square$ non \ $\square$ à améliorer\\[3pt]
\textbf{2. Connecteurs logiques} : au moins quatre, variés et bien placés \hfill $\square$ oui \ $\square$ non \ $\square$ à améliorer\\[3pt]
\textbf{3. Structure de l'essai} : introduction, arguments, conclusion \hfill $\square$ oui \ $\square$ non \ $\square$ à améliorer\\[3pt]
\textbf{4. Lexique du medioambiente} : vocabulaire précis et riche \hfill $\square$ oui \ $\square$ non \ $\square$ à améliorer};
\end{tikzpicture}
\legende{Grille d'auto-évaluation à remplir après chaque thème, version ou essai : cochez honnêtement chaque critère avant de rendre votre copie.}
\end{popfigure}

\begin{experience}[Correction collective en classe]
Par groupes de quatre : chaque élève lit sa version du texte \esp{El río Wouri} à voix haute. Le groupe compare les choix de traduction des cinq locutions et doit \textbf{justifier} chaque équivalent retenu en s'appuyant sur le tableau récapitulatif. Ensuite, un rapporteur présente au tableau la version du groupe, et la classe vote pour la meilleure formulation de chaque locution. Terminez en remplissant individuellement la grille d'auto-évaluation : c'est exactement le regard que portera le correcteur du baccalauréat sur votre copie.
\end{experience}

\begin{savaistu}[Pourquoi tant de locutions avec « a » et « de » ?]
La plupart des locutions adverbiales espagnoles sont nées de la préposition \esp{a} ou \esp{de} suivie d'un nom ou d'un adjectif figé : \esp{a menudo} (littéralement « à fréquent »), \esp{de repente} (« de soudaineté »), \esp{a propósito} (« à propos »). Le français a suivi le même chemin : « à cause de », « de temps en temps », « petit à petit ». Cette parenté de construction est une chance pour vous : les deux langues se ressemblent, à condition de respecter la virgule après la locution initiale et de ne jamais traduire mot à mot (\esp{de repente} ne se traduit jamais par « de repente » !).
\end{savaistu}

\newpage

\coursec{Activité d'intégration}

\courssub{Objectifs de l'activité}

À l'issue de cette activité d'intégration, tu seras capable de :
\begin{itemize}
\item rédiger un \cle{essai argumentatif} court (120 à 150 mots) en espagnol, structuré en introduction, développement et conclusion ;
\item mobiliser correctement au moins quatre \cle{locutions adverbiales} étudiées dans la leçon (\esp{de repente}, \esp{a menudo}, \esp{poco a poco}, \esp{de vez en cuando}, \esp{en seguida}, \esp{a propósito}…) ;
\item utiliser au moins cinq \cle{connecteurs logiques} variés (\esp{en primer lugar}, \esp{además}, \esp{sin embargo}, \esp{por lo tanto}, \esp{en conclusión}…) ;
\item employer le lexique de l'environnement (\esp{el medioambiente}, \esp{la contaminación}, \esp{reciclar}, \esp{los residuos}, \esp{proteger}…) ;
\item évaluer la production d'un camarade à l'aide d'une grille et défendre oralement ta proposition lors d'un débat.
\end{itemize}

\courssub{Situation}

La mairie de ta ville, en partenariat avec le \emph{Foro Juvenil de África Central}, organise un concours d'expression écrite et orale intitulé \esp{« Salvemos nuestro medioambiente »}. Les jeunes de Yaoundé, Douala, Kribi, Bata et Malabo sont invités à proposer des solutions concrètes pour protéger l'environnement dans leur ville. Le texte gagnant sera publié dans le bulletin bilingue du forum et son auteur défendra ses idées lors de la grande finale à Douala.

Voici l'appel à participation, tel qu'il a été affiché dans ton lycée :

\begin{textebox}[Appel à participation — Foro Juvenil de África Central]
\textbf{CONVOCATORIA : « SALVEMOS NUESTRO MEDIOAMBIENTE »}\par
\smallskip
Jóvenes de África Central :\par
\smallskip
Nuestras ciudades se ensucian poco a poco: los residuos de plástico invaden a menudo las calles de Douala, las playas de Kribi se contaminan de vez en cuando y el lago Chad pierde agua día tras día. Sin embargo, ¡todavía podemos actuar!\par
\smallskip
Escribe un ensayo argumentativo (120-150 palabras) con el título:\par
\smallskip
\textbf{« ¿Cómo podemos proteger el medioambiente en nuestra ciudad? »}\par
\smallskip
Tu texto debe contener una introducción con la problemática, dos argumentos con ejemplos locales y una conclusión con una propuesta concreta. Usa al menos cuatro locuciones adverbiales y cinco conectores lógicos.\par
\smallskip
Los mejores ensayos se leerán en voz alta y la clase votará la propuesta más convincente. ¡Participa en seguida!
\end{textebox}

\emph{Traduction :} « Jeunes d'Afrique centrale : nos villes se salissent petit à petit : les déchets plastiques envahissent souvent les rues de Douala, les plages de Kribi se polluent de temps en temps et le lac Tchad perd de l'eau jour après jour. Cependant, nous pouvons encore agir ! Écris un essai argumentatif (120-150 mots) intitulé : "Comment pouvons-nous protéger l'environnement dans notre ville ?" Ton texte doit contenir une introduction avec la problématique, deux arguments avec des exemples locaux et une conclusion avec une proposition concrète. Utilise au moins quatre locutions adverbiales et cinq connecteurs logiques. Les meilleurs essais seront lus à voix haute et la classe votera la proposition la plus convaincante. Participe aussitôt ! »

\emph{Glossaire :} \esp{convocatoria} = appel à candidatures ; \esp{ensuciarse} = se salir ; \esp{los residuos} = les déchets ; \esp{en voz alta} = à voix haute ; « convaincant » se dit \esp{convincente} (attention au faux ami \esp{conveniente} = pratique, avantageux).

\courssub{Consignes}

\begin{enumerate}
\item \textbf{Lecture et repérage.} Lis l'appel à participation. Souligne en rouge les quatre locutions adverbiales qu'il contient et en bleu le connecteur d'opposition. Précise pour chacune sa classe (temps, lieu, manière, quantité) et son équivalent français.
\item \textbf{Préparation du lexique.} Complète mentalement ta fiche de vocabulaire : cinq noms (\esp{la basura}, \esp{el reciclaje}, \esp{la deforestación}, \esp{el vertedero}, \esp{el medioambiente}), trois verbes (\esp{proteger}, \esp{reciclar}, \esp{contaminar}) et deux adjectifs (\esp{contaminado}, \esp{limpio}). Vérifie le genre de chaque nom.
\item \textbf{Plan de l'essai.} Rédige ton plan en trois parties : (a) problématique sous forme de question ; (b) deux arguments, chacun suivi d'un exemple local (les embouteillages et les ordures à Douala, les plages de Kribi, la disparition du lac Tchad, les déchets plastiques à Yaoundé…) ; (c) conclusion avec une proposition concrète.
\item \textbf{Rédaction.} Écris ton essai de 120 à 150 mots. Intègre \textbf{au moins quatre locutions adverbiales} différentes et \textbf{au moins cinq connecteurs logiques} variés. Souligne-les dans ton texte pour faciliter la correction.
\item \textbf{Relecture grammaticale.} Vérifie : les accents (\esp{contaminación}, \esp{además}, \esp{también}), l'accord des adjectifs (\esp{una ciudad limpia}), l'emploi du présent de l'indicatif et du subjonctif après \esp{es importante que} + subjonctif.
\item \textbf{Correction par les pairs.} Échange ta copie avec un voisin. À l'aide de la grille d'évaluation ci-dessous, note son texte sur 10 et écris un conseil d'amélioration en espagnol (\esp{Debes revisar los acentos} ; \esp{Añade un ejemplo local}).
\item \textbf{Débat oral.} Les cinq meilleurs essais sont lus à voix haute. La classe vote la proposition la plus convaincante et justifie son choix en une phrase : \esp{Voto por el texto de… porque sus argumentos son más sólidos}.
\end{enumerate}

\begin{center}
\begin{tabularx}{\linewidth}{|X|c|}
\hline
\rowcolor{popblueL} \textbf{Critère d'évaluation} & \textbf{Points} \\
\hline
Structure respectée (introduction, 2 arguments, conclusion) & 2 \\
\hline
Au moins 4 locutions adverbiales correctement employées & 2 \\
\hline
Au moins 5 connecteurs logiques variés et bien placés & 2 \\
\hline
Lexique de l'environnement riche et exact & 2 \\
\hline
Correction grammaticale (accents, accords, verbes) & 2 \\
\hline
\end{tabularx}
\end{center}

\courssub{Ressources à mobiliser}

\begin{aretenir}[Rappel : les locutions adverbiales par classe]
\begin{itemize}
\item \emph{Temps / fréquence} : \esp{a menudo} (souvent), \esp{de vez en cuando} (de temps en temps), \esp{en seguida} (aussitôt), \esp{de repente} (soudain), \esp{día tras día} (jour après jour).
\item \emph{Manière} : \esp{poco a poco} (petit à petit), \esp{en voz alta} (à voix haute), \esp{a propósito} (exprès ; à propos).
\item \emph{Lieu} : \esp{arriba} (en haut), \esp{en todas partes} (partout), \esp{cerca de aquí} (près d'ici).
\item \emph{Quantité} : \esp{por lo menos} (au moins), \esp{apenas} (à peine), \esp{en absoluto} (pas du tout).
\end{itemize}
\end{aretenir}

\begin{aretenir}[Boîte à outils de l'essai argumentatif]
\textbf{Structure :}
\begin{itemize}
\item \emph{Introducción} : présenter le sujet + problématique (\esp{¿Cómo podemos…?}).
\item \emph{Desarrollo} : deux arguments, chacun illustré d'un exemple local.
\item \emph{Conclusión} : bilan + proposition concrète.
\end{itemize}
\textbf{Connecteurs :} \esp{en primer lugar} (d'abord), \esp{además} / \esp{también} (de plus), \esp{por ejemplo} (par exemple), \esp{sin embargo} (cependant), \esp{por lo tanto} / \esp{por eso} (c'est pourquoi), \esp{en conclusión} (en conclusion).\par
\smallskip
\textbf{Grammaire utile :} \esp{hay que} + infinitif (il faut) ; \esp{es importante que} + subjonctif ; \esp{debemos} + infinitif ; \esp{si} + présent, futur (\esp{si actuamos, mejoraremos la ciudad}).
\end{aretenir}

\begin{attention}[Pièges à éviter]
\begin{itemize}
\item Après \esp{En primer lugar}, mets une virgule : \esp{En primer lugar, debemos reciclar}.
\item \esp{a menudo} signifie « souvent », jamais « au menu » !
\item « Petit à petit » se dit \esp{poco a poco}, et non \esp{pequeño a pequeño}.
\item « De temps en temps » se dit \esp{de vez en cuando}, et non \esp{de tiempo en tiempo}.
\item Après \esp{es importante que}, le subjonctif est obligatoire : \esp{Es importante que reciclemos}, et non \esp{que reciclamos}.
\item N'oublie pas les accents : \esp{contaminación}, \esp{además}, \esp{también}, \esp{plástico}, \esp{árboles}.
\item \esp{convincente} = convaincant ; \esp{conveniente} = pratique, avantageux (faux ami).
\end{itemize}
\end{attention}

\courssub{Proposition de production et corrigé commenté}

\begin{exoresolu}[Modèle d'essai : « ¿Cómo podemos proteger el medioambiente en nuestra ciudad? »]
Rédige un essai de 120 à 150 mots répondant à l'appel du Foro Juvenil, avec au moins quatre locutions adverbiales et cinq connecteurs logiques.
\tcblower
\textbf{Production modèle (environ 145 mots) :}\par
\smallskip
\esp{Nuestras ciudades de África Central se ensucian poco a poco: la basura de plástico invade a menudo las calles de Douala y las playas de Kribi. ¿Cómo podemos proteger el medioambiente en nuestra ciudad?}\par
\smallskip
\esp{En primer lugar, debemos reciclar los residuos. Por ejemplo, en Yaoundé, los jóvenes pueden organizar cada sábado una jornada de limpieza en los mercados. Además, hay que plantar árboles: el lago Chad pierde agua día tras día y la deforestación avanza sin parar. Sin embargo, estas acciones no bastan si los ciudadanos no cambian de mentalidad.}\par
\smallskip
\esp{Por lo tanto, es importante que el ayuntamiento eduque a la población de vez en cuando con campañas en las escuelas. En conclusión, si actuamos en seguida y todos juntos, nuestra ciudad será poco a poco más limpia y nuestra vida será mejor.}\par
\smallskip
\textbf{Commentaire des choix de langue :}
\begin{itemize}
\item \textbf{Locutions adverbiales (5)} : \esp{poco a poco} (manière, « petit à petit »), \esp{a menudo} (fréquence, « souvent »), \esp{día tras día} (temps, « jour après jour »), \esp{de vez en cuando} (fréquence, « de temps en temps »), \esp{en seguida} (temps, « aussitôt »). Elles sont placées naturellement, avant ou après le verbe.
\item \textbf{Connecteurs (6)} : \esp{En primer lugar} ouvre le premier argument ; \esp{Por ejemplo} introduit l'exemple local ; \esp{Además} ajoute le deuxième argument ; \esp{Sin embargo} marque la nuance ; \esp{Por lo tanto} tire la conséquence ; \esp{En conclusión} annonce la fin.
\item \textbf{Grammaire} : \esp{debemos reciclar} et \esp{hay que plantar} expriment la nécessité ; \esp{es importante que el ayuntamiento eduque} emploie correctement le subjonctif présent (\esp{eduque}) ; \esp{si actuamos…, nuestra ciudad será} utilise la structure si + présent, futur.
\item \textbf{Lexique} : \esp{la basura de plástico}, \esp{reciclar los residuos}, \esp{la deforestación}, \esp{jornada de limpieza} : champ lexical de l'environnement riche et précis.
\item \textbf{Ancrage local} : Douala, Kribi, Yaoundé et le lac Tchad rendent l'argumentation concrète et crédible, comme le demande la convocatoria.
\end{itemize}
\end{exoresolu}

\courssub{Bilan de l'activité}

Cette activité t'a permis de réinvestir l'ensemble des savoir-faire de la leçon dans une tâche de communication complète et réaliste :
\begin{itemize}
\item \textbf{Compréhension de l'écrit} : repérage des locutions adverbiales et des connecteurs dans un document authentique ;
\item \textbf{Lexique} : mobilisation du champ lexical de l'environnement dans un contexte camerounais et centrafricain ;
\item \textbf{Grammaire et traduction} : emploi correct des locutions adverbiales, des structures d'obligation et du subjonctif ;
\item \textbf{Expression écrite} : production d'un essai argumentatif structuré, conforme aux exigences de l'examen ;
\item \textbf{Expression orale et coopération} : correction par les pairs, lecture à voix haute et débat argumenté.
\end{itemize}

\begin{experience}[Pour aller plus loin : le mini-débat]
En groupes de quatre, choisissez la meilleure proposition du forum et transformez-la en discours d'une minute. Un élève la présente, les autres posent des questions (\esp{¿Cuánto costará tu propuesta?} ; \esp{¿Quién la va a aplicar?}). Le présentateur doit répondre en utilisant au moins deux locutions adverbiales et un connecteur d'opposition. Cette simulation prépare à l'épreuve orale du baccalauréat.
\end{experience}

\newpage

\coursec{Méthodes et exercices résolus}

\coursec{Leçon E10 : Traducción — Locuciones adverbiales ; Redacción — Ensayo argumentativo sobre el medioambiente}

\courssub{Texte support et découverte}

\begin{textebox}[Article de presse simulé — Cameroun : urgence plastique]
\esp{Yaoundé, 12 mars 2024. — Les récentes pluies diluviennes ont mis à nu l'insuffisance criante de la gestion des déchets plastiques dans la capitale. Les caniveaux du quartier Mvog-Ada, comme ceux de New-Bell à Douala, se sont bouchés \textbf{de golpe}, provoquant des inondations soudaines qui ont endommagé des centaines de foyers. \textbf{De hecho}, selon un rapport du MINEPDED, \textbf{por lo general}, 60\% des déchets solides ne sont pas collectés. \textbf{Poco a poco}, la conscience citoyenne s'éveille : des associations comme « Jeunesse Verte Cameroun » organisent \textbf{a menudo} des opérations « Ville Propre ». \textbf{En seguida}, les autorités annoncent l'acquisition de nouvelles compacteuses, mais les experts insistent : \textbf{a propósito}, la technique ne suffit pas sans changement de mentalité. \textbf{De vez en cuando}, des amendes sont infligées aux contrevenants, mais l'impunité reste \textbf{a menudo} la règle.}
\end{textebox}

\textbf{Traduction / Glossaire :}
« Yaoundé, 12 mars 2024. — Les récentes pluies diluviennes ont mis à nu l'insuffisance criante de la gestion des déchets plastiques dans la capitale. Les caniveaux du quartier Mvog-Ada, comme ceux de New-Bell à Douala, se sont bouchés \textbf{soudain / d'un coup}, provoquant des inondations soudaines qui ont endommagé des centaines de foyers. \textbf{En fait}, selon un rapport du MINEPDED, \textbf{généralement}, 60\% des déchets solides ne sont pas collectés. \textbf{Peu à peu}, la conscience citoyenne s'éveille : des associations comme « Jeunesse Verte Cameroun » organisent \textbf{souvent} des opérations « Ville Propre ». \textbf{Tout de suite / Immédiatement}, les autorités annoncent l'acquisition de nouvelles compacteuses, mais les experts insistent : \textbf{à propos / exprès}, la technique ne suffit pas sans changement de mentalité. \textbf{De temps en temps}, des amendes sont infligées aux contrevenants, mais l'impunité reste \textbf{souvent} la règle. »

\begin{aretenir}[Lexique clé du texte (Medioambiente)]
\begin{itemize}
    \item \textbf{Noms :} \esp{el residuo (plástico/sólido), el desecho, la basura, el vertedero, el caniveal / la cuneta, la inundación, la sequía, la deforestación, la tala ilegal, la caza furtiva, el cambio climático, la capa de ozono, la huella de carbono, la sostenibilidad, la economía circular, el reciclaje, la concienciación, la impunidad, la normativa, la multa, el contenedor, la compactadora}.
    \item \textbf{Verbes :} \esp{contaminar, degradar, agotar, reciclar, reutilizar, reducir, separar, clasificar, colectar, gestionar, legislar, sancionar, concienciar, inundar, obstruir / atascar, talar, cazar}.
    \item \textbf{Adjectifs :} \esp{contaminado, sostenible, renovable, ecológico, biodegradable, tóxico, nocivo, irreversible, urgente, colectivo, individual, institucional}.
    \item \textbf{Locutions du texte :} \esp{de golpe, de hecho, por lo general, poco a poco, a menudo, en seguida, a propósito, de vez en cuando}.
\end{itemize}
\end{aretenir}

\courssub{Les locutions adverbiales : classement et sens}

\begin{methode}[Identifier et traduire une locution adverbiale]
\textbf{Objectif :} Repérer un groupe de mots fonctionnant comme un adverbe (temps, lieu, manière, quantité) et en donner l'équivalent français précis, sans calque littéral.
\begin{enumerate}
    \item \textbf{Repérer le noyau :} Chercher une préposition (\esp{a, de, en, con, sin, por, para}) suivie d'un nom (souvent sans article : \esp{de pronto}, \esp{a veces}) ou d'un adjectif (\esp{de nuevo}, \esp{a solas}).
    \item \textbf{Identifier la valeur sémantique :} Temps (\esp{¿Cuándo?}), Lieu (\esp{¿Dónde?}), Manière (\esp{¿Cómo?}), Quantité (\esp{¿Cuánto?}), Affirmation/Négation/Doute.
    \item \textbf{Consulter la correspondance figée :} La plupart des locutions sont des « blocs » intraduisibles mot à mot. Ex. : \esp{de repente} $\neq$ « de repente » mais \textbf{soudain}.
    \item \textbf{Vérifier le registre :} \esp{A menudo} (soutenu) vs \esp{muchas veces} (courant) ; \esp{de vez en cuando} (ponctuel/irrégulier) vs \esp{a veces} (général).
    \item \textbf{Replacer dans la phrase cible :} Adapter l'ordre des mots (en français, l'adverbe se place souvent après le verbe conjugué ou au début de phrase).
\end{enumerate}
\end{methode}

\begin{propriete}[Classification des locutions adverbiales usuelles (Niveau Terminale)]
\begin{center}
\begin{tabularx}{\linewidth}{|X|X|X|}
\hline
\rowcolor{popblueL} \textbf{TIEMPO (Temps)} & \textbf{LUGAR (Lieu)} & \textbf{MODO (Manière)} \\ \hline
\esp{ahora mismo} (tout de suite) & \esp{a pie} (à pied) & \esp{a mano} (à la main) \\ \hline
\esp{a menudo} (souvent) & \esp{a caballo} (à cheval) & \esp{a máquina} (à la machine) \\ \hline
\esp{a veces} (parfois) & \esp{a la derecha/izquierda} (à dr./g.) & \esp{a propósito} (exprès) \\ \hline
\esp{de golpe} (soudain, brutalement) & \esp{alrededor} (autour) & \esp{a toda prisa} (en vitesse) \\ \hline
\esp{de pronto} (soudain) & \esp{aquí / allí} (ici / là) & \esp{con cuidado} (soigneusement) \\ \hline
\esp{de repente} (soudain) & \esp{cerca / lejos} (près / loin) & \esp{con prisa} (en hâte) \\ \hline
\esp{de vez en cuando} (de temps en temps) & \esp{debajo / encima} (dessous / dessus) & \esp{de buena/mala gana} (de bon/mauvais gré) \\ \hline
\esp{en seguida} (aussitôt, tout de suite) & \esp{delante / detrás} (devant / derrière) & \esp{de memoria} (par cœur) \\ \hline
\esp{en un abrir y cerrar de ojos} (en un clin d'œil) & \esp{en casa} (à la maison) & \esp{de nuevo} (de nouveau) \\ \hline
\esp{hace poco / tiempo} (il y a peu / longtemps) & \esp{en todas partes} (partout) & \esp{despacio} (lentement) \\ \hline
\esp{por lo general} (généralement) & \esp{fuera / dentro} (dehors / dedans) & \esp{paso a paso} (étape par étape) \\ \hline
\esp{poco a poco} (peu à peu) & \esp{por todas partes} (partout) & \esp{por escrito} (par écrit) \\ \hline
\esp{rara vez} (rarement) &  & \esp{por casualidad} (par hasard) \\ \hline
\esp{siempre / nunca} (toujours / jamais) &  & \esp{sin querer} (sans le vouloir) \\ \hline
\esp{tarde / temprano} (tard / tôt) &  & \esp{a solas} (seul, en tête-à-tête) \\ \hline
\end{tabularx}
\end{center}
\begin{center}
\begin{tabularx}{\linewidth}{|X|X|}
\hline
\rowcolor{poporangeL} \textbf{CANTIDAD (Quantité)} & \textbf{AFIRMACIÓN / NEGACIÓN / DUDA} \\ \hline
\esp{a lo sumo} (au maximum) & \esp{desde luego} (bien sûr / certainement) \\ \hline
\esp{a lo menos} (au minimum) & \esp{en absoluto} (absolument pas) \\ \hline
\esp{como máximo / mínimo} (au max / min) & \esp{efectivamente} (effectivement) \\ \hline
\esp{más o menos} (à peu près) & \esp{por supuesto} (bien sûr) \\ \hline
\esp{cuanto más / menos} (plus / moins) & \esp{de ninguna manera} (de aucune façon) \\ \hline
\esp{del todo} (complètement) & \esp{tal vez / quizás} (peut-être $\rightarrow$ Subj.) \\ \hline
\end{tabularx}
\end{center}
\end{propriete}

\begin{exoresolu}[Classification et emploi en contexte]
\textbf{Énoncé :}
\begin{enumerate}
    \item Classez les locutions suivantes dans le tableau ci-dessous : \esp{a pie, de repente, a menudo, a propósito, en seguida, poco a poco, de vez en cuando, a mano, de memoria, por lo general}.
    \item Complétez les phrases avec une locution adaptée (forme correcte).
    \begin{itemize}
        \item[a.] \esp{Mi hermano escribe los poemas \_\_\_\_\_\_\_\_\_\_, no usa computadora.}
        \item[b.] \esp{\_\_\_\_\_\_\_\_\_\_ llueve mucho en la región del Centro en agosto.}
        \item[c.] \esp{Lleguemos \_\_\_\_\_\_\_\_\_\_ al examen, por favor.}
        \item[d.] \esp{El profesor explicó la lección \_\_\_\_\_\_\_\_\_\_ para que todos entendieran.}
    \end{itemize}
\end{enumerate}
\tcblower
\textbf{Solution détaillée :}

\textbf{1. Tableau de classification}

\begin{center}
\begin{tabularx}{\linewidth}{|X|X|X|}
\hline
\rowcolor{popblueL} \textbf{TIEMPO (Temps)} & \textbf{LUGAR (Lieu)} & \textbf{MODO (Manière)} \\ \hline
\esp{de repente} (soudain) & \esp{a pie} (à pied) & \esp{a propósito} (exprès) \\ \hline
\esp{a menudo} (souvent) &  & \esp{a mano} (à la main) \\ \hline
\esp{en seguida} (tout de suite) &  & \esp{de memoria} (par cœur) \\ \hline
\esp{poco a poco} (peu à peu) &  &  \\ \hline
\esp{de vez en cuando} (parfois) &  &  \\ \hline
\esp{por lo general} (généralement) &  &  \\ \hline
\end{tabularx}
\end{center}
\emph{Note :} Aucune locution de « Quantité » ou « Affirmation/Négation » dans cette liste. \esp{A pie} peut être vu comme Manière (moyen de transport) mais classiquement enseigné comme Lieu/Déplacement en Terminale.

\textbf{2. Complétion guidée}
\begin{itemize}
    \item \textbf{[a.]}\esp{Mi hermano escribe los poemas \textbf{a mano}, no usa computadora.}
    \textit{Justification :} Opposition à « usa computadora » (machine) $\rightarrow$ Manière/Instrument : \esp{a mano} (à la main). \esp{A máquina} existerait mais est rare.
    \item \textbf{[b.]}\esp{\textbf{Por lo general} llueve mucho en la región del Centro en agosto.} (Ou \esp{A menudo}).
    \textit{Justification :} Fréquence élevée, généralité climatique $\rightarrow$ Temps (Fréquence) : \esp{Por lo general} / \esp{A menudo}. \esp{De vez en cuando} contredirait « mucho ».
    \item \textbf{[c.]}\esp{Lleguemos \textbf{a tiempo} / \textbf{pronto} al examen, por favor.}
    \textit{Justification :} Consigne de ponctualité. \esp{A tiempo} (à l'heure, locution standard). \esp{En seguida} (tout de suite) possible si l'ordre est immédiat. \emph{Attention :} \esp{Lleguemos} = Subjonctif présent (exhortation 1e pl.).
    \item \textbf{[d.]}\esp{El profesor explicó la lección \textbf{paso a paso} / \textbf{detalladamente} / \textbf{con claridad} para que todos entendieran.}
    \textit{Justification :} Manière méthodique. \esp{Paso a paso} (pas à pas / étape par étape) est une locution très courante. \esp{De prisa} (vite) contredirait le but « para que entendieran ».
\end{itemize}
\end{exoresolu}

\courssub{Correspondances français-espagnol et pièges de traduction}

\begin{exoresolu}[Traduction de phrases à locutions adverbiales (Version/Thème)]
\textbf{Énoncé :} Traduisez en français les phrases 1 à 3, et en espagnol les phrases 4 à 6.
\begin{enumerate}
    \item \esp{Los estudiantes llegaron \textbf{de golpe} a la facultad de Yaoundé.}
    \item \esp{Mi abuela cocina \textbf{a fuego lento} los fines de semana.}
    \item \esp{\textbf{De vez en cuando}, vemos bendskins en el centro de Douala.}
    \item Soudain, la pluie a cessé pendant le match de football.
    \item Peu à peu, les jeunes apprennent à trier les déchets plastiques.
    \item Il a agi exprès pour protéger la forêt tropicale.
\end{enumerate}
\tcblower
\textbf{Solution détaillée :}

\textbf{1. Version (Espagnol $\rightarrow$ Français)}
\begin{itemize}
    \item \textbf{Phrase 1 :} \esp{Los estudiantes llegaron \textbf{de golpe} a la facultad de Yaoundé.}
    \begin{itemize}
        \item \textit{Analyse :} \esp{de golpe} = locution adverbiale de manière (brutalité/suddenness). \esp{llegaron} = prétérit indéfini 3e pers. pl. \esp{a la facultad de Yaoundé} = CC Lieu.
        \item \textit{Traduction :} « Les étudiants sont arrivés \textbf{soudain / d'un coup} à la faculté de Yaoundé. » (Éviter « de coup »).
    \end{itemize}
    \item \textbf{Phrase 2 :} \esp{Mi abuela cocina \textbf{a fuego lento} los fines de semana.}
    \begin{itemize}
        \item \textit{Analyse :} \esp{a fuego lento} = locution de manière (cuisine), litt. « à feu lent ». \esp{cocina} = présent indicatif. \esp{los fines de semana} = locution de temps (régulier).
        \item \textit{Traduction :} « Ma grand-mère cuisine \textbf{à feu doux / à petit feu} les week-ends. »
    \end{itemize}
    \item \textbf{Phrase 3 :} \esp{\textbf{De vez en cuando}, vemos bendskins en el centro de Douala.}
    \begin{itemize}
        \item \textit{Analyse :} \esp{de vez en cuando} = locution de temps (fréquence faible/irrégulière). \esp{vemos} = présent indicatif 1e pl. \esp{bendskins} = mot camerounais (taxi-moto), invariable.
        \item \textit{Traduction :} « \textbf{De temps en temps / Parfois}, on voit des bendskins au centre de Douala. »
    \end{itemize}
\end{itemize}

\textbf{2. Thème (Français $\rightarrow$ Espagnol)}
\begin{itemize}
    \item \textbf{Phrase 4 :} « Soudain, la pluie a cessé pendant le match de football. »
    \begin{itemize}
        \item \textit{Choix :} \esp{De repente} (soudain, imprévu) ou \esp{De golpe}. \esp{De repente} est plus neutre/narratif.
        \item \textit{Temps :} Action ponctuelle, passée $\rightarrow$ \textbf{Prétérit indéfini} : \esp{cesó} (verbe \esp{cesar}, régulier).
        \item \textit{Proposition :} \esp{\textbf{De repente}, la lluvia cesó durante el partido de fútbol.}
    \end{itemize}
    \item \textbf{Phrase 5 :} « Peu à peu, les jeunes apprennent à trier les déchets plastiques. »
    \begin{itemize}
        \item \textit{Choix :} \esp{Poco a poco} (standard) ou \esp{Paulo} (très littéraire, éviter au Bac). \esp{Los jóvenes} (sujet). \esp{Aprenden} (présent indicatif, habitude/généralité). \esp{A separar / a clasificar} (infinitif prépositionnel). \esp{Los residuos / desechos plásticos}.
        \item \textit{Proposition :} \esp{\textbf{Poco a poco}, los jóvenes aprenden a separar los residuos plásticos.}
    \end{itemize}
    \item \textbf{Phrase 6 :} « Il a agi exprès pour protéger la forêt tropicale. »
    \begin{itemize}
        \item \textit{Piège :} « Exprès » $\rightarrow$ \esp{A propósito} (intentionnellement) ou \esp{Adrede} (plus soutenu/littéraire). \esp{A propósito} est le standard attendu.
        \item \textit{Structure :} \esp{Actuó a propósito} + \esp{para} + infinitif (but).
        \item \textit{Proposition :} \esp{Actuó \textbf{a propósito} para proteger la selva tropical.}
    \end{itemize}
\end{itemize}
\textbf{Conclusion :} La maîtrise des locutions figées (\esp{de golpe, a fuego lento, de vez en cuando, de repente, poco a poco, a propósito}) est indispensable pour éviter le mot-à-mot fautif.
\end{exoresolu}

\begin{methode}[Classer les locutions adverbiales par catégorie sémantique]
\textbf{Objectif :} Ranger une liste de locutions dans un tableau (Temps, Lieu, Manière, Quantité, Affirmation/Négation) pour en faciliter la mémorisation et l'emploi correct.
\begin{enumerate}
    \item \textbf{Poser la question test :} Remplacer la locution par \esp{¿Cuándo?} (Temps), \esp{¿Dónde?} (Lieu), \esp{¿Cómo?} (Manière), \esp{¿Cuánto?} (Quantité), \esp{¿Sí/No?} (Aff/Nég).
    \item \textbf{Attention aux "faux amis" de catégorie :} \esp{A propósito} (Manière/Intention) $\neq$ \esp{A propósito de} (Lien/Sujet $\rightarrow$ « À propos de »). \esp{De hecho} (Affirmation/Réalité = « En fait ») $\neq$ \esp{En hecho} (incorrect).
    \item \textbf{Noter la structure interne :} \esp{Prep + Nom (sin art.)} (\esp{de prisa, a pie}), \esp{Prep + Adj} (\esp{de nuevo, a solas}), \esp{Prep + Adv} (\esp{desde ahora, hasta luego}).
    \item \textbf{Mémoriser par paires d'opposés :} \esp{De día / de noche}, \esp{a menudo / rara vez}, \esp{de prisa / despacio}, \esp{por escrito / de palabra}.
    \item \textbf{Créer une phrase-contextue camerounaise :} Ancrer chaque locution dans une réalité locale (ex. \esp{A pie voy al liceo}, \esp{De noche hace calor en Douala}).
\end{enumerate}
\end{methode}

\begin{attention}[Pièges de traduction fréquents (Francophonismes)]
\begin{enumerate}
    \item \textbf{\esp{De repente} vs \esp{De golpe} :} \esp{De repente} = soudain (neutre, narratif). \esp{De golpe} = brutalement, d'un coup (violence, surprise physique). \textit{Ne pas dire « \esp{de repénte} » (accent inexistant).}
    \item \textbf{\esp{A propósito} vs \esp{A propósito de} :} \esp{Lo hice a propósito} (je l'ai fait exprès). \esp{A propósito de tu carta...} (À propos de ta lettre...).
    \item \textbf{\esp{Por lo general} vs \esp{En general} :} \esp{Por lo general} = « généralement » (fréquence). \esp{En general} = « en général / globalement » (portée du propos).
    \item \textbf{\esp{A tiempo} vs \esp{En tiempo} :} \esp{Llegué a tiempo} (j suis arrivé à l'heure / à temps). \esp{En tiempo} = « dans les temps » (délai) ou terme musical. \textbf{Erreur :} \esp{*Llegué en tiempo}.
    \item \textbf{\esp{De hecho} :} « En fait / Effectivement ». Ne pas confondre avec \esp{En hecho} (incorrect) ni traduire par « De fait » (trop soutenu/littéraire).
    \item \textbf{\esp{A menudo} vs \esp{Muchas veces} :} \esp{A menudo} (registre soutenu, écrit). \esp{Muchas veces} (oral, courant). Au Bac écrit : privilégier \esp{a menudo}.
    \item \textbf{\esp{De vez en cuando} :} « De temps en temps / Parfois » (irrégulier). Ne pas utiliser pour « parfois » au sens générique (\esp{a veces}).
\end{enumerate}
\end{attention}

\courssub{L'essai argumentatif : structure et connecteurs}

\begin{methode}[Structurer un essai argumentatif (Ensayo argumentativo)]
\textbf{Objectif :} Produire un texte cohérent, structuré et persuasif sur un thème imposé (ici l'environnement), en respectant le canevas académique hispanophone.
\begin{enumerate}
    \item \textbf{Analyser le sujet :} Identifier les mots-clés (ex. \esp{cambio climático, deforestación, responsabilidad}), la thèse implicite ou explicite, et le type de plan (dialectique : thèse/antithèse/synthèse OU thématique : causes/conséquences/solutions).
    \item \textbf{Rédiger l'introduction (3 mouvements) :}
    \begin{itemize}
        \item \textit{Amorce (Gancho) :} Actualité, citation, statistique générale, observation locale (\esp{En los últimos años...}, \esp{Es innegable que...}).
        \item \textit{Sujet amené / Définition :} Délimiter le thème (\esp{El medioambiente en Camerún...}).
        \item \textit{Problématique + Annonce du plan :} \esp{¿Cómo podemos...?} $\rightarrow$ \esp{En este ensayo, analizaremos primero...; luego...; finalmente...}
    \end{itemize}
    \item \textbf{Développer le corps (2 ou 3 paragraphes arguments) :} 1 Idée directrice = 1 Paragraphe. Structure \textbf{PEEL} : \textbf{P}oint (phrase thèse), \textbf{E}xplication, \textbf{E}vidence (exemple Cameroun/Monde hispanique), \textbf{L}ink (transition).
    \item \textbf{Utiliser les connecteurs logiques (marqueurs de relation) :} Les classer par fonction (voir tableau \cle{Connecteurs}).
    \item \textbf{Conclure (2 mouvements) :} \textit{Bilan synthétique} (\esp{En definitiva, en conclusión, para terminar}) + \textit{Ouverture / Appel à l'action} (\esp{Es urgente la acción colectiva...}, \esp{Depende de nosotros...}).
\end{enumerate}
\end{methode}

\begin{aretenir}[Tableau de référence : Connecteurs logiques (Conectores lógicos)]
\begin{center}
\begin{tabularx}{\linewidth}{|X|X|X|}
\hline
\rowcolor{popblueL} \textbf{Fonction} & \textbf{Connecteurs (Espagnol)} & \textbf{Français / Syntaxe} \\ \hline
\textbf{Addition} & \esp{además, también, asimismo, igualmente, de igual modo, por añadidura} & De plus, aussi, de même. (Souvent en tête de phrase + virgule). \\ \hline
\textbf{Opposition / Concession} & \esp{sin embargo, no obstante, ahora bien, en cambio, por el contrario, aunque + ind./subj.} & Cependant, pourtant, bien que. \esp{Sin embargo} = adv. (point/virgule). \esp{Aunque} = conj. \\ \hline
\textbf{Cause / Explication} & \esp{porque, puesto que, ya que, debido a que, como, pues} & Parce que, car, comme. + \textbf{Indicatif}. \\ \hline
\textbf{Conséquence / Résultat} & \esp{por tanto, por consiguiente, en consecuencia, así que, de modo que, entonces} & Donc, par conséquent, alors. \textbf{Adverbes} $\rightarrow$ Point ou point-virgule avant. \\ \hline
\textbf{Illustration / Exemple} & \esp{por ejemplo, por poner un caso, así, como, v.gr. (verbi gratia)} & Par exemple, ainsi. \\ \hline
\textbf{Ordre / Énumération} & \esp{en primer lugar, en segundo lugar, luego, después, finalmente, para empezar} & En premier lieu, ensuite, enfin. \\ \hline
\textbf{Conclusion / Synthèse} & \esp{en definitiva, en conclusión, en resumen, para terminar, en síntesis} & En définitive, en conclusion, pour finir. \\ \hline
\textbf{Condition / Hypothèse} & \esp{si + ind. (réel), si + imperf. subj. (irréel), con tal de que + subj., siempre que + subj.} & Si, pourvu que, à condition que. \\ \hline
\textbf{But / Finalité} & \esp{para + inf., a fin de que + subj., con el objetivo de + inf.} & Pour, afin que, dans le but de. \\ \hline
\end{tabularx}
\end{center}
\end{aretenir}

\begin{exoresolu}[Transformation et insertion de connecteurs]
\textbf{Énoncé :}
\begin{enumerate}
    \item Reliez les deux phrases par le connecteur logique indiqué (modifiez la ponctuation/majuscules si besoin).
    \begin{itemize}
        \item[a.] \esp{La deforestación avanza. \textbf{(Causa)} El cacao se cultiva sin sombra.}
        \item[b.] \esp{El gobierno prohíbe la caza furtiva. \textbf{(Concesión)} Los furtivos entran en los parques.}
        \item[c.] \esp{Los jóvenes plantan árboles. \textbf{(Consecuencia)} El aire mejora en la ciudad.}
    \end{itemize}
    \item Complétez le paragraphe avec les connecteurs de la liste : \esp{\{por tanto, sin embargo, por ejemplo, además, en conclusión\}}.
    \esp{La contaminación plástica ahoga nuestros ríos. \_\_\_\_\_\_\_\_\_\_, en los mercados de Douala se venden bolsas de un solo uso. \_\_\_\_\_\_\_\_\_\_, muchas terminan en el mar. \_\_\_\_\_\_\_\_\_\_, las tortugas las confunden con medusas. \_\_\_\_\_\_\_\_\_\_, urge una ley que las prohíba. \_\_\_\_\_\_\_\_\_\_, la solución requiere voluntad política y cambio cultural.}
\end{enumerate}
\tcblower
\textbf{Solution détaillée :}

\textbf{1. Relier les phrases (Syntaxe des connecteurs)}

\begin{itemize}
    \item \textbf{[a.] Causa :} \esp{La deforestación avanza \textbf{porque} el cacao se cultiva sin sombra.} (Conjonction \esp{porque} + Indicatif). \textit{Alternative formelle :} \esp{Puesto que el cacao se cultiva sin sombra, la deforestación avanza.}
    \item \textbf{[b.] Concesión :} \esp{El gobierno prohíbe la caza furtiva, \textbf{pero / sin embargo} los furtivos entran en los parques.} (Adversatif coordonné \esp{pero} ou adverbe \esp{sin embargo} $\rightarrow$ point-virgule ou point + majuscule : \esp{...furtiva; sin embargo, los furtivos...}). \textit{Subjonctif requis si « Aunque » :} \esp{Aunque el gobierno prohíba la caza furtiva, los furtivos entran...} (fait réel $\rightarrow$ indicatif possible \esp{prohíbe}, mais concession souvent subjonctif).
    \item \textbf{[c.] Consecuencia :} \esp{Los jóvenes plantan árboles; \textbf{por consiguiente / por tanto / en consecuencia}, el aire mejora en la ciudad.} (Connecteur de conséquence $\rightarrow$ point-virgule ou nouvelle phrase). \textit{Éviter :} \esp{Porque} (cause inverse).
\end{itemize}

\textbf{2. Texte à trous (Cohésion textuelle)}

\esp{La contaminación plástica ahoga nuestros ríos. \textbf{Además}, en los mercados de Douala se venden bolsas de un solo uso. \textbf{Por tanto / Por consiguiente}, muchas terminan en el mar. \textbf{Por ejemplo}, las tortugas las confunden con medusas. \textbf{En conclusión / Por tanto}, urge una ley que las prohíba. \textbf{En conclusión}, la solución requiere voluntad política y cambio cultural.}

\textbf{Justifications :}
\begin{itemize}
    \item \esp{Además} : Ajout d'un argument aggravant (marchés $\rightarrow$ source pollution).
    \item \esp{Por tanto / Por consiguiente} : Conséquence logique (vente $\rightarrow$ fin dans mer).
    \item \esp{Por ejemplo} : Illustration concrète (tortues).
    \item \esp{En conclusión / Por tanto} : Conclusion partielle / Appel à l'action (urgence loi).
    \item \esp{En conclusión} : Synthèse finale globale.
\end{itemize}
\textbf{Attention :} \esp{Por tanto} et \esp{Por consiguiente} exigent un point ou point-virgule avant eux. \esp{Además} en tête de phrase prend une virgule.
\end{exoresolu}

\courssub{Lexique du medioambiente et modèle d'essai}

\begin{exoresolu}[Rédaction de l'introduction et plan détaillé]
\textbf{Énoncé :} Sujet : \esp{« La protección del medioambiente en Camerún: ¿responsabilidad del Estado o de los ciudadanos? »} (La protection de l'environnement au Cameroun : responsabilité de l'État ou des citoyens ?).
\begin{enumerate}
    \item Proposez un plan dialectique en 3 parties (Thèse / Antithèse / Synthèse) avec les idées principales en espagnol.
    \item Rédigez l'\textbf{introduction complète} (environ 80-100 mots) en respectant les 3 mouvements.
\end{enumerate}
\tcblower
\textbf{Solution détaillée :}

\textbf{1. Plan dialectique (Esquema)}

\begin{center}
\begin{tabularx}{\linewidth}{|X|X|}
\hline
\rowcolor{popblueL} \textbf{Parte} & \textbf{Contenido / Ideas fuerza} \\ \hline
\textbf{Introducción} & Amorce : Dégradation visible (plastiques, érosion). Sujet : Partage responsabilité. Problématique : ¿Quién debe actuar primero? Plan : Estado $\rightarrow$ Ciudadano $\rightarrow$ Sinergia. \\ \hline
\textbf{Tesis (Estado)} & \textbf{Argumento 1 :} Marco legal y sanciones (\esp{leyes forestales, multas por tala ilegal}). \textbf{Ejemplo :} Ministerio de Medio Ambiente (MINEPDED), lucha contra caza furtiva en parques (Bouba Ndjida). \textbf{Argumento 2 :} Infraestructuras y gestión de residuos (\esp{recogida de basuras en Yaoundé/Douala, plantas de reciclaje}). \\ \hline
\textbf{Antítesis (Ciudadano)} & \textbf{Argumento 1 :} Cambio de hábitos diarios (\esp{reducir plástico, no tirar basura en desagües $\rightarrow$ inundaciones}). \textbf{Ejemplo :} Iniciativas juveniles \esp{« Limpiemos Douala »}. \textbf{Argumento 2 :} Presión social y voto verde (\esp{exigir cuentas a empresas cacao/madera}). \\ \hline
\textbf{Síntesis (Corresponsabilidad)} & \textbf{Argumento :} Educación ambiental (escuela, medios) + Economía circular. \textbf{Ejemplo :} Proyectos agroforestería cacao sostenible (cooperativas + Estado). \\ \hline
\textbf{Conclusión} & Resumen: Ni solo Estado ni solo pueblo. Llamada: \esp{« El futuro verde de Camerún se construye juntos »}. \\ \hline
\end{tabularx}
\end{center}

\textbf{2. Introduction rédigée (Modelo)}

\begin{textebox}[Introducción - Ensayo argumentativo]
\esp{En los últimos años, las imágenes de los ríos Wouri o Mfoundi asfixiados por botellas plásticas y la tala ilegal en el bosque del Congo se han vuelto tristemente familiares para los cameruneses. Es innegable que la degradación del medioambiente amenaza nuestra seguridad alimentaria y nuestra salud. Ante esta emergencia, surge una pregunta crucial: \textbf{¿recae la protección del entorno principalmente sobre el Estado mediante sus leyes, o sobre la conciencia ciudadana en sus actos cotidianos?} En este ensayo, analizaremos primero el papel irrenunciable de las instituciones públicas; luego, la fuerza transformadora de la responsabilidad individual; para finalmente defender una corresponsabilidad activa como única vía sostenible.}
\end{textebox}

\textbf{Traduction / Glossaire :}
« Ces dernières années, les images des fleuves Wouri ou Mfoundi asphyxiés par des bouteilles en plastique et l'exploitation forestière illégale dans le bassin du Congo sont devenues tristement familières pour les Camerounais. Il est indéniable que la dégradation de l'environnement menace notre sécurité alimentaire et notre santé. Face à cette urgence, surgit une question cruciale : la protection de l'environnement incombe-t-elle principalement à l'État par ses lois, ou à la conscience citoyenne dans ses actes quotidiens ? Dans cet essai, nous analyserons d'abord le rôle irremplaçable des institutions publiques ; ensuite, la force transformatrice de la responsabilité individuelle ; pour finalement défendre une co-responsabilité active comme seule voie durable. »

\textbf{Points de méthode validés :} Amorce ancrée (Wouri, Mfoundi, bassin du Congo), \esp{Es innegable que} (structure impersonnelle + subjonctif), Problématique claire (\esp{¿recae...?}), Annonce de plan explicite (\esp{analizaremos primero... luego... para finalmente...}).
\end{exoresolu}

\courssub{Entraînement guidé : thème, version et rédaction}

\begin{methode}[Rédiger un essai argumentatif complet (Producción escrita)]
\textbf{Objectif :} Mobiliser lexique, grammaire (subjonctif, si-clauses, passif, locutions), structure et connecteurs pour produire un texte de 200-250 mots noté sur 20 (Bac : 10 pts fond / 10 pts forme).
\begin{enumerate}
    \item \textbf{Brouillon (15 min) :} Cerveau en ébullition (mots-clés), Choix du plan, Sélection 2-3 arguments forts + exemples Cameroun/Monde hispanique, Liste connecteurs à caser.
    \item \textbf{Rédaction soignée (45 min) :} Respecter les paragraphes (saut de ligne visible). \textbf{Intro} (1 paragraphe). \textbf{Corps} (2-3 paragraphes distincts). \textbf{Conclusion} (1 paragraphe).
    \item \textbf{Surveillance grammaticale (Check-list) :}
    \begin{itemize}
        \item Accords (GN : \esp{los problemas ambientales son graves}).
        \item Subjonctif après \esp{es necesario que, urge que, exijo que, para que, aunque} (hypothèse).
        \item Si-clauses : \esp{Si + Présent $\rightarrow$ Futur} (réel) / \esp{Si + Imparfait subj. $\rightarrow$ Conditionnel} (irréel présent).
        \item Locutions adverbiales variées (\esp{a menudo, de hecho, en cambio, poco a poco}).
        \item Passif réfléchi (\esp{se prohíbe, se recicla}) ou periphrastique (\esp{es prohibido}).
    \end{itemize}
    \item \textbf{Relecture (10 min) :} Lire à voix basse (détecter les manques d'accords, mots oubliés). Vérifier les accents (\esp{acción, nación, árbol, río, también, sí (adverbe) vs si (conjonction)}).
    \item \textbf{Comptage mots :} Viser 220 ± 10\%.
\end{enumerate}
\end{methode}

\begin{exoresolu}[Essai argumentatif complet : Sujet d'entraînement]
\textbf{Énoncé :} \esp{« Escribe un ensayo argumentativo de unas 220 palabras sobre el tema: "La gestión de los residuos plásticos en nuestras ciudades: ¿problema técnico o cambio de mentalidad?" »}
\tcblower
\textbf{Proposition de rédaction (Modèle corrigé - Niveau Bac Excellent)}

\begin{textebox}[Ensayo argumentativo - Modelo de respuesta]
\esp{La montaña de basura que crece en los vertederos de Nkolfoulou (Yaoundé) o de PK10 (Douala) no es solo un fracaso técnico, es el espejo de nuestros hábitos. Aunque los ayuntamientos invierten en camiones compactadores y contenedores, los canales de drenaje siguen atascados de bolsas plásticas cada temporada de lluvias, agravando las inundaciones. ¿Es, pues, la gestión de residuos un problema de medios o de mentalidad?}

\esp{En primer lugar, la infraestructura es indudablemente necesaria. \textbf{Por consiguiente}, el Estado debe dotarse de plantas de reciclaje modernas y de una recogida selectiva eficaz, como existe en Ruanda o en ciudades europeas. \textbf{Sin embargo}, la experiencia demuestra que la técnica sola no basta: en muchos barrios, los contenedores nuevos son vandalizados o usados como depósitos de escombros.}

\esp{\textbf{Por tanto}, el cambio de mentalidad es la clave profunda. \textbf{De hecho}, la cultura del "usar y tirar" y la falta de educación ambiental desde la escuela primaria perpetúan el ciclo. \textbf{Por ejemplo}, el hábito de tirar la botella por la ventanilla del "bendskin" o del taxi es un acto individual con consecuencia colectiva. \textbf{A propósito}, iniciativas ciudadanas como "Plastic Free July" en Douala demuestran que la concienciación reduce el volumen de basura.}

\esp{\textbf{En definitiva}, no hay dicotomía real: la técnica sin civismo fracasa, y el civismo sin técnica se agota. \textbf{Poco a poco}, una legislación estricta (prohibición de bolsas finas), una economía circular (compra de plástico a recolectores) y una educación obligatoria en las escuelas construirán ciudades limpias. \textbf{Depende de nosotros} transformar el residuo en recurso.}
\end{textebox}

\textbf{Analyse du modèle (Grille d'évaluation) :}
\begin{itemize}
    \item \textbf{Structure :} Intro (Amorce locale Nkolfoulou/PK10 + Problématique + Plan). Corps : 3 paragraphes (Technique / Mentalité / Synthèse). Conclusion (Synthèse + Ouverture action).
    \item \textbf{Connecteurs :} \esp{Aunque} (conc. + ind.), \esp{Por consiguiente}, \esp{Sin embargo}, \esp{Por tanto}, \esp{De hecho}, \esp{Por ejemplo}, \esp{A propósito}, \esp{En definitiva}, \esp{Poco a poco}. (9 connecteurs variés).
    \item \textbf{Grammaire :} Subjonctifs (\esp{exista, prohíba, sea, construyan}) ; Infinitif sujet (\esp{transformar}). Si-clause absente mais \esp{Aunque} + indicatif utilisé. Passif réfléchi (\esp{se agota}). Accords parfaits.
    \item \textbf{Lexique :} \esp{vertedero, canal de drenaje, atascado, contenedor selectiva, planta de reciclaje, vandalizar, escombro, usar y tirar, concienciación, legislación, economía circular, recolector}.
    \item \textbf{Contexte :} Cameroun (Yaoundé, Douala, bendskin, Nkolfoulou, PK10), Référence Rwanda (modèle africain).
    \item \textbf{Comptage :} $\approx$ 235 mots. Dans la norme.
\end{itemize}
\end{exoresolu}

\begin{exercice}[Entraînement Thème / Version — Locutions adverbiales]
\textbf{Consigne :} Traduisez les phrases suivantes.
\begin{enumerate}
    \item \textbf{Version :} \esp{Los voluntarios limpiaron la playa \textbf{a mano} \textbf{poco a poco}. \textbf{De hecho}, recogieron \textbf{por lo general} botellas y bolsas.}
    \item \textbf{Version :} \esp{\textbf{De repente}, se escuchó un ruido \textbf{de golpe}. \textbf{En seguida}, todos salieron \textbf{a propósito} para ver.}
    \item \textbf{Thème :} « Généralement, les étudiants vont à l'université à pied. »
    \item \textbf{Thème :} « Il a lu le rapport par cœur, sans le comprendre. »
    \item \textbf{Thème :} « Peu à peu, la rivière s'est asséchée à cause du climat. »
    \item \textbf{Thème :} « Nous agirons exprès pour changer les choses. »
\end{enumerate}
\end{exercice}

\begin{corrige}[Exercice Thème / Version]
\begin{enumerate}
    \item « Les volontaires ont nettoyé la plage \textbf{à la main} \textbf{peu à peu}. \textbf{En fait}, ils ont ramassé \textbf{généralement} des bouteilles et des sacs. »
    \item « \textbf{Soudain}, on a entendu un bruit \textbf{brutalement / d'un coup}. \textbf{Tout de suite}, tout le monde est sorti \textbf{exprès} pour voir. »
    \item \esp{Por lo general, los estudiantes van a la universidad a pie.}
    \item \esp{Leyó el informe de memoria, sin entenderlo.}
    \item \esp{Poco a poco, el río se secó a causa del clima.}
    \item \esp{Actuaremos a propósito para cambiar las cosas.}
\end{enumerate}
\end{corrige}

\begin{exercice}[Entraînement Rédaction — Sujet d'examen]
\textbf{Sujet :} \esp{« El cambio climático afecta la agricultura en Camerún: ¿qué soluciones propondrían ustedes? »} (Le changement climatique affecte l'agriculture au Cameroun : quelles solutions proposeriez-vous ?).
\textbf{Consignes :}
\begin{itemize}
    \item Rédigez un essai argumentatif d'environ 220 mots.
    \item Utilisez au minimum 6 connecteurs logiques variés (surligne-les).
    \item Intégrez au moins 4 locutions adverbiales étudiées (surligne-les).
    \item Mobilisez le lexique du medioambiente et le contexte camerounais (cacao, café, Nord, Sahel, pluies, sécheresse).
    \item Vérifiez : subjonctif (si-clauses, but, nécessité), accords, accents.
\end{itemize}
\end{exercice}

\begin{corrige}[Exercice Rédaction — Pistes de correction]
\textbf{Plan suggéré (Thématique Causes/Solutions) :}
\begin{itemize}
    \item \textbf{Intro :} Amorce (rendements cacao/café en baisse, Nord menacé par désertification). Problématique : Urgence adaptation. Plan : Constats $\rightarrow$ Solutions techniques $\rightarrow$ Solutions structurelles.
    \item \textbf{Corps 1 (Constats) :} \esp{Por lo general}, pluies irrégulières (\esp{de vez en cuando} torrentielles, \esp{luego} sécheresse). \esp{De hecho}, sols appauvris, parasites. \esp{A propósito}, petits producteurs vulnérables.
    \item \textbf{Corps 2 (Solutions techniques) :} \esp{En primer lugar}, agroforesterie (\esp{por ejemplo}, cacaoyers sous ombre). \esp{Además}, variétés résistantes, irrigation goutte-à-goutte. \esp{Sin embargo}, coût élevé pour paysan.
    \item \textbf{Corps 3 (Solutions structurelles) :} \esp{Por tanto}, État : subventions, assurance récolte, formation. \esp{En cambio}, ONG/Coopératives : encadrement. \esp{Poco a poco}, changement mentalité.
    \item \textbf{Conclusion :} \esp{En definitiva}, résilience = technique + politique + éducation. \esp{Depende de nosotros} (gouvernement, paysan, consommateur).
\end{itemize}
\textbf{Connecteurs cibles :} \esp{Por lo general, De hecho, A propósito, Por ejemplo, Además, Sin embargo, Por tanto, En cambio, Poco a poco, En definitiva, Depende de nosotros}.
\textbf{Locutions cibles :} \esp{Por lo general, De vez en cuando, Poco a poco, A propósito, De hecho, En seguida}.
\end{corrige}

\begin{attention}[Les 5 erreurs qui coûtent des points au Bac]
\begin{enumerate}
    \item \textbf{Confusion \esp{Por qué / Por que / Porqué / Porque} :} \cle{Por qué} (question/exclamation, 2 mots, accent) ; \cle{Porque} (cause, 1 mot, sans accent) ; \cle{Por que} (prep + pronom relatif ou conjonction) ; \cle{Porqué} (nom masculin = « la raison », avec article). \textit{Ex. :} \esp{No sé \textbf{por qué} llueve.} / \esp{Llueve \textbf{porque} hay nubes.} / \esp{El \textbf{porqué} de la inundación es la basura.}
    \item \textbf{Oubli de l'accent sur les adverbes en \esp{-mente} :} L'accent se met \textbf{uniquement sur l'adjectif féminin} d'origine. \esp{Rápida $\rightarrow$ rápidamente} (accent sur \esp{á}). \esp{Fácil $\rightarrow$ fácilmente} (pas d'accent sur \esp{fácil} $\rightarrow$ pas d'accent sur \esp{fácilmente}). \esp{Útil $\rightarrow$ útilmente}. \textbf{Piège :} \esp{Por supuesto} (pas \esp{por supesto}), \esp{De repente} (pas \esp{de repénte}).
    \item \textbf{Calque du français sur « \esp{A tiempo} / \esp{En tiempo} » :} \esp{A tiempo} = « à l'heure / à temps » (ponctualité). \esp{En tiempo} = « dans les temps » (délai) ou terme musical. \textit{Erreur :} \esp{Llegué en tiempo} pour « J suis arrivé à l'heure ». \textbf{Correct :} \esp{Llegué \textbf{a tiempo}.}
    \item \textbf{Mauvaise concordance des temps avec \esp{Si} (Si-clauses) :} \textbf{JAMAIS} d'indicatif présent + conditionnel (\esp{*Si llueve, iría} $\times$). \textbf{Règle stricte :} \esp{Si + Présent $\rightarrow$ Futur} (Réalité) ; \esp{Si + Imparfait Subjonctif $\rightarrow$ Conditionnel} (Hypothèse). \textit{Correct :} \esp{Si llueve, \textbf{iré} a casa.} / \esp{Si \textbf{lloviera}, \textbf{iría} a casa.}
    \item \textbf{Accord du participe passé / adjectif en \esp{-ado/-ido} utilisé comme attribut :} \esp{La calle está \textbf{contaminadA}} (accord avec \esp{calle} f.sg). \esp{Los ríos están \textbf{contaminadOS}} (m.pl). \textit{Erreur fréquente :} \esp{La ciudad está contaminad\emph{o}} (masculin par défaut). \textbf{Règle :} Participe passé = adjectif $\rightarrow$ accord en genre et nombre avec le sujet.
\end{enumerate}
\end{attention}

\newpage

\coursec{Exercices}

\courssub{Je vérifie mes connaissances}

\begin{exercice}[QCM : locutions adverbiales]
Pour chaque phrase, choisis la locution adverbiale qui convient parmi les trois propositions.

\begin{enumerate}
\item \esp{... sonó el timbre y todos los alumnos salieron corriendo del aula.}
\begin{enumerate}
\item \esp{De repente}
\item \esp{Poco a poco}
\item \esp{De vez en cuando}
\end{enumerate}
\item \esp{Mi abuela va al mercado de Mvog-Mbi ... : casi todos los días.}
\begin{enumerate}
\item \esp{a lo lejos}
\item \esp{a menudo}
\item \esp{en vano}
\end{enumerate}
\item \esp{... terminó la lluvia y pudimos jugar al fútbol.}
\begin{enumerate}
\item \esp{Por fin}
\item \esp{Por lo tanto}
\item \esp{Sin querer}
\end{enumerate}
\item \esp{El niño rompió el florero ... : no fue un accidente.}
\begin{enumerate}
\item \esp{a propósito}
\item \esp{sin embargo}
\item \esp{de prisa}
\end{enumerate}
\end{enumerate}
\end{exercice}

\begin{exercice}[Vrai ou faux, en justifiant]
Dis si les affirmations suivantes sont vraies ou fausses, puis justifie ta réponse par une traduction ou un contre-exemple.

\begin{enumerate}
\item \esp{A menudo} signifie « de temps en temps ».
\item \esp{En seguida} et \esp{inmediatamente} sont des synonymes.
\item \esp{Poco a poco} est une locution adverbiale de temps.
\item Dans un essai argumentatif, \esp{sin embargo} sert à introduire un argument contraire.
\item \esp{A propósito} signifie toujours « à propos ».
\end{enumerate}
\end{exercice}

\begin{exercice}[Vocabulaire : le medioambiente]
Trouve l'équivalent espagnol des mots et expressions français suivants, puis construis une phrase complète avec chacun d'eux.

\begin{enumerate}
\item la déforestation
\item le réchauffement climatique
\item les énergies renouvelables
\item recycler les déchets
\item les sacs en plastique
\item protéger la planète
\end{enumerate}
\end{exercice}

\begin{exercice}[Grammaire : à compléter]
Complète les phrases suivantes avec la locution adverbiale qui convient, choisie dans la liste : \esp{de vez en cuando — en seguida — cada vez más — por lo menos — a pie}.

\begin{enumerate}
\item \esp{Vivo cerca del colegio, así que voy ... todos los días.}
\item \esp{... visito a mis primos de Douala, dos o tres veces al año.}
\item \esp{¡Ven ... ! El autobús está a punto de salir.}
\item \esp{La contaminación del aire es ... preocupante en las grandes ciudades.}
\item \esp{Debes plantar ... diez árboles para compensar los que cortaste.}
\end{enumerate}
\end{exercice}

\courssub{Je m'entraîne}

\begin{exercice}[Classement des locutions]
Range les quinze locutions adverbiales suivantes dans un tableau à quatre colonnes (\textbf{temps / lieu / manière / quantité}) et donne l'équivalent français de chacune.

\begin{center}
\esp{de repente — a menudo — poco a poco — de vez en cuando — en seguida — a lo lejos — cerca de — en todas partes — a propósito — sin querer — a pie — por lo menos — cada vez más — apenas — por fin}
\end{center}
\end{exercice}

\begin{exercice}[Reformulation]
Réécris le paragraphe suivant en remplaçant chaque adverbe simple souligné par une locution adverbiale équivalente (\esp{súbitamente} $\rightarrow$ \esp{de repente}, etc.).

\begin{textebox}[Un día en el mercado]
\esp{El sábado pasado fui al mercado con mi madre. \textbf{Inmediatamente}, vimos a un vendedor de frutas que gritaba \textbf{fuertemente}. \textbf{Súbitamente}, empezó a llover y todos corrieron \textbf{rápidamente} a refugiarse bajo los techos. \textbf{Finalmente}, la lluvia se detuvo y pudimos comprar \textbf{tranquilamente} todo lo que necesitábamos.}
\end{textebox}
\end{exercice}

\begin{exercice}[Texte à trous : les sacs plastiques]
Complète cet essai avec les connecteurs et locutions de la liste : \esp{en primer lugar — sin embargo — por lo tanto — en conclusión — de vez en cuando — además}.

\begin{textebox}[Los sacos plásticos: un peligro para nuestro entorno]
\esp{... , los sacos plásticos son muy prácticos y baratos: todo el mundo los usa en los mercados de Yaoundé. ... , representan un grave problema para el medioambiente porque tardan siglos en desaparecer. ... , ensucian las calles y matan a los animales que los tragan. ... , todavía vemos personas que los tiran al suelo ... , sin pensar en las consecuencias. ... , debemos prohibirlos o reemplazarlos por bolsas de tela. ... , proteger la naturaleza es responsabilidad de todos.}
\end{textebox}
\end{exercice}

\begin{exercice}[Appariement français-espagnol]
Associe chaque locution française à son équivalent espagnol, puis traduis les deux phrases qui suivent.

\begin{center}
\begin{tabularx}{\linewidth}{|X|X|}
\hline
\rowcolor{popblueL} \textbf{Français} & \textbf{Español} \\
\hline
1. soudain & a. \esp{de vez en cuando} \\
\hline
2. de temps en temps & b. \esp{por fin} \\
\hline
3. tout de suite & c. \esp{de repente} \\
\hline
4. exprès & d. \esp{en seguida} \\
\hline
5. enfin / finalement & e. \esp{a propósito} \\
\hline
6. au moins & f. \esp{por lo menos} \\
\hline
\end{tabularx}
\end{center}

\begin{enumerate}
\item Soudain, le vent a renversé les poubelles devant le lycée.
\item De temps en temps, nous organisons une journée de nettoyage du quartier.
\end{enumerate}
\end{exercice}

\courssub{Je me prépare au Bac}

\begin{exercice}[Compréhension et étude de langue]
Lis le texte suivant, puis réponds aux questions \textbf{en espagnol}, avec des phrases complètes.

\begin{textebox}[El bosque que desaparece]
\esp{En muchas regiones del mundo, los bosques desaparecen poco a poco. Los agricultores cortan los árboles a menudo para cultivar la tierra, y las empresas explotan la madera sin pensar en el futuro. De repente, un incendio puede destruir en pocas horas lo que la naturaleza construyó durante siglos. Sin embargo, cada vez más jóvenes deciden actuar: plantan árboles, limpian los ríos de vez en cuando y explican a sus familias por qué hay que proteger el medioambiente. Por fin, los gobiernos empiezan a escuchar sus voces y a crear parques naturales. Poco a poco, la esperanza vuelve, pero el camino es todavía largo.}
\end{textebox}

\textbf{Comprensión del texto}
\begin{enumerate}
\item ¿Por qué desaparecen los bosques según el texto? Cita dos razones.
\item ¿Qué hacen los jóvenes para proteger el medioambiente?
\item ¿Cómo sabes que el autor es optimista al final del texto? Justifica tu respuesta.
\end{enumerate}

\textbf{Estudio de la lengua}
\begin{enumerate}
\item Subraya en el texto seis locuciones adverbiales y clasifícalas según su valor (tiempo, lugar, manera, cantidad).
\item Traduce al francés las locuciones \esp{poco a poco}, \esp{de repente}, \esp{sin embargo}, \esp{cada vez más}, \esp{por fin}.
\item Sustituye \esp{a menudo} por un adverbio simple de sentido equivalente.
\end{enumerate}
\end{exercice}

\begin{exercice}[Thème et version]
\textbf{Thème.} Traduis en espagnol les cinq phrases suivantes en utilisant obligatoirement la locution indiquée entre parenthèses.

\begin{enumerate}
\item Soudain, la rivière a débordé et a inondé le quartier. (\esp{de repente})
\item De temps en temps, mon père plante des arbres derrière la maison. (\esp{de vez en cuando})
\item Viens tout de suite, la réunion sur l'environnement commence ! (\esp{en seguida})
\item Il a jeté son sac plastique par terre exprès. (\esp{a propósito})
\item Nous devons recycler au moins la moitié de nos déchets. (\esp{por lo menos})
\end{enumerate}

\textbf{Version.} Traduis en français le paragraphe suivant.

\begin{textebox}[La Amazonía en peligro]
\esp{La selva amazónica desaparece poco a poco ante los ojos del mundo. Los incendios se producen a menudo durante la estación seca y, de repente, miles de hectáreas se convierten en cenizas. Sin embargo, cada vez más países comprenden la importancia de este bosque para el clima del planeta. Por fin, algunos gobiernos han decidido protegerlo con leyes más estrictas.}
\end{textebox}
\end{exercice}

\begin{exercice}[Rédaction type Bac : ensayo argumentativo]
\textbf{Sujet :} \esp{¿Es la contaminación el mayor problema de nuestro planeta? Justifica tu opinión con argumentos y ejemplos concretos.}

Rédige un essai argumentatif en espagnol de \textbf{120 à 150 mots}. Tu respecteras la structure suivante :

\begin{enumerate}
\item \textbf{Introducción :} presenta el tema y tu tesis (una o dos frases).
\item \textbf{Desarrollo :} dos o tres argumentos con ejemplos, organizados con los conectores \esp{en primer lugar}, \esp{además}, \esp{sin embargo}, \esp{por lo tanto}.
\item \textbf{Conclusión :} resume tu posición con \esp{en conclusión} o \esp{para terminar}.
\end{enumerate}

Utilise au moins \textbf{trois locutions adverbiales} étudiées dans la leçon et souligne-les.
\end{exercice}

\newpage

\coursec{Corrigés des exercices}

\courssub{Corrigés des exercices d'entraînement : locutions adverbiales et essai argumentatif}

\begin{corrige}[Exercice 1 — QCM : locutions adverbiales]
\begin{enumerate}
\item \textbf{Réponse a :} \esp{De repente sonó el timbre y todos los alumnos salieron corriendo del aula.} « Soudain, la sonnerie a retenti et tous les élèves sont sortis en courant de la salle. » \esp{De repente} (manière/temps, soudaineté) convient ; \esp{poco a poco} (progression lente) et \esp{de vez en cuando} (fréquence) sont incompatibles avec l'action brusque.
\item \textbf{Réponse b :} \esp{Mi abuela va al mercado de Mvog-Mbi a menudo: casi todos los días.} « Ma grand-mère va souvent au marché de Mvog-Mbi. » La précision « casi todos los días » exige une locution de fréquence : \esp{a menudo}. \esp{A lo lejos} = « au loin » ; \esp{en vano} = « en vain ».
\item \textbf{Réponse a :} \esp{Por fin terminó la lluvia y pudimos jugar al fútbol.} « Enfin, la pluie a cessé et nous avons pu jouer au football. » \esp{Por fin} exprime le soulagement après une attente. \esp{Por lo tanto} = « donc » (conséquence logique) ; \esp{sin querer} = « sans le vouloir ».
\item \textbf{Réponse a :} \esp{El niño rompió el florero a propósito: no fue un accidente.} « L'enfant a cassé le vase exprès. » La négation de l'accident impose \esp{a propósito} (= exprès, intentionnellement).
\end{enumerate}
\end{corrige}

\begin{corrige}[Exercice 2 — Vrai ou faux, en justifiant]
\begin{enumerate}
\item \textbf{Faux.} \esp{A menudo} signifie « souvent », pas « de temps en temps ». Contre-exemple : \esp{Voy a menudo al cine} « Je vais souvent au cinéma ». « De temps en temps » se traduit par \esp{de vez en cuando}.
\item \textbf{Vrai.} \esp{En seguida} et \esp{inmediatamente} signifient tous deux « immédiatement, tout de suite » : \esp{Vengo en seguida} = \esp{Vengo inmediatamente} « Je viens tout de suite ».
\item \textbf{Faux.} \esp{Poco a poco} (« petit à petit ») est une locution adverbiale de \textbf{manière} : elle décrit la façon progressive dont se déroule l'action, pas un moment ni une fréquence. Exemple : \esp{Aprendo español poco a poco} « J'apprends l'espagnol petit à petit ».
\item \textbf{Vrai.} \esp{Sin embargo} (« cependant, néanmoins ») est un connecteur d'opposition : il introduit un argument ou un fait contraire à ce qui précède. Exemple : \esp{El plástico es útil; sin embargo, contamina los océanos.}
\item \textbf{Faux.} \esp{A propósito} a deux sens : « exprès » (\esp{Lo hizo a propósito} « Il l'a fait exprès ») et « à propos / soit dit en passant » pour changer de sujet (\esp{A propósito, ¿viste el partido?}). Il ne signifie donc pas « toujours » à propos.
\end{enumerate}
\end{corrige}

\begin{corrige}[Exercice 3 — Vocabulaire : le medioambiente]
\begin{enumerate}
\item la déforestation = \esp{la deforestación}. \esp{La deforestación destruye el hábitat de muchos animales.} « La déforestation détruit l'habitat de nombreux animaux. »
\item le réchauffement climatique = \esp{el calentamiento global}. \esp{El calentamiento global aumenta cada vez más por culpa de los gases contaminantes.}
\item les énergies renouvelables = \esp{las energías renovables}. \esp{Las energías renovables, como el sol y el viento, no contaminan el aire.}
\item recycler les déchets = \esp{reciclar los desechos / la basura}. \esp{Debemos reciclar los desechos en lugar de quemarlos.}
\item les sacs en plastique = \esp{las bolsas de plástico}. \esp{Las bolsas de plástico tardan siglos en desaparecer de la naturaleza.}
\item protéger la planète = \esp{proteger el planeta}. \esp{Proteger el planeta es el deber de todos los ciudadanos.}
\end{enumerate}
\textbf{Points de vigilance :} \esp{el calentamiento global} prend l'article masculin (\esp{el}) ; \esp{reciclar} est régulier ; ne pas écrire \esp{plásticos} pour « plastique » seul : on dit \esp{bolsas de plástico}.
\end{corrige}

\begin{corrige}[Exercice 4 — Grammaire : à compléter]
\begin{enumerate}
\item \esp{Vivo cerca del colegio, así que voy \textbf{a pie} todos los días.} « J'habite près du collège, donc j'y vais à pied tous les jours. » (lieu/manière de déplacement)
\item \esp{\textbf{De vez en cuando} visito a mis primos de Douala, dos o tres veces al año.} « De temps en temps, je rends visite à mes cousins de Douala. » (fréquence : « dos o tres veces al año »)
\item \esp{¡Ven \textbf{en seguida}! El autobús está a punto de salir.} « Viens tout de suite ! Le bus est sur le point de partir. » (immédiateté)
\item \esp{La contaminación del aire es \textbf{cada vez más} preocupante en las grandes ciudades.} « La pollution de l'air est de plus en plus inquiétante. » (quantité/degré croissant)
\item \esp{Debes plantar \textbf{por lo menos} diez árboles para compensar los que cortaste.} « Tu dois planter au moins dix arbres. » (quantité minimale)
\end{enumerate}
\end{corrige}

\begin{corrige}[Exercice 5 — Classement des locutions]
\begin{center}
\begin{tabularx}{\linewidth}{|X|X|X|X|}
\hline
\rowcolor{popblueL} \textbf{Tiempo} & \textbf{Lugar} & \textbf{Manera} & \textbf{Cantidad} \\
\hline
\esp{de repente} (soudain) & \esp{a lo lejos} (au loin) & \esp{poco a poco} (petit à petit) & \esp{por lo menos} (au moins) \\
\hline
\esp{a menudo} (souvent) & \esp{cerca de} (près de) & \esp{a propósito} (exprès) & \esp{cada vez más} (de plus en plus) \\
\hline
\esp{de vez en cuando} (de temps en temps) & \esp{en todas partes} (partout) & \esp{sin querer} (sans le vouloir) & \esp{apenas} (à peine) \\
\hline
\esp{en seguida} (tout de suite) & & \esp{a pie} (à pied) & \\
\hline
\esp{por fin} (enfin) & & & \\
\hline
\end{tabularx}
\end{center}
\textbf{Remarque :} \esp{de repente} et \esp{por fin} peuvent aussi être classées dans la manière selon les grammaires ; l'essentiel est de justifier le classement par le sens (soudaineté, aboutissement d'une attente).
\end{corrige}

\begin{corrige}[Exercice 6 — Reformulation]
\textbf{Version corrigée du texte :}
\end{corrige}

\begin{textebox}[Un día en el mercado (versión corregida)]
\esp{El sábado pasado fui al mercado con mi madre. \textbf{En seguida}, vimos a un vendedor de frutas que gritaba \textbf{a gritos / con fuerza}. \textbf{De repente}, empezó a llover y todos corrieron \textbf{de prisa} a refugiarse bajo los techos. \textbf{Por fin}, la lluvia se detuvo y pudimos comprar \textbf{con calma} todo lo que necesitábamos.}
\end{textebox}

\begin{corrige}[Exercice 6 — Suite]
« Samedi dernier, je suis allé au marché avec ma mère. Tout de suite, nous avons vu un vendeur de fruits qui criait très fort. Soudain, il a commencé à pleuvoir et tous ont couru vite se réfugier sous les toits. Enfin, la pluie s'est arrêtée et nous avons pu acheter tranquillement tout ce dont nous avions besoin. »

\textbf{Correspondances :} \esp{inmediatamente} → \esp{en seguida} ; \esp{fuertemente} → \esp{a gritos} / \esp{con fuerza} ; \esp{súbitamente} → \esp{de repente} ; \esp{rápidamente} → \esp{de prisa} ; \esp{finalmente} → \esp{por fin} ; \esp{tranquilamente} → \esp{con calma}.
\end{corrige}

\begin{textebox}[Los sacos plásticos: texto completo]
\esp{En primer lugar, los sacos plásticos son muy prácticos y baratos: todo el mundo los usa en los mercados de Yaoundé. Sin embargo, representan un grave problema para el medioambiente porque tardan siglos en desaparecer. Además, ensucian las calles y matan a los animales que los tragan. De vez en cuando, todavía vemos personas que los tiran al suelo sin pensar en las consecuencias. Por lo tanto, debemos prohibirlos o reemplazarlos por bolsas de tela. En conclusión, proteger la naturaleza es responsabilidad de todos.}
\end{textebox}

\begin{corrige}[Exercice 7 — Texte à trous : les sacs plastiques]
\textbf{Ordre logique des connecteurs et locutions :}
\begin{enumerate}
\item \esp{En primer lugar} : ouverture du premier argument (aspect pratique).
\item \esp{Sin embargo} : opposition (le problème écologique).
\item \esp{Además} : ajout d'un argument (rues sales, animaux tués).
\item \esp{De vez en cuando} : fréquence (« todavía vemos personas que los tiran al suelo de vez en cuando »).
\item \esp{Por lo tanto} : conséquence (« debemos prohibirlos... »).
\item \esp{En conclusión} : formule finale (« proteger la naturaleza es responsabilidad de todos »).
\end{enumerate}
\textbf{Traduction :} « Premièrement, les sacs plastiques sont très pratiques et bon marché : tout le monde les utilise dans les marchés de Yaoundé. Cependant, ils représentent un grave problème pour l'environnement car ils mettent des siècles à disparaître. De plus, ils salissent les rues et tuent les animaux qui les avalent. De temps en temps, nous voyons encore des gens les jeter par terre, sans penser aux conséquences. Par conséquent, nous devons les interdire ou les remplacer par des sacs en tissu. En conclusion, protéger la nature est la responsabilité de tous. »
\end{corrige}

\begin{corrige}[Exercice 8 — Appariement français-espagnol]
\textbf{Appariement :} 1-c ; 2-a ; 3-d ; 4-e ; 5-b ; 6-f.

\textbf{Traductions :}
\begin{enumerate}
\item « Soudain, le vent a renversé les poubelles devant le lycée. » $\rightarrow$ \esp{De repente, el viento volcó los cubos de basura delante del instituto.} (Passé simple : \esp{volcó} ; « poubelles » = \esp{los cubos de basura} ; « lycée » = \esp{el instituto}.)
\item « De temps en temps, nous organisons une journée de nettoyage du quartier. » $\rightarrow$ \esp{De vez en cuando, organizamos una jornada de limpieza del barrio.}
\end{enumerate}
\textbf{Point de vigilance :} ne pas traduire « soudain » par \esp{súbitamente} dans le registre courant : \esp{de repente} est la locution attendue ; « enfin » au sens de soulagement = \esp{por fin}, et non \esp{finalmente} (qui marque plutôt l'ordre d'énumération).
\end{corrige}

\begin{corrige}[Exercice 9 — Compréhension et étude de langue]
\textbf{Comprensión del texto (réponses modèles) :}
\begin{enumerate}
\item \esp{Los bosques desaparecen, en primer lugar, porque los agricultores cortan los árboles a menudo para cultivar la tierra y, además, porque las empresas explotan la madera sin pensar en el futuro.} (Les deux raisons : l'agriculture et l'exploitation du bois.)
\item \esp{Los jóvenes plantan árboles, limpian los ríos de vez en cuando y explican a sus familias por qué hay que proteger el medioambiente.}
\item \esp{El autor es optimista al final porque dice que «por fin, los gobiernos empiezan a escuchar sus voces» y que «poco a poco, la esperanza vuelve». Estas frases muestran que la situación mejora.}
\end{enumerate}

\textbf{Estudio de la lengua :}
\begin{enumerate}
\item Six locutions du texte : \esp{poco a poco} (manera), \esp{a menudo} (tiempo/fréquence), \esp{de repente} (tiempo/manière), \esp{sin embargo} (conector de oposición), \esp{cada vez más} (cantidad), \esp{de vez en cuando} (tiempo), \esp{por fin} (tiempo). (On en retient six au choix.)
\item \esp{poco a poco} = petit à petit ; \esp{de repente} = soudain ; \esp{sin embargo} = cependant ; \esp{cada vez más} = de plus en plus ; \esp{por fin} = enfin.
\item \esp{a menudo} → \esp{frecuentemente} (adverbe simple équivalent : « fréquemment »).
\end{enumerate}
\textbf{Barème indicatif (10 pts) :} compréhension 3 $\times$ 1,5 = 4,5 pts ; repérage et classement 2 pts ; traduction des locutions 2,5 pts ; synonyme 1 pt.
\end{corrige}

\begin{corrige}[Exercice 10 — Thème et version]
\textbf{Thème (corrigé) :}
\begin{enumerate}
\item \esp{De repente, el río se desbordó e inundó el barrio.} (Passé simple : \esp{se desbordó}, \esp{inundó} ; « le quartier » = \esp{el barrio}.)
\item \esp{De vez en cuando, mi padre planta árboles detrás de la casa.} (Présent d'habitude ; \esp{detrás de} = derrière.)
\item \esp{¡Ven en seguida, la reunión sobre el medioambiente empieza!} (Impératif \esp{ven} ; \esp{sobre} = sur/au sujet de.)
\item \esp{Tiró su bolsa de plástico al suelo a propósito.} (\esp{tirar} = jeter ; \esp{al suelo} = par terre.)
\item \esp{Debemos reciclar por lo menos la mitad de nuestros desechos.} (\esp{la mitad de} = la moitié de.)
\end{enumerate}

\textbf{Version (traduction modèle) :} « La forêt amazonienne disparaît petit à petit sous les yeux du monde. Les incendies se produisent souvent pendant la saison sèche et, soudain, des milliers d'hectares se changent en cendres. Cependant, de plus en plus de pays comprennent l'importance de cette forêt pour le climat de la planète. Enfin, certains gouvernements ont décidé de la protéger par des lois plus strictes. »

\textbf{Points de vigilance :} \esp{ante los ojos del mundo} = « sous les yeux du monde » ; \esp{se convierten en cenizas} = « se changent en cendres » ; \esp{leyes más estrictas} = « des lois plus strictes ». \textbf{Barème indicatif :} thème 5 $\times$ 2 = 10 pts (locution correcte 1 pt, phrase 1 pt) ; version 10 pts (fidélité 6, qualité du français 4).
\end{corrige}

\begin{textebox}[¿Es la contaminación el mayor problema de nuestro planeta?]
\esp{Hoy en día, muchas personas piensan que la contaminación es el mayor problema de nuestro planeta. En mi opinión, es verdad, y voy a explicar por qué.}

\esp{\textbf{En primer lugar}, la contaminación del aire provoca enfermedades graves, sobre todo en las grandes ciudades, donde la gente respira humo todos los días. \textbf{Además}, el plástico destruye \esp{poco a poco} los océanos: los peces y las tortugas mueren \esp{a menudo} porque tragan bolsas y botellas. \textbf{Sin embargo}, algunos creen que la pobreza es un problema más urgente. Es cierto, pero la contaminación agrava la pobreza: sin agua limpia ni tierra fértil, los campesinos no pueden vivir. \textbf{Por lo tanto}, los dos problemas están unidos.}

\esp{\textbf{En conclusión}, la contaminación es el mayor desafío de nuestro siglo. Debemos actuar \esp{en seguida}, reciclar \esp{por lo menos} una parte de nuestros desechos y proteger la naturaleza, porque el futuro del planeta está en nuestras manos.}
\end{textebox}

\begin{corrige}[Exercice 11 — Rédaction type Bac : ensayo argumentativo]
\textbf{Proposition de rédaction modèle (environ 140 mots) :} voir le texte ci-dessus.

\textbf{Locutions adverbiales soulignées attendues (au moins trois) :} \esp{poco a poco}, \esp{a menudo}, \esp{en seguida}, \esp{por lo menos}.

\textbf{Barème indicatif (20 pts) :} respect de la structure intro/développement/conclusion 4 pts ; connecteurs logiques 4 pts ; au moins trois locutions adverbiales correctement employées 3 pts ; correction grammaticale (accords, temps, accents) 5 pts ; richesse du vocabulaire et pertinence des arguments 4 pts.

\textbf{Points de vigilance :}
\begin{itemize}
\item Ne pas écrire \esp{en primer lugar} sans poursuivre par \esp{además} ou \esp{en segundo lugar} : les connecteurs doivent former une chaîne cohérente.
\item \esp{Sin embargo} s'écrit en un seul mot et introduit toujours une opposition, jamais une conséquence.
\item Attention aux faux amis : « actuellement » = \esp{actualmente / hoy en día}, pas \esp{actualmente} mal orthographié ; « une chance » = \esp{una oportunidad}.
\item Accords : \esp{las bolsas de plástico}, \esp{enfermedades graves}, \esp{leyes estrictas}.
\item Respecter la fourchette 120--150 mots : un essai trop court perd des points de contenu.
\end{itemize}
\end{corrige}

\newpage

\coursec{Fiche bilan}

\begin{popfigure}
\begin{tikzpicture}[
    mindmap,
    concept color=popblue,
    level 1 concept/.append style={level distance=3.5cm, sibling angle=360/6},
    level 2 concept/.append style={level distance=2.5cm},
    every node/.style={font=\small, align=center},
    branch1/.style={concept color=poporange},
    branch2/.style={concept color=popgreen},
    branch3/.style={concept color=poppurple},
    branch4/.style={concept color=poppink},
    branch5/.style={concept color=popgold},
    branch6/.style={concept color=popblue!70!popdark}
]
\node[concept, text=white, minimum size=2.2cm, align=center]{Leçon E10\\ \emph{Traduction \& Redacción}}
    child[branch1] { node[concept, text=white, align=center]{Locutions\\ adverbiales}
        child { node[concept, text=white, scale=0.7] {Temps} }
        child { node[concept, text=white, scale=0.7] {Lieu} }
        child { node[concept, text=white, scale=0.7] {Manière} }
        child { node[concept, text=white, scale=0.7] {Quantité} }
    }
    child[branch2] { node[concept, text=white] {Traduction}
        child { node[concept, text=white, scale=0.7, align=center]{Correspondances\\ Fr/Es} }
        child { node[concept, text=white, scale=0.7, align=center]{Pièges \&\\ Faux amis} }
    }
    child[branch3] { node[concept, text=white, align=center]{Essai\\ argumentatif}
        child { node[concept, text=white, scale=0.7] {Structure} }
        child { node[concept, text=white, scale=0.7, align=center]{Connecteurs\\ logiques} }
    }
    child[branch4] { node[concept, text=white, align=center]{Lexique\\ Medioambiente}
        child { node[concept, text=white, scale=0.7] {Problèmes} }
        child { node[concept, text=white, scale=0.7] {Solutions} }
    }
    child[branch5] { node[concept, text=white] {Entraînement}
        child { node[concept, text=white, scale=0.7] {Thème} }
        child { node[concept, text=white, scale=0.7] {Version} }
        child { node[concept, text=white, scale=0.7] {Rédaction} }
    }
    child[branch6] { node[concept, text=white, align=center]{Évaluation\\ Bac}
        child { node[concept, text=white, scale=0.7] {Compétences} }
        child { node[concept, text=white, scale=0.7] {Checklist} }
    };
\end{tikzpicture}
\legende{Carte mentale de la leçon E10 : locutions adverbiales, traduction et essai argumentatif sur l'environnement}
\end{popfigure}

\coursec{Lexique clé : El medioambiente}

\begin{definition}[El medioambiente : vocabulaire de base]
Le thème de l'environnement est un classique du Baccalauréat (commentaire de texte, thème, version et rédaction). Le tableau ci-dessous regroupe les termes indispensables, avec un exemple en contexte pour chacun.
\end{definition}

\begin{center}
\begin{tabularx}{\linewidth}{|X|X|X|}
\hline
\rowcolor{popblueL} \textbf{Español} & \textbf{Français} & \textbf{Contexte / Exemple} \\
\hline
\esp{el calentamiento global} & le réchauffement climatique & \esp{El calentamiento global amenaza la biodiversidad.} \\
\esp{el cambio climático} & le changement climatique & \esp{Los efectos del cambio climático son cada vez más visibles.} \\
\esp{la contaminación} & la pollution & \esp{La contaminación del aire en las grandes ciudades.} \\
\esp{los gases de efecto invernadero} & les gaz à effet de serre & \esp{Hay que reducir la emisión de gases de efecto invernadero.} \\
\esp{la capa de ozono} & la couche d'ozone & \esp{El agujero en la capa de ozono.} \\
\esp{la deforestación} & la déforestation & \esp{La deforestación de la Amazonía avanza rápidamente.} \\
\esp{la biodiversidad} & la biodiversité & \esp{Proteger la biodiversidad es urgente.} \\
\esp{las energías renovables} & les énergies renouvelables & \esp{Apostar por las energías renovables (solar, eólica).} \\
\esp{el desarrollo sostenible} & le développement durable & \esp{Un modelo de desarrollo sostenible.} \\
\esp{reciclar / la recogida selectiva} & recycler / le tri sélectif & \esp{Debemos reciclar y hacer la recogida selectiva.} \\
\esp{la huella de carbono} & l'empreinte carbone & \esp{Reducir nuestra huella de carbono.} \\
\esp{el residuo / el desecho} & le déchet & \esp{La gestión de los residuos tóxicos.} \\
\hline
\end{tabularx}
\end{center}

\begin{attention}[Genre et emploi : pièges lexicaux]
\begin{itemize}
    \item \esp{el medioambiente} s'écrit aussi \esp{el medio ambiente} (les deux graphies sont admises) ; il est toujours masculin : \esp{el medioambiente está en peligro}.
    \item \esp{la contaminación} est le terme général ; \esp{la polución} existe mais est moins fréquent.
    \item « Menacer » se traduit par \esp{amenazar} (et non par un calque de « menacer » avec \esp{*menazar}, qui n'existe pas).
    \item « Protéger » = \esp{proteger} (g $\rightarrow$ j devant a, o : \esp{protejo, proteja}).
    \item \esp{apostar por} = miser sur, opter pour : \esp{Apostamos por las energías renovables.} « Nous misons sur les énergies renouvelables. »
\end{itemize}
\end{attention}

\coursec{Les locutions adverbiales : classement et sens}

\begin{propriete}[Locutions adverbiales : forme et emploi]
Une \cle{locution adverbiale} (\esp{locución adverbial}) est un groupe de mots figé (le plus souvent \textbf{préposition + nom, adjectif ou adverbe}) qui fonctionne comme un adverbe unique. Elle est \textbf{invariable} et précise le temps, le lieu, la manière ou la quantité.\\
Exemple : \esp{De repente, empezó a llover.} « Soudain, il s'est mis à pleuvoir. »
\end{propriete}

\begin{center}
\begin{tabularx}{\linewidth}{|>{\bfseries}X|X|X|}
\hline
\rowcolor{poporangeL} \textbf{Catégorie} & \textbf{Locuciones españolas (exemples)} & \textbf{Correspondances françaises} \\
\hline
\textbf{TIEMPO (Temps)} & \esp{de repente, de pronto, al instante, en seguida, a menudo, con frecuencia, de vez en cuando, a veces, rara vez, poco a poco, paso a paso, al final, al principio, hoy en día, antaño} & soudain, tout à coup, immédiatement, aussitôt, souvent, fréquemment, de temps en temps, parfois, rarement, peu à peu, étape par étape, finalement, au début, de nos jours, autrefois \\
\hline
\textbf{LUGAR (Lieu)} & \esp{al lado, a la derecha, a la izquierda, enfrente, detrás, encima, debajo, dentro, fuera, por todas partes, en ninguna parte, a lo lejos} & à côté, à droite, à gauche, en face, derrière, dessus, dessous, dedans, dehors, partout, nulle part, au loin \\
\hline
\textbf{MODO (Manière)} & \esp{así, de esta manera, a propósito, por casualidad, a mano, a máquina, en secreto, en público, a escondidas, de buena gana, a la fuerza, con cuidado, a ciegas, al azar} & ainsi, de cette façon, exprès, par hasard, à la main, à la machine, en secret, en public, en cachette, volontiers, de force, avec soin, à l'aveugle, au hasard \\
\hline
\textbf{CANTIDAD (Quantité / Degré)} & \esp{en exceso, en parte, en su mayor parte, por completo, a medias, más o menos, al menos, como máximo, como mínimo, poco a poco} & en excès, en partie, en grande partie, complètement, à moitié, plus ou moins, au moins, au maximum, au minimum, peu à peu \\
\hline
\textbf{AFIRMACIÓN / DUDA / NEGACIÓN} & \esp{desde luego, por supuesto, naturalmente, en efecto, tal vez, quizá(s), acaso, probablemente, en absoluto, de ningún modo, ni mucho menos} & bien sûr, naturellement, en effet, peut-être, probablement, absolument pas, nullement, loin de là \\
\hline
\end{tabularx}
\end{center}

\begin{exemplebox}[Les locutions en contexte]
\begin{itemize}
    \item \esp{De vez en cuando, visitamos el mercado de Mvog-Mbi.} « De temps en temps, nous visitons le marché de Mvog-Mbi. »
    \item \esp{El profesor llegó en seguida.} « Le professeur est arrivé aussitôt. »
    \item \esp{Poco a poco, los alumnos mejoran su expresión oral.} « Peu à peu, les élèves améliorent leur expression orale. »
    \item \esp{Rompió el jarrón a propósito.} « Il a cassé le vase exprès. »
    \item \esp{¿Te gusta el fútbol? --- ¡Desde luego!} « Tu aimes le football ? --- Bien sûr ! »
    \item \esp{No estoy de acuerdo en absoluto.} « Je ne suis absolument pas d'accord. »
\end{itemize}
\end{exemplebox}

\begin{propriete}[Place de la locution adverbiale dans la phrase]
En espagnol, la locution adverbiale se place le plus souvent \textbf{en tête de phrase} (suivie d'une virgule) ou \textbf{après le verbe / en fin de phrase}. Contrairement au français, elle ne s'intercale pas entre l'auxiliaire et le participe passé.\\
Français : « Il a \emph{souvent} mangé ici. » $\rightarrow$ Espagnol : \esp{Ha comido aquí a menudo.} ou \esp{A menudo ha comido aquí.} (jamais \esp{*ha a menudo comido}).
\end{propriete}

\coursec{Correspondances français-espagnol et pièges de traduction}

\begin{attention}[Faux amis et calques structurels]
\begin{itemize}
    \item \esp{De repente} = \textbf{soudain, tout à coup} (et non « de répète », qui n'existe pas en français).
    \item \esp{A propósito} = \textbf{exprès, intentionnellement}. Attention : « à propos / au fait » se dit \esp{por cierto} ou \esp{a propósito} selon le contexte, mais dans le sens de « volontairement », seul \esp{a propósito} convient. Compare : \esp{Lo hizo a propósito.} « Il l'a fait exprès. » / \esp{Por cierto, ¿vienes mañana?} « Au fait, tu viens demain ? »
    \item \esp{De vez en cuando} = \textbf{de temps en temps} (et non « de fois en fois »).
    \item \esp{Poco a poco} = \textbf{peu à peu} ; \esp{paso a paso} = \textbf{étape par étape, progressivement}.
    \item \esp{En seguida} (graphie recommandée ; \esp{enseguida} admise) = \textbf{aussitôt, immédiatement} (et non « en suite »).
    \item \esp{A menudo} = \textbf{souvent} (registre soutenu ; à l'oral on entend aussi \esp{mucho} : \esp{Voy mucho al cine.}).
    \item \esp{En absoluto} et \esp{de ningún modo} = \textbf{absolument pas, nullement} ; ils renforcent une négation : \esp{No lo aceptaré de ningún modo.} « Je ne l'accepterai nullement. »
    \item \esp{A ciegas} = \textbf{à l'aveugle} (et non « à cieques ») ; \esp{a escondidas} = \textbf{en cachette}.
    \item \esp{Tal vez} et \esp{quizá(s)} = \textbf{peut-être} ; ils sont suivis du \textbf{subjonctif} quand le doute est fort : \esp{Quizá llueva mañana.} « Il pleuvra peut-être demain. »
    \item \textbf{Ordre des mots} : ne pas placer la locution entre l'auxiliaire \esp{haber} et le participe passé (voir la règle ci-dessus).
\end{itemize}
\end{attention}

\begin{exoresolu}[Thème guidé : les locutions adverbiales]
Traduisez en espagnol : « Soudain, le marchand a fermé sa boutique et il est parti aussitôt. De temps en temps, il agit ainsi, exprès, pour nous surprendre. »
\tcblower
\textbf{Corrigé :} \esp{De repente, el vendedor cerró su tienda y se fue en seguida. De vez en cuando, actúa así, a propósito, para sorprendernos.}\\
\textbf{Commentaires :} \esp{de repente} en tête de phrase (soudain) ; \esp{en seguida} en fin de proposition (aussitôt) ; \esp{de vez en cuando} (de temps en temps) ; \esp{a propósito} (exprès) ; « pour nous surprendre » = but $\rightarrow$ \esp{para} + infinitif.
\end{exoresolu}

\coursec{L'essai argumentatif : structure et connecteurs}

\begin{methode}[Structure type de l'essai argumentatif en 4 parties]
\begin{enumerate}
    \item \textbf{Introducción} (1 paragraphe) : accroche (\esp{gancho} : citation, question, fait marquant) $\rightarrow$ présentation du thème $\rightarrow$ \cle{thèse} (position claire : \esp{En mi opinión, considero que...}) $\rightarrow$ annonce du plan (\esp{En primer lugar... ; además... ; por último...}).
    \item \textbf{Desarrollo} (2 à 3 paragraphes = 2 à 3 arguments) : une idée force par paragraphe. Structure : \cle{connecteur} + \cle{argument} + \cle{exemple ou explication} + \cle{mini-conclusion}. Utiliser \esp{por ejemplo, en concreto, esto demuestra que...}.
    \item \textbf{Contraargumentación} (facultative mais valorisée) : \esp{Es cierto que... ; sin embargo... ; no obstante...} (reconnaître le point de vue adverse pour mieux le réfuter).
    \item \textbf{Conclusión} (1 paragraphe) : \cle{réaffirmer la thèse} en d'autres termes $\rightarrow$ \cle{synthèse des arguments} $\rightarrow$ \cle{ouverture ou appel à l'action} (\esp{En definitiva, es fundamental que actuemos ya.}).
\end{enumerate}
\end{methode}

\begin{propriete}[Subjonctif et accords dans l'essai]
Dans un essai, on utilise fréquemment le \textbf{subjonctif} après les expressions impersonnelles et les verbes de volonté ou d'opinion niée : \esp{es necesario que protejamos el planeta} ; \esp{quiero que todos reciclen} ; \esp{no creo que sea suficiente}. Après \esp{aunque}, le subjonctif exprime l'hypothèse (\esp{aunque llueva}), l'indicatif le fait réel (\esp{aunque llueve}).
\end{propriete}

\begin{center}
\begin{tabularx}{\linewidth}{|X|X|X|}
\hline
\rowcolor{popgreenL} \textbf{Fonction} & \textbf{Conectores españoles} & \textbf{Français} \\
\hline
\textbf{Ordenar / Énumérer} & \esp{en primer lugar, en segundo lugar, finalmente, para empezar, por un lado... por otro lado} & premièrement, deuxièmement, enfin, pour commencer, d'un côté... de l'autre \\
\hline
\textbf{Ajouter / Renforcer} & \esp{además, también, asimismo, igualmente, incluso, más aún, no solo... sino también} & de plus, aussi, de même, même, qui plus est, non seulement... mais aussi \\
\hline
\textbf{Opposer / Concéder} & \esp{sin embargo, no obstante, aunque, a pesar de (que), por el contrario, en cambio, pero} & cependant, néanmoins, bien que, malgré, au contraire, en revanche, mais \\
\hline
\textbf{Cause / Conséquence} & \esp{porque, ya que, puesto que, como, por tanto, por consiguiente, así que, por eso, de modo que} & parce que, car, puisque, comme, donc, par conséquent, alors, c'est pourquoi, de sorte que \\
\hline
\textbf{Exemplifier / Préciser} & \esp{por ejemplo, en concreto, es decir, o sea, a saber, concretamente} & par exemple, en particulier, c'est-à-dire, c'est-à-dire, à savoir, concrètement \\
\hline
\textbf{Conclure / Résumer} & \esp{en conclusión, en definitiva, en resumen, para terminar, en síntesis} & en conclusion, en définitive, en résumé, pour finir, en synthèse \\
\hline
\end{tabularx}
\end{center}

\begin{attention}[Connecteurs : erreurs fréquentes des francophones]
\begin{itemize}
    \item « D'abord » (premièrement) = \esp{en primer lugar, para empezar} ; \esp{de pronto} signifie « soudain », pas « d'abord ».
    \item \esp{Sin embargo} se place en tête de proposition, suivi d'une virgule : \esp{Sin embargo, debemos actuar.}
    \item \esp{A pesar de} + nom ou infinitif ; \esp{a pesar de que} + verbe conjugué : \esp{a pesar de la lluvia} / \esp{a pesar de que llueve}.
    \item \esp{Como} en tête de phrase = « puisque » : \esp{Como contamina mucho, debemos limitar el coche.} Ne pas le confondre avec \esp{¿cómo?} (comment).
    \item Ne pas traduire « en fait » par \esp{*en facto} : dire \esp{de hecho, en realidad}.
\end{itemize}
\end{attention}

\coursec{Lexique du medioambiente et modèle d'essai}

\begin{exemplebox}[Plan détaillé : \esp{« La protección del medioambiente es responsabilidad de todos »}]
\textbf{Tesis :} \esp{La protección del medioambiente no es solo tarea de los gobiernos, sino un deber ciudadano diario.}
\begin{enumerate}
    \item \textbf{Introducción :} accroche (\esp{« No hay Planeta B »}) ; présentation du thème ; thèse ; annonce du plan.
    \item \textbf{Argumento 1 (individuel) :} \esp{En primer lugar}, les gestes quotidiens (recyclage, transports, consommation). Exemple : \esp{la recogida selectiva en Yaoundé}. \esp{Por tanto}, l'impact individuel est réel.
    \item \textbf{Argumento 2 (collectif) :} \esp{Asimismo}, l'éducation et la sensibilisation à l'école. Exemple : les clubs d'environnement au lycée. \esp{En efecto}, ils changent durablement les mentalités.
    \item \textbf{Contraargumento :} \esp{Es cierto que} les États et les entreprises polluent le plus. \esp{Sin embargo}, la pression citoyenne (vote, boycott) impose la régulation.
    \item \textbf{Conclusión :} \esp{En definitiva}, une responsabilité partagée. Appel : \esp{Actuemos hoy para no arrepentirnos mañana.}
\end{enumerate}
\end{exemplebox}

\begin{textebox}[Modèle d'essai (extrait rédigé)]
\esp{En primer lugar, cada ciudadano puede reducir su huella de carbono con gestos sencillos: reciclar, usar el transporte público y consumir productos locales. Por ejemplo, en Yaoundé, cada vez más familias hacen la recogida selectiva poco a poco. Por tanto, la acción individual tiene un impacto real.}\\
\esp{Es cierto que las grandes empresas contaminan en exceso. Sin embargo, la presión de los ciudadanos obliga a los gobiernos a actuar. En definitiva, es necesario que todos participemos: actuemos hoy para no arrepentirnos mañana.}
\end{textebox}

\textbf{Traduction :} « Premièrement, chaque citoyen peut réduire son empreinte carbone par des gestes simples : recycler, utiliser les transports en commun et consommer des produits locaux. Par exemple, à Yaoundé, de plus en plus de familles font peu à peu le tri sélectif. Donc, l'action individuelle a un impact réel. Il est vrai que les grandes entreprises polluent excessivement. Cependant, la pression des citoyens oblige les gouvernements à agir. En définitive, il est nécessaire que tous participions : agissons aujourd'hui pour ne pas regretter demain. »

\coursec{Entraînement guidé : thème, version et rédaction}

\begin{exercice}[Thème (Fr $\rightarrow$ Es)]
Traduisez : « De nos jours, la déforestation avance rapidement. Peu à peu, les citoyens comprennent qu'il faut agir immédiatement. Bien sûr, les gouvernements doivent aider, mais chacun peut recycler de temps en temps... non, tous les jours ! »
\end{exercice}

\begin{corrige}[Thème]
\esp{Hoy en día, la deforestación avanza rápidamente. Poco a poco, los ciudadanos comprenden que hay que actuar en seguida. Desde luego, los gobiernos deben ayudar, pero cada uno puede reciclar de vez en cuando... ¡no, todos los días!}\\
Points de méthode : \esp{hoy en día} (de nos jours) ; \esp{hay que} + infinitif (il faut) ; \esp{en seguida} (immédiatement) ; \esp{desde luego} (bien sûr) ; \esp{de vez en cuando} (de temps en temps).
\end{corrige}

\begin{exercice}[Version (Es $\rightarrow$ Fr)]
Traduisez : \esp{De repente, nos dimos cuenta de que la contaminación del río era grave. A menudo, los residuos llegan hasta el mercado. No podemos seguir así de ningún modo; es necesario que cambiemos nuestras costumbres.}
\end{exercice}

\begin{corrige}[Version]
« Soudain, nous nous sommes rendu compte que la pollution de la rivière était grave. Souvent, les déchets arrivent jusqu'au marché. Nous ne pouvons absolument pas continuer ainsi ; il est nécessaire que nous changions nos habitudes. »\\
Points de méthode : \esp{darse cuenta de que} = se rendre compte que ; \esp{de ningún modo} renforce la négation ; \esp{es necesario que} + subjonctif = il est nécessaire que + subjonctif en français.
\end{corrige}

\begin{exercice}[Rédaction (production écrite, 15 à 20 lignes)]
Rédigez un essai argumentatif intitulé : \esp{« ¿Es posible proteger el medioambiente en nuestra vida diaria? »}. Consignes : une introduction avec thèse, deux arguments illustrés d'exemples camerounais, un contre-argument réfuté, une conclusion avec appel à l'action, et au moins cinq connecteurs logiques différents.
\end{exercice}

\begin{corrige}[Rédaction : pistes de correction]
Attendus : thèse claire (\esp{En mi opinión, sí es posible...}) ; connecteurs variés (\esp{en primer lugar, además, sin embargo, por tanto, en definitiva}) ; exemples locaux (marchés de Douala, bendskins, tri des déchets, reboisement) ; subjonctif correct après \esp{es necesario que} ; accents et ponctuation espagnole (\esp{¿ ? ¡ !}) respectés.
\end{corrige}

\begin{aretenir}[Checklist avant le Bac -- Leçon E10]
$\square$ Je sais définir une locution adverbiale et donner 3 exemples par catégorie (temps, lieu, manière, quantité).\\
$\square$ Je traduis sans hésiter : \esp{de repente, a propósito, de vez en cuando, poco a poco, en seguida, a menudo, en absoluto}.\\
$\square$ J'évite les calques : \esp{a propósito} (exprès) $\neq$ \esp{por cierto} (au fait) ; \esp{de ningún modo} = nullement.\\
$\square$ Je place correctement la locution dans la phrase espagnole (début ou fin, jamais entre \esp{haber} et le participe passé).\\
$\square$ Je connais la structure en 4 temps de l'essai argumentatif (introduction, développement, contre-argument, conclusion).\\
$\square$ J'utilise au moins 5 connecteurs logiques variés (\esp{además, sin embargo, por tanto, por ejemplo, en definitiva}).\\
$\square$ Je maîtrise le lexique de base de l'environnement (\esp{calentamiento global, energías renovables, huella de carbono, biodiversidad}).\\
$\square$ Je sais écrire une thèse claire, nuancée et personnelle (\esp{En mi opinión / Considero que / Es evidente que}).\\
$\square$ Je développe un argument avec la structure : connecteur + idée + exemple concret + explication.\\
$\square$ Je gère le subjonctif après \esp{es necesario que, quiero que} et l'alternance indicatif/subjonctif après \esp{aunque}.\\
$\square$ Je relis pour vérifier les accents (\esp{á, é, í, ó, ú, ñ, ü}), la ponctuation (\esp{¿ ? ¡ !}) et l'orthographe des locutions.\\
$\square$ Je suis capable de faire un thème (Fr $\rightarrow$ Es) et une version (Es $\rightarrow$ Fr) sur ce thème en 20 minutes chacun.
\end{aretenir}

\findecours

\end{document}
